\section{The local obstruction and the real trace of the actual return}

\begin{statement}{P3K.1}{Proved}
Let

\[
 R=\mathbf C[u,w,X,v]/(uv-Xw),\qquad Z=\operatorname{Spec}R.
\]

Use the small resolution with charts

\[
 U_0=(u,w,z),\quad X=uz,\ v=wz;
 \qquad U_1=(X,v,t),\quad u=Xt,\ w=vt,\quad t=z^{-1}.
\]

Its exceptional curve is the zero section of
\(\mathcal O_{\mathbf P^1}(-1)^{\oplus2}\).  Write
\(L=\pi_{\mathbf P^1}^*\mathcal O(-1)\), with frames
\(e_1=z^{-1}e_0\).  Then the resolution map \(\rho\) satisfies

\[
 \rho_*L\cong I=(u,w),\qquad R^j\rho_*L=0\quad(j>0).
\]

Nevertheless \(I\) is a nonperfect rank-one module at the conifold
point.  The restriction of \(L\) to the link has primitive first Chern
class; it does not extend as a line bundle over the conifold smoothing.
\end{statement}

\begin{proof}
The two affine charts and their affine intersection compute
coherent cohomology.  A Laurent monomial on the intersection has the
form \(u^a w^b z^c e_0\), with \(a,b\geq0\).  It extends over \(U_0\)
when \(c\geq0\).  A monomial from \(U_1\) is

\[
 X^a v^b t^k e_1=u^a w^b z^{a+b-k-1}e_0,
\]

so it extends over that chart precisely for \(c\leq a+b-1\).
Those two integer intervals cover \(\mathbf Z\), including the case
\(a+b=0\).  Thus the first Cech cohomology vanishes.  Higher cohomology
vanishes because the cover has two affine opens and affine intersection.
The intersection of the two spaces of sections is spanned by the
monomials with \(0\leq c\leq a+b-1\).  Each has a factor \(u\) or \(w\)
without a factor \(z\); its remaining factor belongs to the subring
generated by \(u,w,uz,wz\), namely \(R\).  Hence the global sections are
generated by \(u e_0,w e_0\), which on the other chart are \(X e_1,v e_1\).
Sending these generators to \(u,w\) identifies the module with \(I\).
Since the target is affine, these cohomology calculations give all the
stated direct images.

The exact presentation and its matrix factorization are

\[
 R^2\xrightarrow{\Phi}R^2\xrightarrow{[u\ w]}I\longrightarrow0,
 \qquad
 \Phi=\begin{pmatrix}w&v\\-u&-X\end{pmatrix},\quad
 \Psi=\begin{pmatrix}-X&-v\\u&w\end{pmatrix},
\]

\[
 \Phi\Psi=\Psi\Phi=(uv-Xw)\,1_2.
\]

To check the presentation without appealing to a classification,
lift a relation \(au+bw=0\) to the polynomial ring.  Write
\(au+bw=c(uv-Xw)\).  Then
\(u(a-cv)+w(b+cX)=0\).  The coprimality of \(u,w\) gives
\((a,b)=d(w,-u)+c(v,-X)\), exactly the two displayed columns.
For exactness of the periodic tail, lift an element in
\(\ker\Phi\) to a polynomial column \(z\), so that
\(\Phi z=(uv-Xw)y\).  Multiplication by \(\Psi\) and cancellation in
the polynomial domain gives \(z=\Psi y\).  The other kernel is treated
by the same calculation with the matrices reversed.  Consequently

\[
 \cdots\xrightarrow{\Phi}R^2\xrightarrow{\Psi}R^2
 \xrightarrow{\Phi}R^2\longrightarrow I\longrightarrow0
\]

is an infinite free resolution.  At the maximal ideal
\(\mathfrak m=(u,w,X,v)\), every matrix entry lies in \(\mathfrak m\).
After tensoring with the residue field all differentials vanish and
all positive Tor groups have dimension two.  Thus \(I\) has no finite
locally free resolution.  In particular it has not become a zero
neighborhood class merely because \(H^*(\mathbf P^1,\mathcal O(-1))=0\).
It has generic rank one and \(\dim I/\mathfrak m I=2\).

The link of the resolution is the unit sphere bundle of the underlying
real rank-four bundle \(\mathcal O(-1)^{\oplus2}\) on \(S^2\).  Oriented
rank-four bundles on \(S^2\) are classified by
\(\pi_1(SO(4))=\mathbf Z/2\); the class here is
\(c_1\bmod2=(-2)\bmod2=0\).  The link is therefore \(S^2\times S^3\),
and projection to the exceptional \(S^2\) induces an isomorphism in
degree-two cohomology.  The restricted line has \(c_1=-h\), so its
rank-reduced K-class is the primitive generator of
\(\widetilde K^0(S^2\times S^3)=\mathbf Z\).  The smoothing is
\(S^3\times D^3\), whose second integral cohomology vanishes.  The
restriction of a line from that smoothing therefore has zero first
Chern class.  This proves the claimed obstruction to extending this
particular line.  Reversing the chosen Euler-line convention replaces
\(L\) by its dual and changes the generator's sign, without removing
the obstruction.  This establishes only the stated local obstruction;
it does not decide the actual return index.
\end{proof}

All polynomial identities are independently checked in
\path{code/part3_fourfold/conifold_module.py} and exported in the adjacent
JSON file.  The monomial bounds in the proof are symbolic; the finite
loop in the code is an additional sanity check, not its justification.

\begin{statement}{P3K.2}{Proved}
Put \(B=\mathbf C[u,v]\),
\(R'=B[x,y]/(xy-uv)\), and \(J=(x,u)\subset R'\).
The module \(J\) is flat over \(B\).  Its fibre at \((u,v)=(0,0)\) is
the normalization module

\[
 J\otimes_B\mathbf C\cong
 \mathbf C[x]\oplus\mathbf C[y]
 \quad\text{over}\quad\mathbf C[x,y]/(xy).
\]

For an integral projective nodal genus-one curve \(C\) with
normalization \(\nu:\mathbf P^1\to C\), the rank-one torsion-free
degree-zero sheaf \(\nu_*\mathcal O_{\mathbf P^1}(-1)\) has no
cohomology.  Thus the local magnetic module has an exact realization
inside a flat family of torsion-free nodal sheaves.  This supplies a
candidate coefficient object for index excision.  The local lemma
does not by itself compare the actual return symbol; P3K.6–7 below
provide that comparison by a real bordism.
\end{statement}

\begin{proof}
As a \(B\)-module, \(R'\) has the free monomial basis
\(1,x,x^2,\ldots,y,y^2,\ldots\), by reducing each mixed product using
\(xy=uv\).  The submodule \(J\) has the \(B\)-basis

\[
 \{x^n:n\geq1\}\ \cup\ \{u y^n:n\geq0\}.
\]

Spanning follows from \(xy^n=uvy^{n-1}\); independence follows from
the basis of \(R'\) and the fact that \(u\) is a nonzero divisor of
\(B\).  Hence \(J\) is free, in particular flat, over \(B\).
It is presented by the matrix
\(\left(\begin{smallmatrix}u&y\\-x&-v\end{smallmatrix}\right)\).
At \(u=v=0\) the relations become \(x e_2=0\), \(y e_1=0\), proving
the displayed normalization module.  The normalization map is finite,
so the cohomology of its direct image equals that of \(\mathcal O(-1)\),
which vanishes in both degrees.  Since \(\chi(\mathcal O_C)=0\), the
degree \(\chi(\nu_*\mathcal O(-1))-\chi(\mathcal O_C)\) is zero.
Its generic rank is one and it has no subsheaf supported at finitely
many points, so it is rank-one torsion-free.

There is also an exact connection to the local compactified-Picard
kernel.  On the fibre product of two nodal families, written
\(xy=UV\), the diagonal ideal is \((x-U,y-V)\).  Set
\(a=x-U\), \(b=y-V\).  The equation becomes
\(a(V+b)+Ub=0\).  With coordinates
\((x',u',y',v')=(a,b,V+b,-U)\), it is
\(x'y'=u'v'\), and the diagonal ideal is exactly \((x',u')\).
The kernel
\(I_\Delta\otimes\pi_1^*\mathcal O(O)\otimes\pi_2^*\mathcal O(O)\)
therefore has precisely this local matrix module, since the section
lines are trivial near the two nodes.
\end{proof}

For the global algebraic kernel construction, Burban–Kreussler,
*On a relative Fourier–Mukai transform on genus one fibrations*,
arXiv:math/0410349v2, equation (1), pp. 5–7, constructs this sheaf flat
over both factors for a flat projective family of integral Gorenstein
genus-one curves with a section in its smooth locus.  Theorem 2.12
proves the resulting derived equivalence.  The theorem concerns that
specified algebraic family and kernel.  The required identification
with the present circle quotient's symbol still needs its own map.

\begin{statement}{P3K.3}{Proved}
Let \(B=S^3\times S^1\), let \(K\subset S^3\) be the cusp circle,
and set \(T=K\times S^1\).  Give the trivial real normal two-plane
bundle of \(T\) its oriented complex-line structure.  Its K-theory
Thom map and extension by zero give

\[
 i_!:K^0(T)\cong K^0_c(N_T)\longrightarrow K^0(B).
\]

In the rank/degree basis of \(K^0(T)=\mathbf Z^2\), this map kills
the rank generator and maps the degree-one reduced line class to the
primitive top generator of \(\widetilde K^0(B)=\mathbf Z\).  Every
integer is consequently realized by a K-class supported on the cusp
torus and vanishing on its complement.  The spin index on \(B\)
reads exactly this integer.
\end{statement}

\begin{proof}
The oriented normal bundle of a knot in an oriented
three-manifold is trivial: its Euler class lies in \(H^2(S^1)=0\).
The same trivialization pulls back to \(T\).  The K-theory Kunneth
formula, or the cellular filtration and Bott periodicity, gives

\[
 K^0(T)=\mathbf Z\{1\}\oplus\mathbf Z\{\beta\},\quad
 \operatorname{ch}(\beta)=[T]^*,
\]

where \(\beta=[H]-1\) for a line \(H\) of degree one, and

\[
 K^0(B)=\mathbf Z\{1\}\oplus\mathbf Z\{\omega\},\quad
 \operatorname{ch}(\omega)=[B]^*.
\]

The normal Todd class is one, so the Chern character of the Thom
pushforward is the ordinary cohomological Gysin map.  For the rank
generator, \(\operatorname{ch}(i_!1)=\operatorname{PD}(T)=0\), because
\(H^2(B;\mathbf Z)=0\).  Both K-groups are torsion-free, so \(i_!1=0\).
For \(\beta\), the Gysin image integrates to one in degree four;
hence \(i_!\beta=\omega\).  Extension by zero makes its restriction
to \(B\setminus T\) zero.  Finally \(B\) is parallelizable, so its
spin index is \(\int_B\operatorname{ch}_2\), which sends \(\omega\)
to one.
\end{proof}

This proves an explicit limitation of the inference from acyclicity
on regular elliptic fibres: even a known pushforward on the complement
does not choose the extension's top K-class.  It is an actual
nonzero image under a specified Thom/Gysin map, rather than an
unsupported assertion that the local and global calculations are
unrelated.

\subsection{Initial index-excision target, resolved by the real trace below}

The following records the precise target of the initial algebraic
index-excision attack.  P3K.6–7 subsequently construct the required
actual comparison by a real trace and a sphere-bundle bordism.

For the actual conifold replacement mapping torus \(Q^6\), write
\(L_\delta\) for its circle Euler line.  The required integer is

\[
 A=\operatorname{ind}(D_{Q,\mathrm{spin}}\otimes L_\delta).
\]

On the ordinary elliptic pieces, the coefficient line has Picard point
\(P\), and the corresponding degree-zero Dolbeault complexes are
acyclic.  On the compact fixed sheet, P3F.12 supplies
\(R f_*(M+nF)=0\).  P3F.14 supplies the exact local transfer to that
complex.  P3K.1 shows why the conifold neighborhood cannot be discarded
merely from the fibre \(\mathcal O(-1)\)'s cohomology.  P3K.2 gives an
explicit flat torsion-free alternative retaining its link module.

The mathematical object sought at this stage was a relative symbol comparison
between the actual spin symbol twisted by \(L_\delta\) and this
nodal Dolbeault complex, with its two outer trivializations.  It must
identify the primitive meridian class and the degree \(-12\) angular
gauge in P3F.13–15, and must account for the added source-resolution
curves on the outer boundary.  None of P3K.1–3 constructs that relative
symbol map.  The universal Fourier–Mukai theorem supplies the coherent
kernel once a family and pullback map are given; its hypotheses do
not themselves provide the map from this orbit family.

The unsuccessful algebraic attack used the candidate outer morphism
given by the minimal source
resolution of \(UV=Q^{12}\), followed by
\((U,V,Q)\mapsto(U,V,q=Q^{12})\), crossed with the nondegenerate second
direction \(r\ne0\), and then the labelled cyclic descent of P3F.15.
At that stage the descended outer quotient had not been identified
with a proper equivariant modification of the original smooth
family.  P3K.4 supplies the precise multiplicity obstruction to
that proposed map.  The real trace P3K.6–8 supplies the actual
return comparison without using that failed algebraic extension.

\begin{statement}{P3K.4}{Proved}
The free twelve-label construction of P3F.15, when applied to a
reduced \(Q\)-central divisor with no stabilizers at its generic
points, has multiplicity twelve over \(q=Q^{12}\) on every image
component.  A normal space with all those multiplicities equal to
twelve admits no proper bimeromorphic map over the \(q\)-disc to the
actual reduced cusp model.  Thus a common proper toric modification
of these two outer extensions cannot be inferred from their
punctured-family comparison.
\end{statement}

\begin{proof}
A finite quotient by an action with trivial stabilizer is
etale near the generic point under consideration.  For a reduced
component \(D\) of \(Q=0\), its divisorial valuation gives
\(\operatorname{ord}_D(q)=\operatorname{ord}_D(Q^{12})=12\).
The ramification index of the quotient at \(D\) is one, so the order
on its image is also twelve.  A proper bimeromorphic morphism between
normal spaces is an isomorphism at the generic point of the strict
transform of any target prime divisor.  It therefore preserves that
divisor's valuation of \(q\).  A target component of order one would
give a source component of order one, contradicting the displayed
orders.  This proves the assertion about proper modifications.

The exact origin itself detects why the two extensions differ.
For the original default twists, P2.7 gives

\[
 F(s,v)=\left(\frac{\mu(s+12v)}2,
 \frac{C(s+12v)}{12}-\frac{(s+12v)^2}{24}+6v\right),
 \qquad \gamma(s,v)=(6\mu,C-s-12v).
\]

Here \(\mu,C\) are unchanged by the integer cusp deck.  Direct
subtraction, with no limiting argument, yields

\[
 F(s+1,v)-F(s,v)=\frac{\gamma(s,v)}{12}-\frac{\delta}{24},
 \qquad
 F(s+12,v)-F(s,v)=\gamma(s,v)-6\delta.
\]

Thus the partial-origin one-step deck is the fractional gamma
translation that shifts the twelve-cycle.  It is not the reduced
cusp extension furnished by the extending cusp origin.  The exact
identities are exported by \path{code/part3_fourfold/outer_descent.py}.
They do not contradict the angular clutch in P3F.13, nor the partial
filling.  They prevent identifying the freely descended outer model
with the original reduced cusp by the proposed proper-map shortcut.
\end{proof}

\begin{statement}{P3K.5}{Proved}
The integral connecting class from the marked boundary \(T^3\) has square zero and pairs trivially with every extension of the page line.  In particular an odd cap correction does not change the normal Euler square.
\end{statement}

\begin{proof}
Let \(V=V_0\) be the actual regular-binding page, with its marked
boundary \(T^3\), and let \(m_V=c_1(M)\).  Its boundary restriction is
zero.  For the inverse zero-framed Borromean surgery trace \(H\),
\(H^2(H;\mathbf Z)\to H^2(T^3;\mathbf Z)\) is an isomorphism: the
three-handle intersection matrix is zero, and the relevant
Poincare–Lefschetz exact sequence has last term \(H_1(H)=0\).
Consequently any class on \(V\cup H\) that restricts to \(2m_V\) on
\(V\) and to zero on the common boundary restricts to zero on \(H\).
Choose an integral extension \(m\) of \(m_V\), using a boundary
trivialization and the trivial line on \(H\).  The difference between
that class and \(2m\) lies in the image of
\(\delta:H^1(T^3;\mathbf Z)\to H^2(V\cup H;\mathbf Z)\).
For \(k=\delta g\), the relative cup-product formula gives

\[
 k^2=\delta(g\smile k|_{T^3})=0,\qquad
 km=\delta(g\smile m|_{T^3})=0.
\]

Thus its square is always \(4m_{V,\mathrm{rel}}^2\), even when the
gluing correction is odd.  The odd correction is relevant here:
\(V\subset S^5\) is spin and has signature \(-8\), so its boundary
spin structure has Rokhlin invariant eight.  The zero-framed
Borromean trace is spin of signature zero and induces Rokhlin
invariant zero.  Those spin structures cannot be identified.  This
does not change the square calculation, which retains the integral
connecting class rather than setting it to zero.

The following real trace construction uses a different cap and does
not require an embedding of this Borromean trace in the six-sphere.
\end{proof}

\begin{statement}{P3K.6}{Proved}
Let one positive gamma turn denote the actual parameter cover whose
angular clutch is P3F.13.  The orbit family of this cover is the product
\(Y\times S^1_\theta\), with \(Y=S^5\), and its fixed trace is

\[
 T_F=\{(R_\theta(p),\theta):p\in F\},
\]

where \(R_\theta\) is the gamma rotation on the cusp orbit product of
P3F.9.  Use its product trivialization preserving the Seifert gamma
coordinate.  In these coordinates \(R_\theta\) is pure gamma
rotation, including the explicitly trivialized source-translation
suspension.  There is a compact properly embedded oriented four-manifold
\(X\subset Z=Y\times D^2\), with boundary \(T_F\), assembled from the
actual page \(V\) and a product compression trace.  Its complement
has a unique integral class \(x\) with meridian value one.  On the
outer boundary this class is the actual delta-circle Euler class.
On the graph copy of \(V\), its restriction to the framed normal
section is exactly \(c_1(M)\).
\end{statement}

\begin{proof}
\par\medskip\noindent\textbf{The actual orbit trivialization.}\enspace The ordinary and finite
orbit charts are independent of the parameter angle: the variation
of their period and affine-twist columns is in the real delta
direction.  The exact subtraction in P3F.13 leaves, on the cusp
overlap, the gamma translation \(R_\theta\); the residual delta term
is a periodic unitary gauge.  Apply the identity to the exterior
charts and \(R_\theta\) to the cusp chart.  Their overlap maps now
agree pointwise.  Since \(R_1=R_0\), these identifications give a
periodic trivialization of the orbit family.  The fixed sphere in
the cusp chart becomes precisely \(R_\theta F\).  This verifies the
product and the displayed trace for the actual angular family.

\par\medskip\noindent\textbf{The compression.}\enspace Choose the binding at a regular value
\(q=\epsilon>0\) in a small cusp disc.  First write the local
regular-binding Milnor function as
\(u(q)(q-\epsilon)e^{2\pi i12\gamma_{\rm old}}\), with \(u\ne0\),
choose its unit argument, and set
\(\gamma=\gamma_{\rm old}+\arg u(q)/(2\pi\,12)\).
This is an explicit base-dependent circle coordinate change,
commuting with rotation and preserving the integral gamma marking.
Make this unit correction before choosing the short arc and the
handle below.  At its node the source has
coordinates \(xy=q\).  The two normalization discs of the fixed
sphere project to \(x=0\) and \(y=0\) and are separated in the gamma
coordinate.  Keep their values and the rest of the fixed sphere
inside an arc \(J\) of length less than \(1/24\), as allowed by the
relative product construction P3F.9.  Write their actual real gamma
functions as \(g_x(z)\) and \(g_y(z)\), without changing the orbit
map or its page.

Use the following exact real handle in the node chart:

\[
 (z,\phi)\longmapsto
 \bigl(x=z\cos\phi,\quad y=\bar z\sin\phi,\quad
       \gamma=g(z,\phi)\bigr),
 \qquad 0\leq\phi\leq\pi/2,
\]

Here \(g\) interpolates from \(g_x(z)\) to \(g_y(\bar z)\), with
values in \(J\); for example use their convex linear interpolation
in \(2\phi/\pi\).  The two center values are distinct, so
\(g(0,\phi)\) is strictly monotone.  For \(z\ne0\), the pair
\((x,y)\) determines \((z,\phi)\) by its moduli and phases.  At
\(z=0\), gamma determines \(\phi\), and at least one of
\(\sin\phi,\cos\phi\) is nonzero.  These facts also give full
rank to the derivative, so the map is an embedding at the core.
Its exact base value is

\[
 q=xy=|z|^2\cos\phi\sin\phi\geq0.
\]

Restrict to \(|x|,|y|\leq a\) and \(q\leq\epsilon\).
The faces \(\phi=0,\pi/2\) are the two normalization discs at
separated gamma values.  The face \(q=\epsilon\) is the smoothing
annulus, and the remaining faces \(|x|=a\), \(|y|=a\) match the
transported exterior of the source family along the real path.
Their gamma functions extend over that exterior with values in
\(J\), since \(J\) is an interval.  Rounding only these product
corners gives the oriented compression body.  Its two disc-face
orientations are the complex normalization orientations: the
conjugation on the second face is canceled by its opposite boundary
orientation.  Equivalently, reversing it attaches the positive
smoothing one-handle to \(F\).  Its regular boundary is the source
torus embedded as a graph \(\gamma=g_\epsilon(p)\), with the
globally defined real-valued function \(g_\epsilon\) taking values
in \(J\).  Choose a collar-preserving height function to obtain
an embedded three-dimensional compression cobordism

\[
 W\hookrightarrow Y\times[r_0,1],\qquad
 \partial W=F\sqcup(-E_\epsilon).
\]

Its projection lies over \(0\leq q\leq\epsilon\) and inside \(J\),
by the exact displayed formula.  The trace is used in the product
with its height interval.  Topologically \(W\) is a solid torus
with an open three-ball removed.  Its meridian on the incoming
torus is the source vanishing cycle \(a\).  The explicit handle
also fixes the framing and avoids an uncontrolled normal
perturbation of the base coordinate.

\par\medskip\noindent\textbf{Choose the page and embed the graph.}\enspace The weighted circle
has degree twelve on the regular-binding Milnor function.  In the
already chosen unit-corrected coordinate, the page equation is

\[
 \arg(q-\epsilon)+2\pi\,12\gamma=\varphi.
\]

Along the real path \(0\leq q<\epsilon\), its twelve gamma values
are constant.  Choose \(\varphi\) so that all these values miss
\(J\), and so that the short arc \(12J\) lies strictly on one side
of the page near the binding.  This is possible because of the
chosen length of \(J\).  The interior of the actual page \(V\)
then misses the entire projected compression, not just \(F\).
At the binding the initial torus meets \(\partial V\); pushing the
page a small distance to the opposite page-normal side removes this
intersection as well.  Denote this pushoff by \(V^+\).  It is
isotopic to \(V\) in \(Y\setminus F\) and is disjoint from the
projected compression.

On \(\partial V=E_\epsilon\times S^1_\gamma\), define
\(f(p,\gamma)=r_0e^{2\pi i(\gamma-g_\epsilon(p))}\).
This is the exact inverse angular coordinate of the rotating
incoming graph.  Its integral marking is the gamma marking,
because \(g_\epsilon\) is a globally defined real-valued function.
Extend this to a smooth map
\(f:V\to D^2_{r_0}\): on a boundary collar multiply the displayed
value by a cutoff that decreases to zero and is flat at zero,
then set \(f=0\) on the remaining interior.  Its graph in
\(Y\times D^2\) is embedded because its projection to \(Y\) is the
original embedded page.  In the outer annulus attach

\[
 W\times S^1\longrightarrow Y\times D^2,
 \qquad ((y,r),\theta)\longmapsto
 (R_\theta y,r e^{2\pi i\theta}).
\]

The polar coordinate records both \(r\) and \(\theta\), so this
map is an embedding.  Its incoming boundary is exactly the graph
of \(f|_{\partial V}\), and its outgoing boundary is \(T_F\).
Choose the radial derivative of the graph collar nonzero.  Extend
that collar outward and adjust only the outer-annulus collar of
the compression to meet it.  In a normal chart of the common
binding these are graphs over the same radial collar, so smooth
interpolation gives an embedded union.  All this adjustment is
arbitrarily small and leaves the graph inside the closed inner
disc unchanged.  At \(\theta=0\) it is chosen on the compression
side of the page; the fixed opposite pushoff \(V^+\) remains
disjoint.  This is the asserted \(X\).  Novikov additivity gives
\(\sigma(X)=\sigma(V)\): the added product \(W\times S^1\) has zero
signature, as is also seen by reversing its circle factor.

\par\medskip\noindent\textbf{The euler line, including its clutch.}\enspace The integral Thom
sequence for the proper codimension-three embedding \(X\subset Z\)
contains

\[
 H^2(Z;\mathbf Z)=0\longrightarrow H^2(Z\setminus X;\mathbf Z)
 \longrightarrow H^0(X;\mathbf Z)=\mathbf Z
 \longrightarrow H^3(Z;\mathbf Z)=0.
\]

It gives exactly one class \(x\) with the chosen meridian value one.
The identical calculation for
\(T_F\subset Y\times S^1\) identifies
\(H^2((Y\times S^1)\setminus T_F;\mathbf Z)=\mathbf Z\)
by the same meridian.  The actual delta bundle has that value,
so its first Chern class is the restriction of \(x\).  Classification
of principal circle bundles by integral first Chern class supplies
an isomorphism of these specified bundles.  This is an integral
identification on the whole complement; it retains the angular
clutch of degree minus twelve, rather than deleting its gauge.

The normal bundle of the graph copy of \(V\) is the sum of its
oriented page-normal real line and the real tangent plane of the
parameter disc.  The graph derivative gives an upper-triangular
identification with this sum, hence a framing.  Take the section
pointing toward \(V^+\).  Its pushoff first moves its disc coordinate
to zero inside the radius \(r_0\); its \(Y\)-coordinate is off the
graph throughout.  It then moves radially from zero to the outer
parameter point \(\theta=0\).  In the outer annulus the trace at
\(\theta=0\) is exactly \(W\), whose projection misses \(V^+\).
This is a homotopy in \(Z\setminus X\) to the pushoff of the actual
page in the original parameter-zero exterior.  The restriction of
\(x\) along that normal section is therefore

\[
 s^*x=c_1(L_\delta|_V)=c_1(M).
\]

The last equality is the exact circle-line comparison P3F.10.
Boundary collar choices change only its boundary gauge, and their
relative square is zero by the relative cup-product formula.
This completes the construction and the meridian comparison.
\end{proof}

\begin{statement}{P3K.7}{Proved}
The real trace gives a spin bordism carrying the Euler line from the actual covered conifold return to a normal sphere bundle.  Its index is zero.  The original return also has index zero, and its actual Hopf-lift gluing is the standard sphere.
\end{statement}

\begin{proof}
Cap \(X\) along \(F\times S^1\) with \(C=D^3\times S^1\), using
the fixed normal framing chosen for the actual conifold replacement.
Denote the resulting closed four-manifold by \(\bar X\).  Its cap
relative to \(V\) is

\[
 H_a=(W\cup_FD^3)\times S^1
      =D^2_a\times S^1_b\times S^1_\gamma.
\]

The subscripts retain the marked boundary cycles.  This cap kills
\(a\).  The usual elliptic cap kills \(\gamma\); their boundary
markings are not silently equated.  Both have signature zero.

Delete the normal disc bundle of \(X\) from \(Z\), and glue
\((C\times S^2)\times I\) along the resulting boundary collar.
The resulting two closed boundary components are the physical
positive return family \(P_f\) and the normal sphere bundle
\(B=S(\bar\nu)\to\bar X\).  On the outer side this is precisely
the replacement of the fixed \(F\times D^3\) by \(D^3\times S^2\),
performed in every parameter fibre.  On the inner side it completes
the normal sphere bundle over the cap.  We now specify its two
boundary orientations and the return convention explicitly.

Let the circle-first quotient convention be
\(\Omega_{X_6}=d\delta\wedge\Omega_Y\).  In the fixed normal
complex two-plane the weights are \((-1,+1)\).  Conjugating the
negative-weight complex coordinate reverses its real orientation;
contracting the positive Hopf volume by the circle vector reverses
the outward cone convention once more.  The two signs give
\(\Omega_Y=\Omega_F\wedge\Omega_{\nu,\mathrm{Hopf}}\), where
the meridian sphere has the positive Hopf orientation.  The graph
part of \(X\) has the complex page orientation, and the oriented
compression has \(\partial X=F\times S^1_\theta\), in that order.
Give \(Z=Y\times D^2\) its product orientation.  Since \(\dim Y=5\),
the outward boundary orientation is

\[
 \Omega_{\partial Z}=-\Omega_Y\wedge d\theta.
\]

The inner boundary of the deleted normal disc bundle is oriented
\(-\Omega_{\bar X}\wedge\Omega_{S^2}\).  The product cap preserves
these signs.  Thus the actual oriented bordism has boundary
\(-P_f\sqcup(-B)\), with \(P_f\) oriented by the fibre followed
by positive parameter, and \(B\) by the complex page followed
by the positive meridian sphere.

Jin's mapping torus \(M_f\) identifies \((z,0)\sim(f(z),1)\).
The physical positive family of the return \(f\), which identifies
\((z,1)\sim(f(z),0)\), is therefore \(P_f=M_{f^{-1}}\).
Characteristic-number additivity under inversion, together with
the two boundary signs just computed, gives

\[
 (k,\ell,A)(P_f)=-(k,\ell,A)(B),\qquad
 (k,\ell,A)_{\mathrm{Jin}}(f)=(k,\ell,A)(B).
\]

This records both reversals.  The following integrations on \(B\)
therefore compute Jin's invariants of the specified physical
positive return \(f\), with the stated positive meridian class.

The Euler line extends over the bordism.  On
\(F\times S^1\times S^2\) its Chern class is the meridian generator,
by the positive-weight SU(2) trivialization of P3F.7.  It therefore
extends over \(D^3\times S^1\times S^2\).  An ambiguity in the
boundary bundle isomorphism is a map to \(S^1\); its homotopy class
lies in \(H^1(S^2\times S^1\times S^2)=\mathbf Z\), and the
restriction from \(H^1(D^3\times S^1\times S^2)\) is an isomorphism.
It extends as well.  The bordism is spin: \(Y\times D^2\) is spin,
and the two possible spin structures on the gluing boundary differ
only on the same \(S^1\) factor, where either extends over the cap.
This identifies the actual spin return symbol and Euler line with
the inner boundary symbol by spin bordism, with no analytic corner
or contraction of exceptional curves.

Write \(e_v\) for the vertical Euler class of \(S(\bar\nu)\), and
define \(b\in H^2(\bar X;\mathbf Z)\) by

\[
 e_v=2x-\pi^*b.
\]

Existence and uniqueness follow from the sphere-bundle Gysin
sequence, since both classes have fibre integral two.  Since the
normal bundle of the graph page is framed and its chosen section
has \(s^*e_v=0\), P3K.6 gives \(b|_V=2c_1(M)\).  On the product
cap \(C\), \(b=0\).  Its restriction on the intermediate product
compression is zero as well: \(H^2(W\times S^1)\) injects into the
cohomology of its two boundary components.  This is seen from
\(H_1(W)=\mathbf Z\), \(H_2(W)=\mathbf Z\), and the boundary
torus and sphere, whose products with the circle generate both
classes of \(H_2(W\times S^1)\).  On the two boundaries the relevant
restriction is zero.  Equivalently,
\(H^2(H_a)\to H^2(T^3)\) is injective, so the vanishing boundary
restriction of \(b\) gives its vanishing on this whole cap.

Extend \(M\) trivially across \(H_a\), with class \(m\).  The
Mayer–Vietoris difference \(b-2m\) belongs to
\(\delta H^1(T^3;\mathbf Z)\).  As in P3K.5, its square and its
product with \(m\) vanish.  Therefore

\[
 \langle b^2,[\bar X]\rangle=4m_{V,\mathrm{rel}}^2=-8,
 \qquad
 \sigma(\bar X)=\sigma(V)=-8,
 \qquad
 \langle p_1(T\bar X),[\bar X]\rangle=-24.
\]

These use the actual page invariants P3F.10.  No identification of
the two different marked caps is needed.

Finally \(T_{\mathrm{vert}}B\oplus\mathbf R\cong\pi^*\bar\nu\),
so \(e_v^2=\pi^*p_1(\bar\nu)\).  With \(\pi_*x=1\), this gives

\[
 \pi_*x^2=b,\qquad
 \pi_*x^3=\frac{p_1(\bar\nu)+3b^2}{4},\qquad
 p_1(TB)=\pi^*(p_1(T\bar X)+p_1(\bar\nu)).
\]

The actual spin return index is consequently

\[
 A(f)=\frac{4\int_Bx^3-\int_Bp_1(TB)x}{24}
     =\frac{3\int_{\bar X}b^2-\int_{\bar X}p_1(T\bar X)}{24}
     =\frac{3(-8)-(-24)}{24}=0.
\]

The normal Pontryagin class cancels.  Thus no unproved normal
framing or spin claim about the elliptic cap enters the result.
For the original parameter return, P3F.11 states that this covered
return has index six times the original integer.  Since that
integer is integral, the original index is also zero.  The actual
Hopf-lift gluing P3F.8 and the extending \(a\)-powers P3F.4 then
identify the clutched sphere as standard.
\end{proof}

\begin{statement}{P3K.8}{Proved}
The same actual return index vanishes by a second calculation: vertical sphere quantization, an explicit marked cap switch, and acyclicity of the elliptic Picard line.  This calculation does not evaluate the page signature, its line square, or its first Pontryagin number.
\end{statement}

\begin{proof}
The actual bordism of P3K.6–7 also gives an analytic computation of
\(A\) that does not evaluate the signature of the page, the square
of its line, or a Pontryagin number.  Keep the inner sphere bundle
\(\pi:B=S(\bar\nu)\to\bar X\) and the integral class
\(e_v=2x-\pi^*b\).

\par\medskip\noindent\textbf{Vertical quantization.}\enspace The spin Dirac symbol on \(B\),
twisted by the Euler line of class \(x\), is a spin-c symbol of
determinant \(2x\).  The oriented vertical two-plane has its
canonical complex spin-c symbol, of determinant \(e_v\).  Dividing
these two spin-c structures gives the pullback of a base spin-c
structure of determinant \(b\).  This follows directly from the
stable tangent splitting into vertical and horizontal bundles and
the displayed determinant identity.  Fibrewise the vertical symbol
is the Dolbeault symbol of \(\mathcal O_{\mathbf P^1}\), not that of
\(\mathcal O(-1)\).  Its kernel consists of constant functions and
its cokernel is zero.  The constants are fixed by every oriented
sphere transition, so its family index is the trivial line, with
its actual transition functions.  This symbol composition can be
verified directly at the operator level.  Give the fibres their
round metrics and use an \(SO(3)\) connection on \(\bar\nu\).
The vertical Dolbeault Dirac operator \(D_v\) is invariant under
this structure group, so the horizontal connection commutes with
\(D_v\).  Its squared restriction orthogonal to the constants
has a positive uniform lower bound.  If \(\epsilon_v\) is the
vertical grading and \(D_h\) the horizontal base spin-c Dirac
operator on the bundle of vertical spinors, the operator

\[
 D_t=tD_v+\epsilon_vD_h,\qquad t>0,
 \qquad D_t^2=t^2D_v^2+D_h^2
\]

is invertible on that orthogonal complement.  The square identity
uses \(D_v\epsilon_v=-\epsilon_vD_v\) and the commuting invariant
connection.  On the constant vertical kernel it is precisely the
base Dirac operator.  The displayed operator has the same elliptic
symbol as the twisted total-space spin Dirac operator after a
positive rescaling of the vertical metric; zeroth-order terms do
not change its index.  Consequently

\[
 A=\operatorname{ind}(D_{\bar X,b}).
\]

This is a factorization of the specified actual symbol.  It keeps
the meridian Dirac index one that is essential at the conifold
replacement.

\par\medskip\noindent\textbf{Switch the marked cap.}\enspace The compression cap is
\(H_a=D^2_a\times S^1_b\times S^1_\gamma\), whereas the regular
elliptic cap is \(H_\gamma=S^1_a\times S^1_b\times D^2_\gamma\).
Use their common marked \(T^3\) boundary.  Their closed difference
is

\[
 H_a\cup_{T^3}(-H_\gamma)
 =S^1_b\times
   \bigl((D^2_a\times S^1_\gamma)
       \cup_{T^2}(S^1_a\times D^2_\gamma)\bigr)
 \cong S^1\times S^3.
\]

The indicated meridian-longitude exchange is the usual two-solid-
torus decomposition of \(S^3\); its meridian intersection is one.
Both cap determinant classes vanish and their boundary spin-c
structures agree: on \(T^3\), a spin-c structure with determinant
class zero is unique because \(H^2(T^3;\mathbf Z)\) has no
two-torsion.  Glue with this boundary identification.  The closed
difference has \(H^2=0\), so it has its spin structure as a spin-c
structure and extends over \(S^1\times D^4\).  Its Dirac index is
zero by spin-c bordism.

For completeness, additivity here uses the cut-and-paste bordism
obtained by taking the cylinders on the common page and on the
two caps, and joining their common boundary collars by a pair of
pants.  The three boundary components are the old capped page,
the oppositely oriented new capped page, and the oppositely
oriented closed cap difference.  The same spin-c structure on the
common collar extends over this product construction.  Thus the
index of \((\bar X,b)\) is equal to that of the compact sectioned
elliptic page \(S\) with the resulting determinant \(b_\gamma\).

\par\medskip\noindent\textbf{Identify the base spin-c structure.}\enspace The two integral
vanishing cycles needed here can be exhibited directly.  In the
source marking of P3F.12, take

\[
 A_E=\begin{pmatrix}0&1\\-1&-1\end{pmatrix},\qquad
 B_E=\begin{pmatrix}0&-1\\1&0\end{pmatrix}.
\]

The nodal monodromy \(A_EB_E\) has image of its nilpotent part
generated by \(a=(0,1)^t\).  The twelve unramified cusp lifts of
the regular-binding page include the transport \(A_Ea=(1,-1)^t\).
The two columns \(a,A_Ea\) have determinant minus one.  Van Kampen
therefore kills the whole fundamental group of a torus fibre.
The section then identifies the fundamental group of the compact
elliptic page with that of its genus-three base, and identifies
the fundamental group of the open page with that of the base
with one boundary component.  In particular \(H_1(S;\mathbf Z)
\cong\mathbf Z^6\), so \(H^2(S;\mathbf Z)\) has no torsion.

On the marked boundary \(T^3\), the two source loops die in
\(H_1(V)\), and the base boundary loop is a product of
commutators and also dies.  Hence the restriction
\(H^1(V)\to H^1(T^3)\) is zero.  The image of
\(H^1(H_\gamma)\) is the span of the two source characters.
The Mayer–Vietoris kernel of restriction to \(V\) and
\(H_\gamma\) is consequently the infinite cyclic group generated
by the fibre class \(F\).  It is primitive, since the section
intersects a fibre once.

The elliptic canonical class \(K_S\) is pulled back from the base,
so it restricts to zero on both the page with boundary and the
regular cap.  The actual determinant and the canonical spin-c
determinant twisted by the Picard line thus satisfy

\[
 b_\gamma-\bigl(-K_S+2c_1(M)\bigr)=kF
\]

for an integer \(k\).  Both are characteristic classes.  Since
\(F\) is primitive, their common mod-two class makes \(k=2n\)
even.  Therefore

\[
 b_\gamma=-K_S+2c_1(M+nF).
\]

There is no two-torsion ambiguity in the spin-c structure on
\(S\).  Its Dirac operator is exactly the Dolbeault spin-c
operator twisted by \(M+nF\).

\par\medskip\noindent\textbf{Acyclicity from explicit sections.}\enspace The source
primitive-generator identification fixes, up to sign,

\[
 P_*=(0,\sqrt2\,t^2(t-1))
 \quad\hbox{on}\quad
 y^2=x^3-3t^3(t-1)x+2t^4(t-1)^2.
\]

We use that identification only to specify the line in the
construction.  The following direct calculation of its fibrewise
nontriviality does not use the height, the signature, or the square
of the page line.  In the ordinary affine charts the displayed
point is finite and hence different from the zero section.  At the
\(IV^*\) fibre take \(t=z^3\), \(x=z^4X\), \(y=z^6Y\).
Then

\[
 P'_*=(0,\sqrt2(z^3-1))\longrightarrow(0,-\sqrt2)
 \quad\hbox{on }Y^2=X^3+2.
\]

At the \(III\) fibre take \(t=1+z^4\), \(x=z^2X\),
\(y=z^3Y\).  This gives

\[
 P'_*=(0,\sqrt2\,z(1+z^4)^2)\longrightarrow(0,0)
 \quad\hbox{on }Y^2=X^3-3X.
\]

Both limiting points are finite smooth points of the good-reduction
elliptic curves and are different from the point at infinity.
At the nodal fibre put \(q=1/t\), \(X=q^2x\), \(Y=q^3y\).
The section is \((0,\sqrt2(1-q))\), with limiting point
\((0,\sqrt2)\) on \(Y^2=(X-1)^2(X+2)\).  It is smooth and
different from both the node \((1,0)\) and the zero section.
The same statements hold for \(-P_*\); fibre inversion gives
the exact comparison between these two choices.  Pullback under
the actual regular-binding cover preserves these explicit
sections and their nonidentity specializations.  Every smooth
elliptic fibre therefore carries a nontrivial degree-zero line,
whose zeroth and first cohomology vanish.  On each of the twelve
nodal fibres the line has degree zero and a nonidentity
multiplicative Picard parameter.  The normalization sequence is
a two-term complex
\(\mathbf C\xrightarrow{1-\lambda}\mathbf C\), where
\(\lambda\ne1\); it is acyclic.  The bundle \(nF\) is pulled
back from the base and does not affect any of these restrictions.
All source identities and root-chart substitutions are checked
exactly in \path{code/part3_fourfold/direct_section.py} and its recorded
outputs.
Proper cohomology and base change give the vanishing below.
More precisely, Stacks Project Lemma 30.22.1, tag 07VK,
applies to this proper flat projective elliptic fibration and
its line bundle: on each Noetherian affine base open it provides
a perfect complex computing the direct image and commuting with
arbitrary derived base change.  Every fibre complex is zero by
the preceding calculation.  A bounded complex of finite projective
modules with this property is zero: at any local ring, cancel
every differential entry that is a unit; the remaining minimal
complex has zero differentials after reduction to the residue
field, so its zero fibre cohomology forces every remaining rank
to be zero.  This proves

\[
 R f_*(M+nF)=0.
\]

Consequently the Dolbeault index is zero.  Combining the actual
symbol factorization and cap comparison gives

\[
 A=\operatorname{ind}(D_{\bar X,b})
  =\operatorname{ind}\bar\partial_{S,M+nF}=0.
\]

The computation uses fibrewise acyclicity and the explicit cap
switch; it does not reuse the numerical page characteristic
calculation in P3K.7.  The integral cover argument P3F.11 again
gives zero for the original parameter return.
\end{proof}

\begin{statement}{P3K.9}{Proved}
For the rotating framing specified below, the one-gamma-turn
return has Jin characteristic numbers \((k,\ell)=(-6,-24)\).
After removing an extending constant-rotation loop, this framing
descends through the sixfold parameter cover.  With the descended
framing the original positive return has

\[
 k=-1,\qquad \ell=-4.
\]

Every other tubular framing differs by the explicitly described
changes \((k,\ell)\mapsto(k+n,\ell+4n)\) of P3F.11.  The asserted
individual values include the framing; their index is independent
of that choice.
\end{statement}

\begin{proof}
\par\medskip\noindent\textbf{Frames on the two pieces.}\enspace Give the graph copy of the page
its normal frame consisting of its oriented normal in \(Y\),
followed by the two Cartesian parameter directions.  The graph
derivative changes this frame only by a triangular shear, whose
off-diagonal block contracts linearly to zero.  On the compression
use a normal frame of \(W\subset Y\times I\), then transport it
by the actual gamma rotation \(R_\theta\).  This normal rank-three
bundle is trivial.  Indeed \(W\) is a solid torus minus a ball and
has trivial tangent bundle, while \(TY\oplus\mathbf R\) is trivial;
the Whitney formula gives \(w_2(\nu W)=0\).  Over a three-complex,
an oriented rank-three bundle is classified by \(w_2\), since
\(\pi_1(SO(3))=\mathbf Z/2\) and \(\pi_2(SO(3))=0\).
At the sphere end, its comparison with the fixed positive-Hopf
frame is a map \(S^2\to SO(3)\).  This extends over \(D^3\) because
\(\pi_2(SO(3))=0\).  Make this extension at parameter zero and
transport it by \(R_\theta\).  The resulting frame extends over
the whole marked compression cap \(H_a\).

We compute the transition between these two actual frames.  Use
radian angular coordinates in this paragraph.  Let \(\phi\) be
the fixed outward-collar phase, namely the Milnor-page phase plus
\(\pi\), and let \(d=12\) be its degree along the gamma circle.
At the binding its outward page ray has angle
\(\alpha=\phi-d\gamma\) in the complex \(q-\epsilon\) plane.
The exact graph boundary from P3K.6 is
\(\gamma=\theta+g_\epsilon(p)\), where
\(p\in T^2\) is the source point and \(g_\epsilon\) is now the
radian-valued lift of that graph.  No source-dependent change of
the binding phase is made.

After quotienting by the source tangent directions and the tangent
to this boundary graph, the four remaining coordinates are the
real and imaginary \(q-\epsilon\) directions, the parameter radial
direction, and the relative angular direction.  The chosen page
collar has positive parameter radial derivative, pointing from the
page towards the compression.  We use the outward page collar and
the inward compression tangent with this same direction.  In the
four indicated coordinates its tangent is
\(u=(\cos\alpha,\sin\alpha,\rho,0)\), with \(\rho>0\).
The explicit quotient morphism, from this four-space modulo
\(\mathbf Ru\) to the reference three-space, is

\[
 (v_q,v_r,v_\eta)\longmapsto
 \left(v_q-\frac{v_r}{\rho}(\cos\alpha,\sin\alpha),v_\eta\right).
\]

The Cartesian parameter axes have radial-relative coordinates
\((\cos\theta,-\sin\theta)\) and
\((\sin\theta,\cos\theta)\), after a positive scaling of the
relative angular coordinate.  Project them by this formula,
reverse both parameter axes, and apply the positive reference
rescaling \(\operatorname{diag}(\rho,\rho,1)\); normalize the
page-normal column by the positive factor \(\rho^{-1}\).
The graph normal frame in
this reference space, after constant oriented changes of frame and
positive rescalings, is

\[
 G(\alpha,\theta)=
 \begin{pmatrix}
 -\sin\alpha&\cos\theta\cos\alpha&\sin\theta\cos\alpha\\
 \cos\alpha&\cos\theta\sin\alpha&\sin\theta\sin\alpha\\
 0&\sin\theta&-\cos\theta
 \end{pmatrix}.
\]

Here the first column is the oriented page normal, and the other
two columns are the Cartesian parameter directions projected along
the collar tangent.  The coefficients of the source derivatives
of \(g_\epsilon\) enter only the source tangent directions already
quotiented out; they give triangular shears before this quotient.
The positive radial collar coefficient can be rescaled to one.
During collar rounding the transverse vector remains in the same
open radial half-space.  This space is contractible, so this
rounding contributes no additional three-dimensional clutching
class.  These operations specify a homotopy of the actual normal
transition, not merely an equality of its characteristic numbers.

Direct multiplication gives \(G\in SO(3)\) and
\(G(\alpha,\theta)=R_z(\alpha)G(0,\theta)\), where \(R_z\)
rotates the first two reference axes.  The compression-cap frame
in these same coordinates has the form \(h(p)\), independent of
\(\theta\): both the ambient product coordinates and this frame
are transported by \(R_\theta\).  Consequently the full transition
on the marked \(T^3\) is homotopic to

\[
\begin{aligned}
 h(p)^{-1}G(\phi-d(\theta+g_\epsilon(p)),\theta)
 &=\bigl[h(p)^{-1}R_z(-d g_\epsilon(p))\bigr]\\
 &\quad\cdot\bigl[R_z(\phi-d\theta)G(0,\theta)\bigr].
\end{aligned}
\]

The first factor is pulled back from the source \(T^2\), and the
second from the parameter circle.  The cubic Maurer--Cartan form
of a product of two maps is the sum of their cubic forms plus an
exact form; this identity follows by expanding
\((ab)^{-1}d(ab)=b^{-1}(a^{-1}da)b+b^{-1}db\) and using the
Maurer--Cartan equations.  Each of the two cubic forms here is
zero by dimension.  Its integral over \(T^3\) is therefore zero.
The Pontryagin clutching formula can be obtained directly by using
the interpolating connection \(t g^{-1}dg\) on the collar between
two framed pieces: its Chern--Weil integral is a fixed multiple
of that cubic integral.  We conclude

\[
 \langle p_1(\bar\nu),[\bar X]\rangle=0.
\]

This calculation allows nontrivial one-dimensional and
two-dimensional frame data.  It proves the required vanishing
without replacing the actual normal bundle by a trivial bundle.
The matrix identities and ensuing exact arithmetic are exported
in \path{code/part3_fourfold/normal_clutch.py} and its JSON and text outputs.

\par\medskip\noindent\textbf{The integers and their descent.}\enspace P3K.7 now gives, with
its full orientation comparison,

\[
 k_{\gamma}=\frac{0+3(-8)}4=-6,
 \qquad \ell_{\gamma}=-24+0=-24.
\]

We must compare the rotating framing to framings pulled back from
one original parameter circuit.  The fundamental group of the
normal-frame gauge space is

\[
 \pi_1\operatorname{Map}(S^2,SO(3))
 =\pi_1(SO(3))\oplus\pi_3(SO(3))
 =\mathbf Z/2\oplus\mathbf Z.
\]

Evaluation at one point, with its constant-map section, gives the
displayed splitting.  Pullback through the positive sixfold circle
cover sends \((\varepsilon,n)\) to \((0,6n)\).  Choose any
framing over the original parameter circle and write \(k_0\) for
its Jin integer.  Its sixfold pullback has integer \(6k_0\).
By the exact framing-to-return morphism P3F.11, the free part of
the difference between this pullback and the rotating framing is
\(-6-6k_0\), a multiple of six.  The remaining possible
\(\mathbf Z/2\) part is represented by a constant-axis rotation
loop, independent of the fixed-sphere point.  It extends over the
fixed \(D^3\), contributes zero to the cubic clutch integral, and
changes neither characteristic number.  Remove that extending
loop.  The resulting rotating framing is the sixfold pullback,
up to homotopy, of a framing over the original circle.  We take
this as the descended rotating-framing convention.

The replacement diffeomorphism for the pullback is the sixth
iterate of that original framed replacement.  Jin's integral
mapping-torus homomorphism, in the convention of P3F.11, gives
\(6k=-6\) and \(6\ell=-24\).  This proves the stated
\((k,\ell)=(-1,-4)\), including the exact descent of the
normalization.
\end{proof}

\begin{statement}{P3K.10}{Proved}
Let \(S\) be the minimal regular elliptic surface with section \(O\)
whose affine equation is

\[
 y^2=x^3-3f(t)x+2g(t),\qquad
 f=(t^2+1)^2+t,\qquad
 g=(t^2+1)^3+\frac32t(t^2+1)+\frac38.
\]

Put \(G=(t^2+1,\sqrt3/2)\), \(R=8G\),
\(M=\mathcal O_S(R-O)\), and \(Q=G\).  Choose any one of the four
nodal fibres as the simple zero of the linearization.  At the
ordinary fibre \(t=0\), perform the multiplicity-one logarithmic
gluing whose integral circle clutch is \(\widehat\gamma\).
The resulting construction of Engel's Remark 8.5 gives a proper
holomorphic family of compact complex manifolds diffeomorphic to
the standard \(S^6\), over a punctured scalar disc
\(0<|\lambda|<c\).

For the positive \(\lambda\)-loop, with the complex orientation and
the Eells--Kuiper convention of P3F.22, the return satisfies

\[
 A=-3,\qquad \mu(\Sigma_\Phi)=-\frac3{28},\qquad
 [\Phi]=25\in\mathbf Z/28.
\]

In particular this is a generator.  The positive finite parameter
pullbacks of this same family realize every oriented class in
\(\Theta_7\).  For every nonzero such class the mapping torus is
not diffeomorphic to \(S^1\times S^6\).
\end{statement}

\begin{proof}
\par\medskip\noindent\textbf{The source and its section.}\enspace More generally put
\(R_0=t^2+a\),
\(f=R_0^2+bt\), and
\(g=R_0^3+\frac32btR_0+\frac38b^2\).  Exact expansion gives

\[
 g^2-f^3=
 \frac34b^2a(t^2+a)^2+
 \frac18b^3t^3+\frac98b^3at+\frac9{64}b^4.
\]

At \(a=b=1\), this is
\((48t^4+8t^3+96t^2+72t+57)/64\).  Its four roots are simple
and none is a root of \(f\).  At infinity the minimal coefficients
are units and the discriminant has order eight.  Thus the fibres
are exactly \(I_8+4I_1\).  The point \(G\) satisfies the equation
by direct substitution.

Rationality can be checked without identifying the surface by a
table entry.  Set \(u=x-(t^2+a)\).  The equation becomes

\[
 y^2=3(ut-b/2)^2+u^2(u+3a).
\]

On \(u\ne0\), the two rational functions

\[
 A_0=\frac{y-\sqrt3(ut-b/2)}u,
 \qquad B_0=\frac{y+\sqrt3(ut-b/2)}u
\]

satisfy \(A_0B_0=u+3a\).  Conversely,

\[
 u=A_0B_0-3a,\qquad
 t=\frac{B_0-A_0}{2\sqrt3}+\frac b{2u},\qquad
 y=\frac u2(A_0+B_0),\qquad x=t^2+a+u.
\]

These are inverse rational maps on nonempty open sets.  The smooth
minimal model is therefore a rational elliptic surface.

The section \(G\) is disjoint from \(O\).  At infinity, with
\(z=1/t\), its minimal coordinates are
\(X=1+az^2\), \(Y=(\sqrt3b/2)z^3\).
The two local nodal branches have orders three and five when
\(a=b=1\): expanding the square root of
\((1+az^2)^2+bz^3\) cancels the order-three term on exactly
one branch, whose next nonzero term has order five.  Their product
has order eight.  Hence \(G\) meets component three of \(I_8\),
up to reversing the cyclic numbering.  Its height is

\[
 h(G)=2-\frac{3(8-3)}8=\frac18.
\]

The root lattice has rank seven and determinant eight.  Shioda--Tate
therefore gives Mordell--Weil rank one.  The unimodular surface
lattice and the height discriminant formula give
\(h(G)=m^2|T|^2/8\), where \(m\) is its index in the free
Mordell--Weil group and \(T\) its torsion subgroup.  Thus
\(m=|T|=1\).  This identifies the actual generator and the
torsion-free \(I_8\) source, including the component marking.

The section \(R=8G\) meets the identity component at \(I_8\).
Its height is eight, so the same intersection formula gives
\(R\cdot O=3\).  Exact doubling locates all three intersections:
the reduced denominator of \(x(8G)\) is

\[
 (6t+1)^2(360t^2+282t+199)^2.
\]

The three roots are distinct and this denominator is coprime to
the discriminant.  The numerator is coprime to it as well, as
checked in the saved exact group-law calculation.  Consequently
all three intersections are ordinary and simple.  In particular
the Picard point of \(M\) is nonidentity at every \(I_1\) and
at the smooth identity component of \(I_8\).

\par\medskip\noindent\textbf{The actual marked local system.}\enspace We give a comparison
with the displayed source, rather than merely exhibit matrices
having the correct local spectra.  At \(a=0,b=1\) the source has
\(I_9+3I_1\).  Its \(j\)-map has degree twelve and ramification
profiles \(3^4\), \(2^6\), and \((9,1,1,1)\) over
\(0,1728,\infty\).  Indeed \(f=t(t^3+1)\) has four simple
roots, \(g=t^6+\frac32t^3+\frac38\) has six simple roots,
and \(g^2-f^3=(8t^3+9)/64\); these polynomials are pairwise
coprime.  Their indicated ramification already totals twenty-two,
so Riemann--Hurwitz leaves no further branch points.

The associated trivalent modular graph is unique.  It has four
vertices and six edges, with three monogon faces.  Each monogon
uses a loop at a vertex and consumes two of its three half-edges.
The three loops therefore occupy three distinct vertices, whose
remaining half-edges must join the fourth vertex.  This is the
three-star with a loop on each leaf.  Its cyclic ordering is unique
up to orientation-preserving graph isomorphism.  This proves
uniqueness of the permutation cover with these profiles by the
usual path-lifting classification of covers of the three-punctured
sphere.

For \(v=(p,q)^t\), write

\[
 D_v=\begin{pmatrix}1-pq&p^2\\-q^2&1+pq\end{pmatrix},
 \qquad T=D_{(1,0)}.
\]

The subgroup generated by
\(T^9,D_{(-3,1)},D_{(0,1)},D_{(3,1)}\) has the same cover:
its twelve-coset permutation table, exported by the program, has
exactly the three profiles above.  The products of the displayed
matrices equal the identity.  Thus the modular covering comparison
identifies this tuple with the actual \(I_9+3I_1\) source, up to
the explicitly permitted choice of marked vanishing paths and a
single integral symplectic basis.  The lift from projective
monodromy is fixed by the trace-two local \(I_n\) matrices.

Near infinity the deformation in \(a\) is analytically
\(xy=z^8(z+c a+O(a^2,az))\), with \(c\ne0\), after absorbing
units.  Hence it splits the \(I_9\) into an \(I_8\) and an
\(I_1\) with the same vanishing cycle.  Transport to \(a=1\)
along a path avoiding the finitely many other discriminant
collisions gives the actual five-point tuple

\[
 E_0=T^8,\quad E_1=T,\quad E_2=D_{(-3,1)},\quad
 E_3=D_{(0,1)},\quad E_4=D_{(3,1)}.
\]

Such transport is an isomorphism of the integral local systems
over the punctured surfaces and carries the section \(G\).
The path choice only changes the marking by the corresponding
Hurwitz moves.  No assertion about the topology depends on choosing
a particular branch of this transport: the program verifies all
four choices of the nodal zero below.

In the cohomology convention \((\psi,u,w,\delta)\) of P3F.21,
the weighted Jacobi lift for \(M=8G-O,Q=G\) is

\[
 J(E,k,z)=
 \begin{pmatrix}
 1&cj-ah&dj-bh&z\\
 0&a&b&8j\\
 0&c&d&8h\\
 0&0&0&1
 \end{pmatrix},
 \quad E=\begin{pmatrix}a&b\\c&d\end{pmatrix},
 \quad k=\binom jh.
\]

The component of \(G\) fixes \(k_0=(3,0)^t\).  An integral
change of Abel--Jacobi lift also sets \(k_4=0\).  Here is the
simultaneous gauge explicitly.  Adding the common integral lift
\((a,b)^t\) changes the first cocycle by
\((E_0-I)(a,b)^t=8b(1,0)^t\), so \(b\) fixes its component-three
representative.  The fourth coefficient changes by \(-a+3b\);
varying \(a\) then sets it to zero without changing the first.
The remaining
cocycles have the forms
\(k_1=r(1,0)^t\), \(k_2=j(-3,1)^t\),
\(k_3=h(0,1)^t\).  The nonzero-linearization local fillings
have rank-one nilpotent image; this forces their central entries
to be \(9,8r^2,8j^2,8h^2\), respectively.  At the scalar-zero
\(I_1\), the central entry is \(-1\).  The actual global product
must be the identity.  Its source columns give
\(j=2h\), \(r=-3h-3\), and its central entry gives

\[
 16(h+1)(4h+5)=0.
\]

The unique integral solution is \(h=-1\).  This proves the
actual section-lift identification; it is not a bounded search
for a convenient tuple.  The resulting matrices are

\[
\begin{aligned}
 T_0&=\begin{pmatrix}1&0&3&9\\0&1&8&24\\0&0&1&0\\0&0&0&1\end{pmatrix},
 &T_1&=\begin{pmatrix}1&0&0&0\\0&1&1&0\\0&0&1&0\\0&0&0&1\end{pmatrix},\\
 T_2&=\begin{pmatrix}1&2&6&32\\0&4&9&48\\0&-1&-2&-16\\0&0&0&1\end{pmatrix},
 &T_3&=\begin{pmatrix}1&1&0&8\\0&1&0&0\\0&-1&1&-8\\0&0&0&1\end{pmatrix},\\
 T_4&=\begin{pmatrix}1&0&0&-1\\0&-2&9&0\\0&-1&4&0\\0&0&0&1\end{pmatrix}.
\end{aligned}
\]

Moving the linearization zero to any of the other three \(I_1\)
fibres adds one to the central entry of \(T_4\) and subtracts
one at the chosen fibre.  The central matrices commute, so the
global product is still the identity.  Every integer calculation
below has been executed for all four choices.

\par\medskip\noindent\textbf{Line, fillings, and the ordinary logarithmic gluing.}\enspace
The line \(M\) has multidegree zero on every component of
\(I_8\).  Translation by \(Q\) therefore preserves its
multidegree.  The vertical line
\(\operatorname{Hom}(t_Q^*M,M)\) has zero intersection with
every fibre component and is a whole-fibre pullback.  Its degree,
computed by restriction to \(O\), is

\[
 (8G)\cdot O+1-(7G)\cdot O+(-G)\cdot O
 =3+1-3+0=1.
\]

Thus it is \(\mathcal O(F)\), and its section may have its
sole simple zero at any chosen nodal fibre.  Scaling the induced
lift by sufficiently small \(\lambda\) makes it uniformly
contracting.  On the complement of that fibre its quotient is a
smooth Hausdorff complex threefold, as follows directly from the
norm contraction and a compact annular fundamental domain.

At \(I_8\), the Tate and suspension periods, in a multidegree-zero
line frame, have valuation vectors \((8,0)\) and \((-3,0)\).
If they are denoted by \(h\) and \(g\), respectively, take

\[
 H=h^{-1}g^{-3},\qquad K=h^3g^8.
\]

The change has determinant one and their first valuations are
\(1,0\).  The second multiplier of \(K\) is \(\lambda^8\)
times a nonvanishing unit.  Thus the ray fan with consecutive
vertices \((j,1)\) is the actual primitive rank-one filling.
Equivalently it is the contracting quotient of the smooth source
line complement; \(Q\) cyclically permutes all eight components.
This equivalence retains the source elliptic map and its section
line.  The same primitive rank-one model applies at the other
three nonzero-linearization nodal fibres.  At the chosen zero
the two valuations are \((1,0),(0,1)\).  The square triangulation
has unimodular cones and gives the smooth rank-two cusp with
charts \(q=z_0z_1z_2\), its toric period identifications, and
proper reduced central fibre.  These are the specified local
models, with all valuation changes integral.

The base point \(p=0\) is ordinary.  On its smooth torus disc,
and on the universal logarithm cover of the scalar parameter,
let \(\gamma(z,v)\) be the holomorphic period of the primitive
suspension cycle. Glue a second copy of this smooth torus disc
by the following punctured-disc isomorphism, directed from the
cap copy to the ordinary copy:

\[
 \Psi_v(\zeta,z)=
 \left(\zeta+\frac{\log z}{2\pi i}\gamma(z,v),z\right),
 \qquad \lambda=e^{2\pi i v}.
\]

A logarithm branch change adds an actual lattice period, so this
is a single-valued holomorphic map on the torus family, with
inverse given by subtraction.  The cap is smooth, its central
fibre has multiplicity one, and its boundary clutch is exactly
\(\widehat\gamma\). Indeed a cap meridian with fixed cap coordinate
is carried to the ordinary base meridian plus \(\gamma\), the
\(m=1,n=1\) marking used below. This verifies the allowed logarithmic
transformation at an ordinary fibre even though its multiplicity
is one.  All pieces are proper over their closed base discs and
agree by the displayed holomorphic maps on collars; their finite
gluing is a compact Hausdorff complex threefold.

\par\medskip\noindent\textbf{Integral topology.}\enspace Put \(H_i=T_i^{-t}\) on the
homology lattice \(\Lambda\).  The images of \(H_i-I\) have
Smith factors \((1),(1),(1),(1),(1,1)\).  The primitive toric
collapses kill exactly these vanishing lattices.  We first fix
an actual global origin for the auxiliary untwisted family, before
the ordinary logarithmic transformation.  The restriction
\(M|_O\) has degree \(R\cdot O-O^2=4\).  Choose its section
with divisor \(4t_c\).  Its nonzero values over the other base
points give an origin in the quotient.  Near \(t_c\), its line
character is \(V=q^4\) times a holomorphic unit.  Subtracting
four gamma periods changes this to the nonvanishing bounded
character \(\lambda^{-4}e^{-8\pi iC}\) times a unit and changes
the source coordinate by the bounded translation \(-4\Gamma\).
It therefore extends in the dense torus of the cusp component.
This proves that the origin is global, including all five singular
fibres.  Every base meridian based at this origin contracts in
its local filling.

Van Kampen begins with the affine torus-bundle group over the
punctured sphere.  After these base meridians are killed, all
generators lie in the abelian fibre lattice, and the conjugation
relations impose the images of \(H_i-I\).  Their stacked relation
matrix has Smith factors \((1,1,1)\) and quotient \(\mathbf Z\)
measured by \(\psi\).  The ordinary log cap adds a vector with
\(\psi\)-value one.  The augmented Smith factors are
\((1,1,1,1)\), proving \(\pi_1=1\), with every affine origin
correction retained in the kernel already killed.  In the chosen
global-origin marking the complete relations are the abelian
fibre relations, \((H_i-I)\Lambda=0\) for \(0\leq i\leq4\),
and \(\widehat\gamma=0\).  Changing an origin only adds a vector
of \(\ker\psi\), which is precisely the primitive subgroup
generated by those five image lattices.

The complete cohomology calculation is integral.  For degree \(q\)
let \(W_i=\bigwedge^qT_i\), let \(L_i\) be the integer column
lattice of \(W_i-I\), and use its Hermite basis.  The marked
parabolic complex is

\[
 \mathbf Z^{\binom4q}\xrightarrow{d^0}
 \bigoplus_{i=0}^4 L_i\xrightarrow{d^1}
 \mathbf Z^{\binom4q},
\]

where \(d^0(a)=((W_i-I)a)_i\) and
\(d^1((b_i))=\sum_i(W_0\cdots W_{i-1})b_i\).
This is obtained by cutting the marked punctured sphere and
attaching each invariant stalk; the product relation gives
\(d^1d^0=0\).  The code exports every matrix in every degree.
For degrees one, two, and three, their ranks are respectively
\((3,3),(5,5),(3,3)\), and every nonzero Smith factor of
both differentials is one.  The middle groups vanish integrally;
the end groups are free of rank one.

There are no finite quotient caps in this example.  At each
primitive rank-one filling the collapse is the cellular pinching
of one circle in \(T^2\), crossed with the other elliptic factor;
it identifies cohomology with the full invariant lattice.  At
the rank-two filling the two valuation vectors and triangulation
identify the model integrally with the cusp of Engel's Lemma 7.4.
For an independent local chain verification, use the actual
dual-moment-cell quotient and blowup-collapse comparison of P1.11
in rank two, with \(G=3I_2-\mathbf1\mathbf1^t\).  Its translated
simplex counts are \((1,3,2)\), with twelve codimension-one
incidences.  For a simplex \(\sigma\), put
\(Q_\sigma=\mathbf Z^2/L_\sigma\).  The differential from
\(\bigoplus_{\dim\sigma=p}\bigwedge^q Q_\sigma\) is precisely
the signed integral quotient matrix (1.11.1), now with
\(0\leq p,q\leq2\).  The program
\path{code/part3_fourfold/rank_two_cusp.py} exports all these matrices
and all their Smith transformations.  Their total homology is
free of ranks \((1,2,4,2,1)\).  The collapse is the block sum
of \(\bigwedge^q(\mathbf Z^2\to Q_\sigma)\); its induced
homology matrices are onto with only unit Smith factors.  The
only nontrivial nonidentity block on the associated graded groups,
in total degree two, is

\[
 \begin{pmatrix}0&0&-1&1\\1&0&0&0\\0&-1&0&1\end{pmatrix}.
\]

Support leaves no possible higher differential.  The finite
filtration therefore makes the actual collapse onto and primitive
in every degree.  Its dual image is monodromy invariant and
primitive, with the same rank as the full invariant lattice, so
the two lattices agree.  This separately verifies every local
specialization assertion of the cited lemma needed here.
These exact local comparisons
give \(R^qf_*\mathbf Z=j_*\bigwedge^q\Lambda^*\) everywhere.

The surviving invariant generators, in degrees zero through four,
are

\[
 1,\quad\psi,\quad\omega=8uw+\psi\delta,
 \quad\eta_3=\psi uw,\quad\eta_4=\psi uw\delta.
\]

The corresponding coinvariant generators are
\(1,[\delta],\eta_2,[uw\delta],[\eta_4]\), with
\([\psi\delta]=8\eta_2\), \([uw]=\eta_2\), all other
degree-two basis classes zero.  The only nonzero Leray columns
are zero and two.  There is no unaccounted intrinsic cusp
differential.  Inversion in the untwisted torus family about the
global origin just constructed extends over every local filling:
on primitive nodal charts
\(X_i=Z/q^i\), \(Y_i=q^{i+1}/Z\), it is the actual holomorphic
map \((X_i,Y_i)\mapsto(Y_i,X_i)\) in chart \(-i-1\).
The other elliptic coordinate is inverted as well.  On the
rank-two square fan the lattice map \((v,h)\mapsto(-v,h)\)
sends each cone to its centrally inverted cone; its dual monomial
maps are holomorphic chart isomorphisms and agree on overlaps.
These are integral automorphisms of the specified fans.
The bounded origin characters only add a toric unit translation,
which extends too.  Thus this is an involution of the actual
untwisted compact family.  It acts by \((-1)^q\) on \(R^q\).
Naturality of Leray then makes every untwisted \(d_2\) satisfy
\(d_2=-d_2\).  Its targets are the free groups just computed,
so every one is zero integrally.

Replacing the ordinary cap changes only its boundary translation.
The cellular prism for this actual translation changes \(d_2\)
by contraction with its integral clutch.  Consequently the
translation prism of the ordinary
logarithmic clutch gives
\(d_2(a)=H[\iota_{\widehat\gamma}a]\), where \(H\) is the
positive base class.  The four contractions are

\[
 [\iota_\gamma\psi]=1,\quad
 [\iota_\gamma\omega]=[\delta],\quad
 [\iota_\gamma\eta_3]=\eta_2,\quad
 [\iota_\gamma\eta_4]=[uw\delta].
\]

For an arbitrary integral affine origin correction, the same
calculation is exactly its \(\psi\)-coordinate; source and line
circle coordinates contribute zero to all four projections.
Thus every displayed differential is multiplication by one for
the actual glued family.  There are no higher differentials,
because the base has dimension two.  Only
\(H^0=H^6=\mathbf Z\) remains, without torsion or extension
ambiguity.  Together with simple connectivity this proves that
the total manifold is a homotopy six-sphere, and \(\Theta_6=0\)
identifies its smooth oriented type with the standard sphere.

\par\medskip\noindent\textbf{The closed positive parameter loop.}\enspace Write \(t_c\) for
the chosen nodal zero and put
\(h(t)=(t-t_c)/t\).  Since
\(\gamma(z,v+1)=\gamma(z,v)+\delta\), the smooth logarithmic
gluing changes by \((\log t)/(2\pi i)\delta\) under the
positive parameter deck.  The global ordinary deck is
\(-\log h(t)\delta/(2\pi i)\): near \(t=0\) it has exactly
that value up to a holomorphic unit logarithm, which extends on
the cap. More explicitly, write
\[
 D_{\rm ord}=-\frac{\log h(t)}{2\pi i}\delta,\qquad
 D_{\rm cap}=-\frac{\log(t-t_c)}{2\pi i}\delta.
\]
The second logarithm is holomorphic on the chosen smooth cap.
The directed gluing satisfies the exact compatibility
\[
 D_{\rm ord}\circ\Psi_v=\Psi_{v+1}\circ D_{\rm cap},
 \qquad
 D_{\rm ord}=D_{\rm cap}+\frac{\log t}{2\pi i}\delta .
\]
Here the \(D\)'s denote the respective translations from \(v\)
to \(v+1\). At the cusp use \(q=h(t)\), \(s=\log q/(2\pi i)\).
Then the actual periods are

\[
 \gamma(s,v)=\binom{\Gamma(q)}{C(q)+s+v},\qquad
 \delta=\binom01,
\]

with \(\Gamma,C\) holomorphic at zero, and the deck translates
by \(-s\delta\).  Its line-character multiplier \(q^{-1}\)
extends through the cusp: compose an inverse gamma lattice shift
with the bounded source translation \(\Gamma\) and the
nonvanishing multiplier \(\lambda e^{2\pi iC}\).  At every
other original fibre \(h\) is a unit.  This proves that the
glued family descends holomorphically over the punctured
\(\lambda\)-disc, with the stated positive closing map.

Here \(\gamma\) is the inward period of the actual contracting
suspension. In the physical Abel and line coordinates it is
\[
 B(a,V)=
 \bigl(a+\Gamma(q),\,e^{2\pi iC(q)}q\lambda V\bigr),
 \qquad \Gamma(q)=\log_{\rm Tate}(-G).
\]
For small \(q,\lambda\), its positive real direction decreases
\(\log|V|\). This fixes its orientation relative to the source
page. With \(\delta\) the positive phase of \(V\), the complex
orientation gives
\[
 \operatorname{or}(Y)
   =-\operatorname{or}(S)\,d\log|V|
   =\operatorname{or}(S)\,d\gamma .
\]
Thus the positive gamma in the Hopf page below is this inward
period. The normalized smooth norm coordinate increases in
this direction.

For a completely marked angular comparison use the auxiliary
logarithmic cover \(q=Q_0^2\), and the two origins

\[
 o=\binom{-\Gamma s}{-Cs-s^2/2-sv},\qquad
 \sigma=\binom{\Gamma v}{Cv+(v^2-v)/2}.
\]

The first has parameter difference \(-s\delta\); the second
has difference \(-s\delta+\gamma\).  Under \(s\mapsto s+2\),
their differences are \(-2\gamma-2\delta\) and zero.  Thus
they descend on this auxiliary cover.  In contrast, the ordinary
origin has one-step difference \(-\gamma-\delta/2\); we do
not impose a false one-step periodicity.  The cusp origin extends
by its bounded unit characters.  Subtracting the ordinary origin
removes the parameter deck and gives the one-node parameter chain
\(X_i=V/\lambda^i\), \(Y_i=\lambda^{i+1}/V\),
\(X_iY_i=\lambda\), over every nonzero \(q\).  Its origin
extends.  These are the two extending collar origins.  No proper
two-variable corner or descent of such a corner is used.

Their exact difference is

\[
 F=\binom{\Gamma(s+v)}{C(s+v)+(s+v)^2/2-v/2},
\]

and hence

\[
 F(s,v+\theta)-F(s,v)
 =\theta\gamma(s,v)+\frac{\theta(\theta-1)}2\delta.
\]

The additional line gauge is periodic, equals one at both
endpoints, and vanishes in the circle quotient.  The double cusp
deck changes this angular difference by \(2\theta\delta\),
so the retained line gauge has winding two.  In particular the
actual positive \(\lambda\)-loop makes one positive gamma turn,
with no parameter cover.  Every identity is executed in
\path{code/part3_fourfold/os45_origins.py}.

More explicitly, the ordinary and cusp coordinates are
\(\zeta-o\) and \(\zeta-\sigma\), so the cusp-to-ordinary
transition is translation by \(F=\sigma-o\). Its angular
increment is the positive translation displayed above. This
identifies the moving fixed sphere in the constant orbit with
the positive inward gamma motion used by the real trace.

\par\medskip\noindent\textbf{The actual page and the return index.}\enspace The line circle is
free away from the selected rank-two cusp.  Indeed its equality
with an iterate of the contracting lift forces that iterate to
be zero by the norm inequality, and then forces the circle
element to be the identity.  This includes the whole \(I_8\)
neighbourhood, where translation permutes the components.  The
ordinary logarithmic cap is free as well.  At the single cusp
the primitive fan gives the one fixed rational curve with normal
weights \((-1,+1)\).  The orbit-pair argument P3F.7 therefore
applies to this proved sphere: its quotient is the standard
\((S^5,S^2)\), and its Euler line has positive meridian one.

The gamma circle over the base has the one ordinary clutch of
degree one, so its total space is the standard Hopf \(S^3\).
For precision, in the two standard Hopf charts the normalized
sections \((1,t)/\sqrt{1+|t|^2}\) and
\((1/t,1)/\sqrt{1+|1/t|^2}\) have angular coordinates related
by \(\theta_\infty=\theta_0+\arg t\), exactly the clutch
above.  A regular Hopf binding has a complex-oriented disc page
whose projection to \(\mathbf P^1\) has degree one.  Its inverse
image in the actual orbit is the elliptic page \(V\), with
boundary \(T^3\); its regular cap is the displayed rational
elliptic surface \(S\).  The affine origin obstruction on the
disc with marked singular values vanishes by the free-meridian
cochain argument of P3F.10.  Thus this comparison is with the
actual page, including the retained \(I_8\), rather than with
an unrelated elliptic surface.

The actual circle line restricts to \(M=8G-O\), up to a base
pullback, by its fibrewise Picard point.  At \(I_8\) both lines
have multidegree zero, so the same normalized fibrewise comparison
extends across its cycle of components.  On the base with boundary
the pullback is topologically trivial.  Therefore the actual
relative page invariants are

\[
 s_\delta=M^2=-2-2(3)=-8,
 \qquad p_{\rm rel}=3\sigma(S)=-24.
\]

Near the selected \(I_1\) the source translation is in the smooth
identity component.  The product coordinate of P3F.9 and the real
phase-preserving compression of P3K.6 consequently apply with
page degree \(d=1\).  The page phase can be chosen to avoid the
short gamma arc of that compression.  Its exact graph boundary,
integral Euler-line restriction, and oriented complement bordism
are therefore the ones proved in P3K.6–7.  The angular computation
above identifies this bordism with the actual positive scalar
return.  With the two orientation reversals of P3K.7 retained,

\[
 A=\frac{12s_\delta-p_{\rm rel}}{24}
   =\frac{-96+24}{24}=-3.
\]

The typed Hopf transfer and Eells--Kuiper normalization P3F.22 give
the asserted oriented class \(25\).  Pullback by
\(\lambda=z^n\) is a proper holomorphic family on the punctured
disc and its positive loop is exactly the \(n\)-fold iterate.
The classes \(-3n\pmod{28}\), for positive \(n\), exhaust
\(\mathbf Z/28\).  The mapping-torus obstruction P3F.24 applies
to each nonzero class and to both spin structures.  This proves
all the assertions of the statement.
\end{proof}

\begin{statement}{P3K.11}{Proved}
Fix logarithmic orders \(3K,4L\), with
\(\gcd(3K,4L)=1\), and primitive logarithmic coefficients
\(n_0,n_1\), so \(\gcd(n_0,3K)=\gcd(n_1,4L)=1\).
For the sphere members impose
\(a=4L n_0+3K n_1=\pm1\), with zero additional cusp collar.
Use the period coordinate \(c\) and scalar \(u\) of the original
family.  The closed higher parameter is
\[
 w_h=\exp\!\left(-\frac{\pi i c}{KL}\right),
 \qquad u=\kappa w_h^{2KL},\qquad \kappa\ne0.
\]
In particular, the original expression \(w=\exp(-\pi i c)\)
does not give one closed higher loop when \(KL>1\).
The positive sixfold cover \(w_h=r^6\) gives the signed
gamma turn \(a\gamma\).  For the positive-gamma sphere marking
\(a=1\), with the descended rotating framing of P3K.9, the Jin
integers of the positive \(w_h\)-loop are
\[
                  (k,\ell)=(-KL,-4KL).
\]
The integer return index is zero for either sphere sign.
These statements identify
both the parameter and the framing in the supplied higher-order
formulas.
\end{statement}

\begin{proof}
\par\medskip\noindent\textbf{The integral orbifold complex.}\enspace We use the marked
lattice and period convention of P5.6.  Its cohomological
monodromy matrices are
\[
 T_1=\begin{pmatrix}
 1&0&-6&2\\0&-1&1&1\\0&-1&0&1\\0&0&0&1
 \end{pmatrix},\qquad
 T_2=\begin{pmatrix}
 1&6&0&-3\\0&0&-1&1\\0&1&0&0\\0&0&0&1
 \end{pmatrix}.
\]
Put \(A_i=T_i^{-t}\), \(A_0=(A_1A_2)^{-1}\),
\(N_1=1+A_1+A_1^2\), and
\(N_2=1+A_2+A_2^2+A_2^3\).  The lattice local system
on the higher root base is the pullback of the original one.
Consequently its finite cyclic norm operators are exactly
\(K N_1\) and \(L N_2\).  This assertion retains the
ineffective stabilizers: the generator of the additional
order-\(K\) or order-\(L\) stabilizer acts as the identity
on the lattice, so the full cyclic sum repeats the original
sum precisely \(K\) or \(L\) times.

Here is a complete block specification of the integer complex
used in the computation.  Let
\[
 C^0=\Lambda\oplus\Lambda\oplus\Lambda\oplus\mathbb Z^2,
 \quad C^1=\Lambda^7,\quad C^2=\Lambda^5,
 \quad C^3=\Lambda^2,
 \qquad \Lambda=\mathbb Z^4,
\]
and let \(J:\mathbb Z^2\to\Lambda\) have columns
\(e_3,e_4\).  In the displayed ordered summands, the three
differentials are
\[
\begin{split}
 d^0(a,b,c,z)={}&\bigl((A_1-1)a,(A_2-1)a,
 (A_1-1)b,(A_2-1)c,\,a-b,a-c,a-Jz\bigr),\\
 d^1(p,q,r,s,u,v,w)={}&\bigl(KN_1r,LN_2s,
 p-r-(A_1-1)u,\ q-s-(A_2-1)v,\\
 &\hspace{39mm}-A_0p-A_0A_1q-(A_0-1)w\bigr),\\
 d^2(x,y,z_1,z_2,z_0)={}&\bigl((A_1-1)x,(A_2-1)y\bigr).
\end{split}
\]
The pair-of-pants and the three annuli give the last three
summands of \(C^2\).  The first two are the degree-two terms
of the finite cyclic resolutions; the cusp contributes its
invariant lattice \(J\mathbb Z^2\) in degree zero.  These
are precisely the gluing maps of the lattice sheaf, including
the cyclic degree-two terms.  In particular, no finite norm
relation is replaced by a rational invariant projector.
Direct multiplication gives \(d^1d^0=d^2d^1=0\).

Write
\[
 \delta=(0,0,0,1)^t,\quad
 \epsilon_1=(1,2,-4,0)^t,\quad
 \epsilon_2=(1,3,-3,0)^t.
\]
The two invariant lattices have primitive bases
\(B_i=(\delta,\epsilon_i)\), and the norm matrices in
these bases are
\[
 B_1^{-1}N_1=
 \begin{pmatrix}0&2&1&3\\3&0&0&0\end{pmatrix},
 \qquad
 B_2^{-1}N_2=
 \begin{pmatrix}0&2&2&4\\4&0&0&0\end{pmatrix}.
\]
Here \(B_i^{-1}\) denotes the inverse coordinate map on its
image, so these displayed expressions are integral maps
\(\Lambda\to\mathbb Z^2\).  A primitive basis of
\(\ker d^2\) is \(\operatorname{diag}(B_1,B_2,I_{12})\).
The exact computation eliminates twelve unit pivots in the
matrix of \(d^1\) in that basis.  Every operation is an
elementary operation over \(\mathbb Z[K,L]\); the exported
left and right matrices have determinant \(\pm1\).
The remaining matrix has the same image as the three columns
\[
 R_0=
 \begin{pmatrix}
 K&0&0\\0&0&-3K\\0&2L&0\\0&0&4L
 \end{pmatrix}.
\]
This last image equality is also certified integrally: the
export gives matrices \(C,Q\) with
\(R_{\rm tail}=R_0C\), \(CQ=I_3\), and
\(R_{\rm tail}Q=R_0\).  Thus it is a proof for symbolic
\(K,L\), in addition to the executed Smith forms.
Equivalently, adjoining the eliminated generator \(h_c\),
the presentation is
\[
\begin{gathered}
 K\delta_1=0,\qquad 2L\delta_2=0,\\
 3K\epsilon_1=h_c=4L\epsilon_2.
\end{gathered}
\]
Since \(\gcd(3K,4L)=1\), the last relation gives a free
rank-one group and no torsion.  Moreover \(K\) is odd and
\(\gcd(K,L)=1\).  Therefore
\[
 H^2(\widehat B,\mathcal H)
 =\mathbb Z\oplus\mathbb Z/K\oplus\mathbb Z/(2L)
 =\mathbb Z\oplus\mathbb Z/(2KL).
\]
The complete matrices and elementary transformations are in
\path{code/part3_fourfold/higher_parameter.json}.
The same script computes the full integer complex, rather than
only this reduced presentation, for
\((K,L)=(1,1),(1,2),(1,5),(3,1),(5,1),(5,2),(7,5),(13,7)\).
Their torsion orders are respectively
\(2,4,10,6,10,20,70,182\), with free rank one in every case.

\par\medskip\noindent\textbf{The actual torsor period.}\enspace In the fixed period
coordinates, changing \(c\) to \(c'=c-b\) is the integral
shear
\[
 S_b=I+bE_{41},\qquad
 \Pi_c=\Pi_{c'}S_b,\qquad b\in\mathbb Z.
\]
It fixes \(\delta\) and sends each \(\epsilon_i\) to
\(\epsilon_i+b\delta\).  The logarithmic vectors are
\(r_i\delta+n_i\epsilon_i\).  Their actual logarithmic
annulus cocycles therefore change the characteristic class
by
\[
             b(n_0\delta_1+n_1\delta_2).
\]
This follows either by inserting the two vectors in the
displayed cochain complex, or by taking the logarithms in the
gluing maps \(x\mapsto x+\log(s_i)v_i/(2\pi i)\).
The coefficient \(r_i\) is fixed by the shear and contributes
zero to the difference.  Primitivity implies
\(\gcd(n_0,K)=\gcd(n_1,2L)=1\), so the displayed torsion
element has order exactly \(2KL\).  A change of logarithm
branch changes its cochain representative by a coboundary
and leaves this conclusion unchanged.

For completeness, the zero of this characteristic obstruction
is realized by actual holomorphic translations.  The vector
bundle of Lie algebras on the higher root base is the pullback
of the original bundle.  The added stabilizers act trivially
on its fibres.  Taking their invariants is exact in
characteristic zero, and the coarse pushforward is unchanged.
Consequently the explicit original calculation
\(H^1(\widehat B,\mathcal V)=H^2(\widehat B,\mathcal V)=0\)
continues to hold.  In the exponential sequence
\(0\to\mathcal H\to\mathcal V\to\mathcal V/\mathcal H\to0\),
the connecting map from torsors to \(H^2(\mathcal H)\) is
therefore an isomorphism.  When \(2KL\mid b\), solve the
integer coboundary equation above and subtract its logarithms
on the overlapping charts.  The remaining vector-valued
cocycle is a coboundary by \(H^1(\mathcal V)=0\); exponentiating
its chartwise primitives gives the required translation
isomorphism.  Reversing these primitives is its inverse.
At the cusp the shear changes the bounded beta term by an
integer; its toric factors \(\exp(2\pi i C\lambda)\), fan,
and monomial transition maps are identical.  Thus these chart
translations extend through the actual filling, as in the
explicit sheaf construction of P5.6.  Conversely a marked
translation isomorphism makes the characteristic difference
zero, so \(2KL\mid b\).  This proves the asserted exact
period in the marked family, with an extending morphism in
both directions.

\par\medskip\noindent\textbf{The loop and its characteristic numbers.}\enspace The scalar
relation is \(u=\kappa\exp(-2\pi i c)\).  Quotienting the
marked parameter by its just-proved period \(2KL\) gives
\(w_h=\exp(-\pi i c/(KL))\), and consequently
\(u=\kappa w_h^{2KL}\).  Put \(d=12KL\),
\(u=\kappa r^d\), and \(v=\log r/(2\pi i)\).
We retain the sign of the actual origin calculation.  The
finite-cap parameter descent is translation by
\[
 H(t)\delta=-\frac{4L n_0\log t+3K n_1\log(t-1)}{2\pi i}\delta.
\]
At the cusp it is \((as+U(q))\delta\), where \(U\) is a
single-valued holomorphic unit logarithm.  Write the period
\(\gamma=(\Gamma,C-s-dv)^t\).  On \(q=Q^d\), the
actual extending origins are
\[
\begin{split}
 o&=\binom{-a\Gamma s/d}{a s^2/(2d)-aCs/d+v(as+U)},\\
 \sigma&=\binom{a\Gamma v}{vU+aCv-ad v^2/2+ad v/2}.
\end{split}
\]
Their parameter differences are respectively \(H\delta\)
and \(H\delta+a\gamma\).  Their \(Q\)-differences are
integral lattice periods because \(d\) is even.  The
bounded cusp expression for \(\sigma\), and the invariant
zero coordinate after subtracting \(o\), give precisely the
extending sections in the node charts used in P2.2--3.
This checks the same morphisms with \(12\) replaced by
\(d\); it retains the signed twist.  Subtracting gives
\[
 F=\sigma-o
 =a\binom{\Gamma(s+dv)/d}{C(s+dv)/d-(s+dv)^2/(2d)+dv/2},
\]
and hence
\[
 F(s,v+\theta)-F(s,v)
 =a\theta\gamma+\frac{ad}{2}\theta(1-\theta)\delta.
\]
The last term is a periodic unitary gauge with both endpoint
values zero.  Thus the quotient clutch is the signed gamma
loop \(a\gamma\).  Equating the scalar expressions gives
the exact sixfold map \(w_h=r^6\), independently of the sign.

For \(a=1\), the page calculation P3F.17 gives
\[
 s_\delta=-2KL,\qquad p_{\rm rel}=-24KL.
\]
The normal-frame calculation of P3K.9 applies with page degree
\(12KL\).  Its transition factors into a two-dimensional
factor and a circle factor, exactly as in that proof;
the cubic clutch integral remains zero.  Thus
\(p_1(\bar\nu)=0\), whereas the Euler determinant has
square \(b^2=4s_\delta=-8KL\).  The oriented sphere-bundle
comparison P3K.7 gives
\[
 k_\gamma=\frac{3(-8KL)}4=-6KL,
 \qquad \ell_\gamma=-24KL.
\]
To descend the framing, choose any framing over the
\(w_h\)-circle.  Its sixfold pullback differs from the
rotating framing by a class in
\(\pi_1\operatorname{Map}(S^2,SO(3))=\mathbb Z/2\oplus\mathbb Z\).
The free difference is \(-6KL-6k_0\), which is divisible
by six.  Remove the constant-axis torsion loop, which extends
over the fixed \(D^3\) and changes neither characteristic
number.  The remaining class is the sixfold pullback of an
integral framing class, exactly as proved in P3K.9.  The
resulting descended rotating framing therefore gives
\(6k=-6KL\) and \(6\ell=-24KL\).  This proves the
stated pair and \(A=(4k-\ell)/24=0\).  For \(a=-1\),
the just-computed clutch traverses the gamma loop inversely.
The actual real trace of P3K.6--7 applies to this source
page with the corresponding boundary orientation; the index
of the inverse loop is the negative of the forward index.
That forward index is zero since
\(12s_\delta-p_{\rm rel}=0\).  Consequently the covered
index is zero also for \(a=-1\), and integral additivity
implies that the original \(w_h\)-loop has index zero.
The displayed individual pair is assigned to the stated
positive-gamma marking, so it involves no implicit reversal
of a negative sphere datum.

The original expression \(w\) and the higher expression
\(w_h\) satisfy \(w=w_h^{KL}\).  A single positive
\(w\)-loop only advances \(c\) by its original period two;
the higher characteristic obstruction vanishes for that
advance exactly when \(KL=1\).  This explicitly accounts
for the otherwise different-looking cover counts.  Changing
the final normal framing still changes \((k,\ell)\) by
\(n(1,4)\), as in P3F.11; the displayed values use the
descended rotating convention throughout.
\end{proof}

\begin{statement}{P3K.12}{Proved}
In the rank-one Mumford model above a smooth elliptic source fibre,
every double curve of the central cycle is elliptic and fixed by the
line circle.  A smooth cyclic descent, with the actual faithful root
character on the base, retains an elliptic fixed curve.  The same
conclusion holds after potential good reduction of the source and
for any number of permuted double curves.  Such a descended local
piece cannot occur in a smooth circle action on \(S^6\).
\end{statement}

\begin{proof}
Let \(E\) be one of the double curves.  At a point of
\(E\) the smooth threefold has coordinates \((e,x,y)\), with
\(e\) a local elliptic coordinate and root-base map
\[
                         z=xy.
\]
The line circle acts by \((e,x,y)\mapsto(e,\tau x,\tau^{-1}y)\).
For a higher vanishing order the refined Mumford chain has these
same charts at each double curve.  A finite root deck transformation
commutes with this circle, so a stabilizer at the chosen point
preserves the two normal weight lines separately.  A nonidentity
element of the stabilizer acts faithfully on \(z\), since it is
a subgroup of the cyclic root group acting by multiplication on
the root coordinate.

Write a cyclic stabilizer of order \(h\) on its tangent space as
\[
 (e,x,y)\longmapsto
 (\zeta^{a}e,\zeta^{b}x,\zeta^{c}y),
 \qquad b+c\equiv1\pmod h,
\]
after choosing a generator for its faithful base character.  The
equation on the characters follows from the quadratic leading term
of \(z=xy\); it is invariant under changes of local coordinates.
An element which acts nontrivially on the elliptic tangent cannot
be a pseudoreflection.  Indeed, a pseudoreflection with nontrivial
elliptic character would have to fix both normal lines pointwise,
and would therefore fix \(z\).  Faithfulness makes that element
the identity, a contradiction.  Hence every pseudoreflection lies
in the kernel of the elliptic tangent character.

We spell out the smooth-quotient criterion used here.  A finite
analytic action fixing the origin is holomorphically linearized by
the germ
\[
 F(v)=\frac1{|H|}\sum_{g\in H}(Dg_0)^{-1}g(v).
\]
It is equivariant, its derivative is the identity, and the inverse
function theorem supplies its inverse.  Averaging over the finite
group commutes with completion of the invariant local ring.  Thus
the analytic quotient is smooth exactly when the algebraic linear
quotient has a regular local ring at the origin.  Satriano,
Corollary 1.8, applies to that linear action; constant finite
groups are stable by Proposition 2.2.  It says that the stabilizer
of a smooth quotient is generated by pseudoreflections
\cite[Corollary 1.8 and Proposition 2.2]{Satriano2012}.
If the elliptic tangent character were nontrivial, the preceding
kernel argument would contradict this criterion.  It is therefore
trivial for every point stabilizer.

A finite-order automorphism of an elliptic curve which fixes a
point and has derivative one is the identity: the same averaging
linearization in one dimension proves identity near that point,
and analytic continuation proves identity on the connected curve.
Consequently the effective action of the setwise stabilizer of
\(E\) is free on \(E\).  Its quotient is again an elliptic
curve, by the unramified Riemann--Hurwitz formula.  If the root group
permutes double curves, choose one representative from any orbit;
its image is precisely the quotient by its setwise stabilizer.
It is a smooth elliptic curve fixed by the descended line circle.
Possible stabilizer elements acting identically on \(E\) can
ramify in a normal direction, but they do not change this curve
quotient or its genus.

Choose a prime not dividing the two nonzero normal circle weights
along the elliptic fixed curve.  Average a Riemannian metric and
use its equivariant normal exponential map; in that neighbourhood
the fixed set of the prime-order subgroup is exactly the zero
section.  Thus the elliptic curve is a connected component of
that finite-subgroup fixed set.  The Smith--Floyd inequality
\cite[Theorem 4.4, p.~146]{Floyd1952} bounds its total mod-prime
Betti number by that of \(S^6\), which is two.  The theorem's
locally compact, finite-dimensional Hausdorff hypotheses apply to
this compact smooth manifold; its Cech homology agrees with
singular homology here and on the smooth fixed submanifold.
An elliptic fixed component already has total Betti number four,
a contradiction.  This proves the
assertion without assuming that the finite descent is unramified
in the normal directions.
\end{proof}

\begin{statement}{P3K.13}{Proved}
In a retained nonzero-linearization \(I_n\) neighbourhood, the
circle quotient is the mapping torus of the actual source
translation.  If its global oriented orbit space has
\(H_2(Y;\mathbb Z)=0\), that translation is transitive on the
\(I_n\) components.  In particular this conclusion applies to
the retained neighbourhoods of the semifree sphere members with
the single standard fixed \(\mathbb P^1\).
\end{statement}

\begin{proof}
Let \(S_D\) be the smooth elliptic-surface
neighbourhood of the fibre, and let \(Q:S_D\to S_D\) be
translation by the specified section.  The actual nonzero
linearization gives \(M^\times/(\lambda\phi)\).  Quotient
first by the unit line circle.  A Hermitian norm identifies
\(M^\times/S^1\) with \(S_D\times\mathbb R\), where
the second coordinate is log norm.  The lifted translation becomes
\((p,t)\mapsto(Qp,t+\rho(p))\), after reversing the chosen
generator if necessary, with \(\rho(p)>0\).  This inequality
holds uniformly on a compact smaller neighbourhood by contraction.
Its fundamental region is
\(0\le t\le\rho(p)\), with the upper boundary identified to
the lower boundary by \(Q\).  Reparametrizing each interval
smoothly to \([0,1]\), with compatible collar germs at its
ends, gives an orientation-preserving diffeomorphism to
\[
 T_Q=(S_D\times[0,1])/(p,1)\sim(Qp,0).
\]
This is the actual orbit map; no period-only replacement has been
used.  The global orientation on this neighbourhood is fixed by
the complex orientation of \(S_D\), followed by the positive
suspension coordinate.  Reversing that convention would multiply
all intersection numbers below by the same sign and would not
affect their nonvanishing.

Label the fibre components cyclically by \(C_j\),
\(j\in\mathbb Z/n\), and let \(Q\) shift labels by
\(c\).  Put \(g=\gcd(n,c)\) and \(h=n/g\).  Its
component orbits have the integral cycle sums
\[
 Z_r=\sum_{j\equiv r\;({\rm mod}\;g)}C_j,
 \qquad 0\le r<g.
\]
Suppose \(g>1\).  No two components in one orbit are adjacent
in the cycle.  Hence each \(Z_r\) is a disjoint union of smooth
oriented two-spheres in the slice \(S_D\times\{0\}\).
Its translation mapping torus
\[
 D_r=(Z_r\times[0,1])/(p,1)\sim(Qp,0)
\]
is a smooth closed oriented three-manifold in \(T_Q\),
because \(Q\) permutes these disjoint spheres by holomorphic
isomorphisms.  It defines an integral three-cycle.  These cycles
are compact and lie in the interior of the local
neighbourhood, so inclusion defines classes in the global
\(H_2(Y;\mathbb Z)\) and \(H_3(Y;\mathbb Z)\).

The local intersection pairing is
\[
                    [Z_r]\cdot[D_s]=Z_r\cdot_S Z_s.
\]
To verify it, represent the two-cycle by a small transverse
perturbation in a source slice and represent the mapping-torus
cycle in a product collar of that slice.  Its extra tangent
direction is the positive suspension direction.  The local
five-dimensional intersection sign is therefore exactly the
complex-surface intersection sign.  The same calculation using
a normal pushoff computes the diagonal self-intersections.
All perturbations are compactly supported in this neighbourhood,
so the calculation also equals the global homological pairing.

Let \(B_n\) be the intersection matrix of the \(I_n\)
cycle: diagonal entries are \(-2\), successive distinct
components meet once, and the two \(I_2\) components meet
twice.  Put \(B_1=(0)\).  Summing by the orbits gives
\[
                 (Z_r\cdot Z_s)_{r,s}=h B_g.
\]
This identity follows by counting the \(h\) equal edges
between adjacent orbits.  The quadratic form is
\[
 (t_0,\ldots,t_{g-1})B_g(t_0,\ldots,t_{g-1})^t
       =-\sum_{r\in\mathbb Z/g}(t_r-t_{r+1})^2.
\]
Thus its rank is \(g-1\), and it has a nonzero entry whenever
\(g>1\).  Such an entry is impossible when
\(H_2(Y;\mathbb Z)=0\), because every displayed global
two-cycle would then have zero pairing with every closed
three-cycle.  Hence \(g=1\), proving transitivity.
The program \path{code/part3_fourfold/vertical_scope.py} exports
every orbit matrix for \(2\le n\le12\), all shifts, and
their exact Smith forms.  The displayed cycle and intersection
arguments prove the general formula, independently of that
finite check.
\end{proof}

\begin{statement}{P3K.14}{Proved}
Let \(S\) be one of the rational elliptic surfaces with
\({\rm MW}(S)=\mathbb ZG\).  Write the generic Picard point
of \(M\) as \(aG\), and the translating section as \(bG\).
For a literal Remark 8.5 line with
\(M\otimes t_{bG}^*M^{-1}=\pi^*\mathcal O(k)\),
\[
                         k=ab\,h(G).
\]
At a retained \(I_n\) on which the translation is transitive,
the whole-fibre linearization forces the generic Picard section
to meet the identity component.  At that fibre every vertical
correction in \(M\) is a whole-fibre divisor.  In particular,
for the transitive component markings in the rank-one table,
\(n\mid a\).
\end{statement}

\begin{proof}
On the generic fibre,
\(M\otimes\mathcal O(aG-O)^{-1}\) is trivial.  A nonzero
rational trivializing section has divisor supported on fibres,
so
\(M=\mathcal O(aG-O+D)\) for an integral vertical divisor
\(D\).  Choose a positive integer \(N\) such that
translation by \(NbG\) fixes every irreducible fibre
component.  Iterating the actual line isomorphism telescopes to
\[
 M\otimes t_{NbG}^*M^{-1}=\pi^*\mathcal O(Nk).
\]
The pullback of \(D\) by this translation equals \(D\)
as a divisor, so it cancels from the left-hand side.  Intersecting
with the zero section computes its degree as
\[
 (aG)\cdot O+1-(aG-NbG)\cdot O+(NbG)\cdot O.
\]
In the section-height formula of P3F.2, the local corrections
for \(aG\) and \(aG-NbG\) agree, since \(NbG\) meets
the identity component everywhere; those for \(NbG\) are
zero.  Substitution cancels these local terms and the constant
terms, leaving
\[
 \frac{a^2-(a-Nb)^2+(Nb)^2}{2}\,h(G)
                         =Nab\,h(G).
\]
Division by \(N\) proves the asserted degree, with no
restriction on the allowed vertical divisor.

For the local statement, intersect the whole-fibre line relation
with every \(C_j\).  It gives
\(M\cdot C_j=M\cdot t_{bG}(C_j)\).  Under a transitive
permutation these degrees are equal.  Their sum is the degree
of \(M\) on the fibre, namely zero, so every degree is zero.
If \(aG\) meets \(C_r\), the component-degree equation is
\[
                   e_r-e_0+B_n d=0,
\]
where \(d\) is the vector of coefficients of \(D\) at this
fibre.  For every integral vector \(d\),
\(\sum_j j(B_nd)_j\) is divisible by \(n\), as follows
by subtracting consecutive cyclic differences.  Applying this
functional to the displayed equation gives \(r\equiv0\pmod n\).
Thus \(aG\) meets the identity component.  The equation
now reduces to \(B_nd=0\).  The displayed quadratic form
in P3K.13 shows that \(d\) is constant, so this divisor is
an integer multiple of the full fibre.  The actual generator
markings derived in P3F.2 are: components \(3,5\) at \(I_8\);
\(2,5\) at \(I_7\); \(2,3\) at the OS55 \(I_5\)
and \(1,4\) at the OS56 \(I_5\); \(1,3\) at \(I_4\);
\(1,2\) at \(I_3\); and \(1\) at \(I_2\).
P3F.2 obtains these labels from the enumerated rank-one source
lattices and the local height contributions, rather than choosing
an arbitrary component marking.  Every displayed label is a unit
in its \(\mathbb Z/n\).  Thus \(r=a[G]=0\) is exactly
\(n\mid a\).  This also proves that the local removal of
the vertical divisor is multiplication by the actual meromorphic
base factor \(t^j\), a holomorphic invertible map on the
punctured annulus.  Its exponent is the integral delta collar
change; no fractional or unrecorded gauge is introduced.
\end{proof}

\subsection{Remaining scope after these reductions}

P3K.15 below transports the component cycles, the line, and its integer
delta character through a supplied affine cap on an actual extending
root line.  Applying it to a proposed new cap still requires constructing
that root source, its line, and the complete cyclic affine norm.  Partial
roots may leave a reduced reducible fibre, and zero-gamma finite kernels
must be retained.  Thus neither the retained-neighbourhood obstruction
nor the cap comparison alone is a complete enumeration of the literal
arbitrary-site scope.

\begin{statement}{P3K.15}{Proved}
Every smooth sphere member with the standard fixed sphere and
normal circle weights \((-1,+1)\) has a semifree line-circle
action.  For a cyclic affine log cap at a nonzero-linearization
site, written on a smooth root source as a linearized line quotient, its actual circle orbit is
a mapping torus of a smooth source-surface quotient.  The finite
zero-gamma subgroup and its line character are retained in this
comparison.  If the global orbit has \(H_2(Y;\mathbb Z)=0\),
the resulting monodromy is transitive on the components of its
central fibre.  The descended line has multidegree zero there,
and a vertical change with the same generic Picard point is an
integer multiple of the primitive whole-fibre divisor.
\end{statement}

\begin{proof}
\par\medskip\noindent\textbf{Semifreeness from the actual fixed sphere.}\enspace Let
\(F\cong S^2\) be the fixed sphere whose two normal weights
are \(-1,+1\).  For every prime \(p\), an invariant normal
exponential map identifies the \(\mu_p\)-fixed set in a
neighbourhood of \(F\) with \(F\) itself.  Thus \(F\)
is an entire connected component of \(X^{\mu_p}\).
It contributes two to the total mod-\(p\) Betti number.
Floyd's inequality \cite[Theorem 4.4]{Floyd1952} bounds the
total by two.  Any point of \(X^{\mu_p}\setminus F\)
would belong to another component and contribute at least one,
which is impossible.  Every nontrivial finite subgroup of the
circle contains a subgroup of prime order.  Consequently no
finite nontrivial stabilizer occurs outside \(F\), proving
semifreeness without an assumption on the distant log caps.

\par\medskip\noindent\textbf{The lifted group and its finite kernel.}\enspace Write the
actual root source as \(S'\to\Delta_z\), a smooth surface,
and its generic-degree-zero line as \(N\).  Its scalar
suspension lift is a holomorphic line automorphism \(B\) of
\(N^\times\), covering the source translation \(Q\).
For sufficiently small nonzero scalar it is uniformly contracting
on a compact smaller source neighbourhood.  The actual order-\(m\)
affine cap action on \(N^\times/\langle B\rangle\) lifts
to a holomorphic line map \(A\) on the root source, with
\[
 AB=BA,\qquad A^m=B^n,
\]
and with root action \(z\mapsto\zeta_m z\).  Here \(n\)
is the integral gamma coefficient of the actual cyclic norm.
More explicitly, on the vector cover write the affine action
as \((p,R)\).  Its \(m\)-th power translates by
\[
            p+Rp+\cdots+R^{m-1}p
                 =n\gamma+\lambda_0,
\]
where \(\lambda_0\) is an elliptic-and-delta period.  This
is the actual order-\(m\) descent condition.  The latter
period acts trivially on \(N^\times\), whereas the gamma
period is exactly \(B\); hence \(A^m=B^n\).
For the supplied caps with invariant vector \(v\),
\(p=\Pi v/m\) and the norm is exactly \(\Pi v\), so
\(n=\psi(v)\).  For a noninvariant affine translation the
displayed norm, rather than its initial gamma coefficient,
is required.  The root deck preserves this gamma point in
the semiabelian group, so its affine lift commutes with \(B\).
This retains any delta phase of \(A\); it is not discarded
from the line lift.

First take \(\gcd(m,n)=1\).  Choose integers \(r,s\)
with \(rn+sm=1\), and set \(H=A^rB^s\).  The displayed
relations give
\[
                  A=H^n,\qquad B=H^m.
\]
These are identities of actual line maps, not merely identities
of their homological actions.  Thus the original cap quotient
\(N^\times/\langle A,B\rangle\) is exactly
\(N^\times/\langle H\rangle\).  In a root choice for the
scalar, \(A\) contains its \(n/m\)-th power, while \(B\)
contains its first power.  Hence \(H\) contains the positive
power \(1/m\), times a line map bounded on the chosen compact
source neighbourhood.  Taking the scalar smaller makes \(H\)
uniformly contracting.  This conclusion also holds when some of
\(n,r,s\) are negative; the exponent is the exact Bezout
sum \(rn/m+s=1/m\).

For the nonprimitive case put \(g=\gcd(m,n)\),
\(m'=m/g\), \(n'=n/g\), and define
\[
                  C=A^{m'}B^{-n'}.
\]
Then \(C^g=1\), and the root-base character shows that its
order is exactly \(g\).  It is the complete finite subgroup
of the abelian group generated by \(A,B\).  Choose
\(rn'+sm'=1\).  Modulo \(C\), the map
\(H=A^rB^s\) satisfies
\[
             \bar A=\bar H^{n'},\qquad
             \bar B=\bar H^{m'}.
\]
This describes the actual finite quotient followed by the
residual contracting suspension.  In particular it does not
replace \((m,n)\) by a reduced pair while forgetting its
finite kernel.

The marked outer meridians are retained under this change of
generators.  In the ordered generators \((C,H)\), they are
exactly
\[
                   A=C^sH^{n'},\qquad
                   B=C^{-r}H^{m'}.
\]
The matrix with columns \((m',-n')^t\) and \((r,s)^t\)
has determinant one.  It transforms the sole relation
\((m,-n)^t\) to \((g,0)^t\).  Thus even for an order-one
cap the integer annular clutch is still the specified meridian;
it has not been lost in the local mapping-torus description.
\path{code/part3_fourfold/finite_cap_group.py} exports these matrices
and checks them for all \(1\le m\le24\),
\(-2m\le n\le2m\), as well as retaining the general
Bezout proof above.

Let \(\bar S=S'/\langle C\rangle\).  A stabilizer of a
point \(p\in S'\) acts on the complex line \(N_p\) by
a root of unity.  If that root were nontrivial, the corresponding
unit line rotation would identify two points of the cap quotient
and produce nontrivial finite line-circle isotropy.  Semifreeness
therefore forces every source-point stabilizer to act trivially
on its line fibre.  The line descends to a line bundle
\(\bar N\) on \(\bar S\).  This descent can be seen
without a categorical shortcut: choose a local frame nonzero at
the point and average it over its stabilizer.  Its value remains
nonzero because the stabilizer acts trivially on that fibre.
The invariant frame identifies the quotient line with a product
over the local source quotient.  Translates of these frames give
the transition functions on the quotient.

The finite source quotient is smooth.  Indeed the line-complement
quotient \(N^\times/\langle C\rangle\) is smooth, since
its further quotient by the residual contracting infinite cyclic
group is the given smooth cap and that quotient is locally a
covering.  In the invariant frames just constructed it is
\(\bar S\times\mathbb C^*\).  A local analytic product
with \(\mathbb C^*\) is smooth exactly when its base is
smooth, as follows from the local dimension and cotangent-space
dimension.  This proves smoothness of \(\bar S\).
The source quotient may be branched in a normal direction; the
argument only uses the proved triviality of the line stabilizers.

\par\medskip\noindent\textbf{The real orbit map and its line.}\enspace On \(\bar N\)
choose a Hermitian metric and write log norm as \(t\).
The contracting residual map becomes
\[
(p,t)\longmapsto(hp,t-\rho(p))
\]
on the line-circle quotient, with
\(\rho(p)=-\log(\|H(v)\|/\|v\|)>0\).  It is
bounded away from zero on a compact smaller neighbourhood.
The upper face \((p,\rho(p))\) is identified by this actual
map with \((hp,0)\).  As in P3K.13, this interval fundamental
region is smoothly
reparametrized to the unit interval with compatible end collars.
This gives an explicit orientation-preserving diffeomorphism
from the actual circle orbit of the cap to
\[
           T_h=(\bar S\times[0,1])/(p,1)\sim(hp,0).
\]
The principal circle line is also identified: it is the unit
circle bundle of \(\bar N\) on this mapping torus, glued at
the ends by the unit normalization of the actual line map
\(\bar H\).  Both its scalar phase and any finite-kernel
character have therefore been transported, not reconstructed
from the orbit alone.

\par\medskip\noindent\textbf{The component intersection obstruction through the cap.}\enspace
Let the connected central fibre of \(\bar S\) be
\(\sum_i m_i C_i\), and let \(B=(C_i\cdot C_j)\).
All \(m_i\) are positive.  The relation \(B m=0\) and
nonnegative off-diagonal intersections give the exact identity
\[
 x^tBx=-\sum_{i<j}(C_i\cdot C_j)m_i m_j
       \left(\frac{x_i}{m_i}-\frac{x_j}{m_j}\right)^2.
\]
Connectedness of the fibre implies that the kernel is exactly
the line spanned by \(m\).  The map \(h\) permutes its
components and preserves the multiplicities.  If there are
\(r\) component orbits, the invariant coefficient space has
dimension \(r\).  The restriction of \(B\) to it has
rank \(r-1\).

We realize this restricted form by actual cycles in \(T_h\),
including the cases where components within one orbit meet.
Let \(z\) be an integral two-cycle representing an orbit-sum
divisor.  Its class is \(h\)-invariant, so choose a compact
singular three-chain \(c\) in \(\bar S\) with
\(\partial c=h_*z-z\).  The product trace of \(z\)
over the suspension interval, corrected by \(-c\) in the
seam slice, is a closed integral three-cycle \(D_z\) in the
mapping torus.  A test two-cycle in a slice halfway along the
interval misses the correcting chain.  Its intersection with
\(D_z\) is exactly its surface intersection with \(z\),
with the complex-surface orientation followed by the positive
suspension direction.  Thus all entries of the invariant
intersection form are realized by compact two- and three-cycles
in the actual cap orbit.  Inclusion in the global orbit preserves
these intersection numbers.  If \(H_2(Y;\mathbb Z)=0\),
every such number is zero, forcing \(r=1\).

Finally \(\bar H\) is a line automorphism, so the degrees
\(\bar N\cdot C_i\) are invariant under the component
permutation.  Transitivity makes them all equal.  Their positive
multiplicity-weighted sum is the generic fibre degree, zero;
hence every one is zero.  For two such line extensions with the
same generic Picard point, let \(L\) be their ratio line.
The proper source map to the disc has coherent direct image
\(\pi_*L\) by proper direct-image coherence.  It has rank
one because the generic restriction is the trivial line on a
connected elliptic curve.  It is torsion-free: multiplication
by a nonzero base function is injective on the line sheaf on
the smooth source, hence also on its direct image.  On the
disc this coherent sheaf is therefore free of rank one after
shrinking.  A nonzero local generator pulls back to a section
of \(L\), nonzero on generic fibres, whose divisor \(D\)
is vertical.  This gives the required meromorphic line
comparison from the proper source itself, without extending
an arbitrary punctured-disc scalar.  The component-degree
equations now say \(BD=0\).
Thus \(D\) is an integer multiple of the primitive integral
fibre vector \(F_{\rm prim}\).  If the actual fibre is
\(\nu F_{\rm prim}\), this need not yet be a base pullback.
On the normalization of the root which removes that multiplicity,
\(\mathcal O(kF_{\rm prim})\) is represented by the
base monomial \(z^k\).  Its deck character is
\(\zeta_\nu^k\), so the comparison records exactly the
integer delta character of the log cap.  This proves the stated
vertical rigidity with its full remaining collar ambiguity.
\end{proof}

The statement concerns caps supplied with their actual extending
linearized root line and affine action at a nonzero-linearization
site.  Scalar-zero Mumford fillings are treated separately.  For any proposed broader
cap, constructing these maps remains necessary before applying
the statement.  The OS50 quadratic-root lead is kept separately
until its source marking and full integer topology are certified.

\begin{statement}{P3K.16}{Proved}
For the whole-fibre \(I_1\) cusp construction of P3F.9, together
with the actual nonzero-linearization affine caps of P3K.15, the
delta orbit has a global map to a closed Seifert three-manifold
over the original base \(\mathbb P^1\).  Its regular fibres are
connected elliptic tori.  Every sphere member has at most two
nontrivial effective gamma multiplicities.  For two such
multiplicities, write the marked local pairs as
\((m_1,n_1),(m_2,n_2)\), with each pair primitive, and let
\(b\) be the remaining integral gamma clutch.  The exact restriction is
\[
       m_1m_2b+m_2n_1+m_1n_2=\pm1.                 \tag{P3K.16.1}
\]
In particular \(\gcd(m_1,m_2)=1\).  Order-one clutches and the
finite zero-gamma subgroups are retained in the source and line
data; this restriction does not remove them.
\end{statement}

\begin{proof}
\par\medskip\noindent\textbf{Construction of the map.}\enspace On an ordinary source disc,
the delta orbit of the contracting line quotient is the source
mapping torus of P3K.13.  Projection of that source surface to
its disc, together with the suspension coordinate, gives the
map to the disc times the gamma circle.  Straightening the norm
roof is over the source and its base.  On the whole-fibre \(I_1\)
cusp, the actual product of P3F.9 preserves the Seifert gamma
coordinate, so its map is exactly
\[
       E_{I_1}\times S^1_\gamma\longrightarrow
                    \Delta_q\times S^1_\gamma.
\]
This includes the nodal fibre and the circle-fixed sphere.
The specified annular period maps intertwine these projections;
their integral gamma shifts are their recorded Seifert clutches.

Here is the map at an arbitrary supplied cap, including its
finite kernel.  Use the actual maps \(A,B,C,H\) of P3K.15,
and write \(g=\gcd(m,n)\), \(m'=m/g\), \(n'=n/g\).
The root source has coordinate \(z\), with original base
coordinate \(q=z^m\).  The finite map \(C\) has faithful
base character of order \(g\).  Hence the smooth source
quotient \(\bar S=S'/\langle C\rangle\) has base coordinate
\(u=z^g\), and \(q=u^{m'}\).  If
\(rn'+sm'=1\), then \(\bar H\) acts on this base by
\[
                 u\longmapsto e^{2\pi i r/m'}u.
\]
Projecting the actual mapping torus of \(\bar S\) to the
mapping torus of this disc rotation gives the cap map.
The target is a solid torus: the explicit map
\((u,t)\mapsto(e^{2\pi i rt/m'}u,e^{2\pi it})\)
trivializes its disc bundle.  The regular Seifert gamma orbit
has multiplicity \(m'\), and its primitive marked meridian
relation is
\[
                     m'a-n'\gamma=0.             \tag{P3K.16.2}
\]
Indeed these two loops are the actual \(\bar A=\bar H^{n'}\)
and \(\bar B=\bar H^{m'}\) of P3K.15, not an inferred
replacement for their affine lifts.  This transports the stated
outer meridians and glues the target solid torus to the ordinary
target charts.  Doing so at all sites constructs the closed
three-manifold \(B\) and the map \(p:Y\to B\).

Off the finitely many discriminant and cap-core circles, the
map is a smooth proper submersion with connected elliptic
fibres.  Connectedness at a cap can be checked before any
topological conclusion: the \(g\) roots of a nonzero \(u\)
are permuted transitively by \(C\), so their quotient fibre
is one elliptic curve.  The same fact is preserved by the
suspension.  The finite-root and ordinary maps above agree
with the annular source maps, so there is no disconnected
generic-fibre substitution in this construction.

\par\medskip\noindent\textbf{Fundamental groups and the number of sites.}\enspace Every
loop in the smooth three-manifold \(B\) can be moved off the
finite union of discriminant and core circles: perturb its
finitely many smooth arcs transversely to these curves, whose
dimension sum is two, less than three.  Over the resulting
loop, a choice of horizontal complement to the submersion
kernel lifts the loop from a chosen point of its fibre.  The
lift exists for the whole compact interval because the
submersion is proper over that interval.  Its endpoint lies
in the same connected elliptic fibre as its starting point.
Join them by a path in that fibre.  The resulting closed loop
projects to the original loop followed by a constant path.
Thus \(p_*:\pi_1(Y)\to\pi_1(B)\) is surjective.
For a sphere member, P3F.8 gives \(Y=S^5\), so
\(\pi_1(B)=1\).

Remove discs around the nontrivial effective Seifert sites.
The gamma-circle bundle over the remaining punctured sphere
is trivial, since its base retracts onto a graph.  Choose a
section, retain the marked boundary loops \(a_i\), and write
its integer discrepancy from the actual cap framings as \(b\).
Van Kampen applied to the actual meridians (P3K.16.2) gives
\[
 \pi_1(B)=\left\langle a_1,\ldots,a_s,\gamma\ \middle|\
 [a_i,\gamma]=1,\quad a_i^{m_i}=\gamma^{n_i},\quad
                 a_1\cdots a_s\gamma^b=1\right\rangle.
                                                        \tag{P3K.16.3}
\]
Order-one sites can be included in this presentation and
eliminated using their actual relation \(a_i=\gamma^{n_i}\).
Their integers are then added to \(b\); their source and
line lifts are still retained in P3K.15.

Quotient (P3K.16.3) by \(\gamma\).  With at least three
nontrivial multiplicities, further killing all but three
generators maps this quotient onto
\[
       \langle x,y,z\mid x^{m_1}=y^{m_2}=z^{m_3}=xyz=1\rangle.
\]
This group is nontrivial for every \(m_i\ge2\), as the
following explicit real matrix representation proves.  On
\(\mathbb R^3\), take the symmetric bilinear form with matrix
\[
 G=\begin{pmatrix}
 1&-\cos(\pi/m_3)&-\cos(\pi/m_2)\\
 -\cos(\pi/m_3)&1&-\cos(\pi/m_1)\\
 -\cos(\pi/m_2)&-\cos(\pi/m_1)&1
 \end{pmatrix},
\]
and set \(R_i(v)=v-2(v,e_i)_G e_i\).  Direct substitution
gives \(R_i^2=I\).  The span of each pair \(e_i,e_j\)
has positive definite Gram matrix, since its determinant is
\(\sin^2(\pi/m_k)>0\).  It therefore has a one-dimensional
orthogonal complement even when \(G\) itself is degenerate.
On the pair span the two reflections have angle \(\pi/m_k\),
and their product has characteristic polynomial
\(z^2-2\cos(2\pi/m_k)z+1\); on the complement it is
the identity.  Thus \(R_2R_3,R_3R_1,R_1R_2\) have
orders exactly \(m_1,m_2,m_3\), and their product is
exactly the identity by cancellation of adjacent reflections.
Their first member is nonidentity.  This gives the promised
nontrivial representation, contradicting
\(\pi_1(B)=1\).  Thus \(s\le2\).

\par\medskip\noindent\textbf{The exact determinant and the marked base.}\enspace For two
sites eliminate \(a_2\) from (P3K.16.3).  The remaining
generators \(a_1,\gamma\) commute, and their relation matrix is
\[
 \begin{pmatrix}
      m_1&-n_1\\
      -m_2&-m_2b-n_2
 \end{pmatrix}.
                                                        \tag{P3K.16.4}
\]
Triviality of the group is therefore exactly (P3K.16.1).
Any common divisor of \(m_1,m_2\) divides its left-hand
side, which proves their coprimality.  With one exceptional
site the same calculation is \(|m_1b+n_1|=1\); with none
it is \(|b|=1\).  Equivalently one may insert multiplicity-one
solid tori and use the displayed two-site matrix in all cases.

The resulting target is the actual marked Hopf three-sphere,
with both signs retained.  Here is an explicit Heegaard
comparison, rather than an appeal to a classification by
fundamental group.  In the middle torus with basis
\((a_1,\gamma)\), its two meridians are
\[
 \mu_1=(m_1,-n_1),\qquad
 \mu_2=(-m_2,-m_2b-n_2).
\]
Put \(\delta=\det(\mu_1,\mu_2)=\pm1\).
The integer matrix with columns \(\mu_1,\delta\mu_2\)
has determinant one and sends these meridians to the two
coordinate meridians after inversion.  This torus map
extends across each solid torus, because in meridian-longitude
coordinates every required boundary shear extends by
\[
 (\rho e^{i\theta},e^{i\varphi})\longmapsto
 (\rho e^{i(\varepsilon\theta+k\varphi)},
                         e^{i\varepsilon'\varphi}),
 \qquad \varepsilon,\varepsilon'\in\{\pm1\}.
\]
The two target solid tori are the standard decomposition of
\(S^3\subset\mathbb C^2\).  In these coordinates the
original gamma vector is exactly \((\delta m_2,m_1)\).
Consequently its signed weighted Hopf action and its two
exceptional orders have both been transported explicitly.

Finally the replacement \(a_i\mapsto a_i\gamma^{k_i}\)
changes the integer records by
\[
         n_i\mapsto n_i+m_i k_i,\qquad
         b\mapsto b-\sum_i k_i.
\]
This is also an identity of actual cap quotients: replace
\(A_i\) by \(A_iB_i^{k_i}\), which generates the same
group with \(B_i\).  Its cyclic norm changes by
\(m_i k_i\gamma\); its finite kernel \(C_i\) is unchanged.
Thus residue representatives for \(n_i\) can be chosen
without deleting the recorded finite kernel, delta character,
or source-line map.  The determinant is invariant under
exactly these transformations.
\end{proof}

The scope of this reduction is the actual whole-fibre cusp
projection specified at the start.  It does not extend the old
gamma coordinate through the component-supported OS49 cusp;
that construction uses its separate mixed \(K\) marking.
The Seifert projection also does not recover the finite-kernel
source quotient or the line character from \(B\) alone.
Those data remain part of every cap in the enumeration.

An arbitrary periodic unimodular fan need not have this fixed-sphere
model.  The present map and its use of P3F.8 are not asserted for
such a fan without its separate orbit comparison.

\begin{statement}{P3K.17}{Unresolved; scalar-fan audit correction}
\par\noindent\textbf{Audit correction, 17 November 2026.} The OS49 extra fixed-normal boundary must retain its own outer collar arrows. The whole-I1 obstruction does not compute that separate correction. The claimed complete two-fan equality is reopened pending that explicit outer-marking audit. The formulas and earlier argument below are retained as the audited prior claim, not as current proved conclusions. See P3K.61--62 and the correction ledger.
\medskip

The negative- and positive-diagonal OS49 families have equal
actual return index for every cap order \(m=3k\),
\(k>0\), every signed primitive twist \(n=\pm1\),
and every additional integral delta twist \(d\delta\).
Their index is zero, their oriented return class is zero, and
their actual mapping tori are products as in P3F.24.
The comparison is an oriented bordism carrying the actual
Euler line; it imposes no spin hypothesis on the comparison
bordism.  The endpoint index structures are the prescribed
ones from the original families.
\end{statement}

\begin{proof}[Preserved earlier argument; scope corrected above]
\par\medskip\noindent\textbf{The real master space and its actual resolution maps.}\enspace

Order the Cox coordinates as \((z_0,z_4,z_1,z_2)\), corresponding
to rays \(u_0,u_4,u_1,u_2\).  Let an auxiliary circle act with
weights \((1,1,-1,-1)\), and let the physical delta circle have
weights \((-1,0,0,+1)\).  These actions commute.  Write
\(a=(z_0,z_4)\), \(c=(z_1,z_2)\), and
\[
       \mu=\|a\|^2-\|c\|^2.
\]
The auxiliary circle is free on \(\mathbb C^4\setminus\{0\}\).
Consequently its quotient is a smooth oriented real seven-manifold.
The residual delta action is semifree.  Indeed every nonzero
two-by-two minor of the matrix with rows
\((1,1,-1,-1)\) and \((-1,0,0,1)\) is \(\pm1\).
On a coordinate support of rank two this proves that the
two-torus has trivial stabilizer.  A rank-one support lies in
one of the two pairs \((z_4,z_1)\), \((z_0,z_2)\),
and the residual circle fixes it pointwise.  Subtracting the
auxiliary lift on these two supports gives normal weights
\((-1,+1)\) in each case.  Thus the circle orbit has the
ordinary rank-three Hopf normal model along its fixed locus,
with no additional finite-quotient strata.
Its positive moment levels are the negative-diagonal small resolution,
and its negative moment levels are the positive-diagonal small
resolution.  Here are the actual maps proving these identifications.
On the positive level use the complex quotient
\[
   \{a\ne0\}\big/\bigl((a,c)\sim(ta,t^{-1}c)\bigr)
       \longrightarrow
   \{([a],ac^t):\operatorname{im}(ac^t)\subset\mathbb Ca\}.
\]
Its target is the total space of \(\mathcal O(-1)^{\oplus2}\)
on \(\mathbb P^1_a\).  Every complex orbit meets a prescribed
positive moment level in exactly one auxiliary-circle orbit:
the required positive real factor solves
\[
 r^2\|a\|^2-r^{-2}\|c\|^2=\tau,
 \qquad
 r^2=\frac{\tau+\sqrt{\tau^2+4\|a\|^2\|c\|^2}}
                 {2\|a\|^2}.
\]
This is smooth for \(a\ne0\) and \(\tau>0\), including
\(c=0\).  The negative-level map is the same formula with
the two pairs interchanged.  Both maps commute with delta.
The positive stability condition excludes \(z_0=z_4=0\),
so its two fan triangles have common edge \(u_1u_2\);
the negative condition has common edge \(u_0u_4\).
Thus these are the specified two fans with their original
coordinate identifications on the common rank-one matrix locus.

Restrict to \(|\mu|\le\epsilon\) and a large fixed radius,
and remove a ball of radius \(\rho\) with \(\rho^2<\epsilon\).
The two moment faces are the truncated small resolutions, the
outer face is their common exterior times an interval, and the
new inner face is
\[
             S^7/S^1_{\rm aux}=\mathbb {CP}^3.
\]
This is an actual real bordism relative to the common exterior.
Conjugating the two auxiliary-negative coordinates makes the
auxiliary action scalar-positive; two complex conjugations
preserve the real orientation.  In the resulting homogeneous
coordinates, put
\[
       A=(z_0,\overline z_2),\qquad
       B=(z_4,\overline z_1).
\]
Delta acts by \((A,B)\mapsto(e^{-it}A,B)\).

\par\medskip\noindent\textbf{The circle orbit, its line, and the explicit joint cap.}\enspace

The orbit map on this inner face is
\[
 [A,B]\longmapsto
 \left([A],[B],\frac{\|A\|^2}{\|A\|^2+\|B\|^2}\right),
\]
with the irrelevant factor collapsed at either endpoint.
Thus the actual orbit is \(\mathbb {CP}^1*\mathbb {CP}^1\),
identified with \(S^5\subset\mathbb R^3\oplus\mathbb R^3\)
by \((p,q,\phi)\mapsto(\cos\phi\,p,\sin\phi\,q)\).
The two fixed spheres are its two join ends.  On the free orbit
\(S^2_A\times S^2_B\times(0,1)\), the circle line is exactly
\[
           L=\mathcal O(1)_A\boxtimes\mathcal O(-1)_B,
           \qquad x=h_A-h_B.                         \tag{P3K.17.1}
\]
Indeed a nonzero pair determines the relative tensor of its
two tautological lines; delta acts with negative weight on
that tensor.  Taking the associated line with the physical
positive delta character gives the displayed inverse tensor.
This fixes the sign rather than choosing it from the two
normal models independently.

The two base charges are therefore \(+1,-1\).  In particular
the line does not extend over either individual conifold cap:
the restriction image of
\(H^2(D^3_A\times S^2_B)\) in the first boundary has no
\(h_A\) term, and the other cap has no \(h_B\) term.
We retain this obstruction and instead join the fixed spheres.

Take \(C=S^2_P\times I\), with meridian sphere \(S^2_Q\),
and the line \(\mathcal O(1)_P\boxtimes\mathcal O(1)_Q\)
on \(C\times S^2_Q\).  At its first end use
\[
              (A,B)=(P,\iota Q),
\]
and at its second end use
\[
              (A,B)=(Q,\iota P),                    \tag{P3K.17.2}
\]
where \(\iota\) is the antipodal map of \(\mathbb {CP}^1\).
Both maps pull (P3K.17.1) to \(p+q\).  These are actual
line isomorphisms: the orthogonal-complement map on lines in
\(\mathbb C^2\), followed by its determinant pairing,
identifies \(\iota^*\mathcal O(1)\) complex-linearly with
\(\mathcal O(-1)\); the second end then uses the canonical
interchange of the two tensor factors.  All maps preserve
the physical principal-circle character.

The resulting closed five-dimensional orbit correction is the
mapping torus
\[
 Q_{\rm corr}=T_g(S^2\times S^2),\qquad
           g(A,B)=(\iota B,\iota A).
\]
The map \(g\) preserves \(h_A-h_B\) and the product
orientation.  Its line is glued by the specified tensor
isomorphisms, so this statement retains the actual line,
not merely an invariant cohomology class.

\par\medskip\noindent\textbf{The actual parameter gauge on the master space.}\enspace

For the closed OS49 parameter \(\lambda=r^m\), with
\(m=3k\) and signed primitive twist \(n=\pm1\), the
extra twist \(d\delta/m\) contributes only a parameter-independent
scalar to the finite-cap line map.  Its source coordinate is unchanged;
the same scalar extends in the Cox lift and across the joint cap.
It therefore introduces no parameter-crossing line class.  The
period-origin calculation in P3F.32 gives the primitive
\(K\)-turn and the angular line winding \(2mn\).
The \(K\)-period on the ordinary Tate coordinates is
\[
      Z\longmapsto a(q)Z,\qquad
      V\longmapsto q^n r^{2mn}b(q)V,
\]
after incorporating the signed traversal (so the bounded units
include the corresponding signed powers), where those units
have parameter-independent logarithms.  Reindexing
the common fan by its integral vertical period removes the
displayed \(q^n\) by the signed integral vertical translation
\(n\).  On the resulting common Cox charts the
bounded map has the explicit lift
\[
 \begin{aligned}
 z_0&\longmapsto\bigl(a(q)b(q)r^{2mn}\bigr)^{-1}z_0,
       &z_1&\longmapsto a(q)z_1,\\
 z_2&\longmapsto b(q)r^{2mn}z_2,
       &z_4&\longmapsto z_4.
 \end{aligned}                                      \tag{P3K.17.3}
\]
The product of its four multipliers is one, so it fixes the
base monomial \(q=z_0z_1z_2z_4\).  Its torus vector is
\(\log a\,(u_1-u_0)+\log(br^{2mn})\,(u_2-u_0)\),
which is exactly the marked \((Z,V)\) action.  Thus its
parameter-dependent factor is one global physical delta
action; it is not two separately chosen gauges at the two
ends of the exceptional geometry.

For the exact quotient, label the unit-square flop cells by
\(v\in\mathbb Z^2\).  Their period lattice is
\(\Gamma=\langle(2,1),(-3,-1)\rangle=\mathbb Z^2\),
and \(3(2,1)+2(-3,-1)=(0,1)=K_{\rm trop}\).
It acts freely and transitively on these cells.  Their diagonal
curves are disjoint, since each has the two maximal cones
belonging to its own square.  Choose disjoint equivariant
tubular neighbourhoods and insert one master chart in each.
The quotient of their disjoint union by \(\Gamma\) is
exactly one representative chart with trivial stabilizer.
There is no residual \(K\)-circle on this correction.
The signed parameter return shifts its cell label by
\(n(0,1)=n(3h+2g)\).  Returning to the chosen representative
by that unique period gives (P3K.17.3).  Reversing the
representative convention inverts the one common gauge at
both ends and leaves the comparison unchanged.

Fixing \(|r|\), its angular part is
\(e^{2mn i\theta}\) in that physical circle.  The remaining
unit multipliers are independent of \(\theta\); their
chosen logarithms give equivariant isotopies on the Cox
charts.  The positive real rescalings in Section 1 transport
these isotopies to the moment levels and to the inner face.
Explicitly, follow a Cox map by its unique positive auxiliary
rescaling back to the original value of \(\mu\).  Its
implicit derivative is
\(2r^2\|a\|^2+2r^{-2}\|c\|^2>0\) away from the
removed origin, so this is smooth through the real bordism.
The bounded units can be changed to their values at \(q=0\)
using \(a(tq),b(tq)\), \(0\le t\le1\), with an
exterior cutoff.  The maps stay invertible and commute with
both circles; they are independent of the parameter angle.
This change introduces no angular character on the inner face.
These maps need not preserve a fixed Euclidean radius.
Use transported truncations: under the moment-preserving
equivariant isotopy \(G_t\), take the inner and outer
faces to be \(G_t\) of the original faces.  They stay
embedded and disjoint, while the moment faces stay at
\(\mu=\pm\epsilon\).  Glue these transported charts
and collars by their exact maps.  For the remaining
constant-unit clutch \(G_0\), its chosen logarithms give
\(G_0^t\), and
\((x,t)\mapsto(G_0^t x,e^{2\pi it})\) trivializes
the mapping torus.  The angular delta factor preserves
the round face and is trivial on its circle orbit.
Consequently the two maps (P3K.17.2) have the same angular
gauge, which extends by that constant circle multiplication
over all of \(C\times S^2_Q\).  No difference of two
angular gauges is introduced on its ends.  The closed
correction over the parameter circle therefore carries
a line pulled back from \(Q_{\rm corr}\), after this
explicit gauge identification.  Its first Pontryagin class
is pulled back from the same five-dimensional factor.
The potential mixed line class can also be checked directly.
Cover the correction by its free-orbit cylinder and joint-cap
cylinder; their overlap has two components.  The actual
transition functions there are
\(\varphi_1(\theta)=g(\theta)\varphi_{1,0}\) and
\(\varphi_2(\theta)=g(\theta)\varphi_{2,0}\), with
the same scalar from (P3K.17.3).  Change the cap frame by
\(g(\theta)^{-1}\).  Both angular factors disappear
simultaneously, leaving parameter-independent transitions.
This explicitly excludes a line class on
\(S^1_{\rm joint}\times S^1_{\rm parameter}\).
Both six-dimensional characteristic numbers vanish:
\[
  \langle x^3,[Q_{\rm corr}\times S^1]\rangle=0,
  \qquad
  \langle p_1x,[Q_{\rm corr}\times S^1]\rangle=0.       \tag{P3K.17.4}
\]

\par\medskip\noindent\textbf{The additional normal-sphere boundary in the surgery bordism.}\enspace

The fixed locus in the master bordism must also be retained.
There are exactly two relevant auxiliary-quotient fixed charts:
\((z_4,z_1)\) and \((z_0,z_2)\).  Each has auxiliary
weights \((1,-1)\), and its quotient is \(\mathbb R^3\)
by the explicit invariant map
\[
      (z_+,z_-)\longmapsto
       \bigl(2\Re(z_+z_-),2\Im(z_+z_-),
                         |z_+|^2-|z_-|^2\bigr).
\]
The inverse orbit is determined by
\(|z_\pm|^2=(\sqrt{\mu^2+4|z_+z_-|^2}\pm\mu)/2\).
The two moment faces in each chart are discs.  Removing the
inner master-space ball removes an interior three-ball in
each of these fixed charts.  Away from the flop curve the
fixed sphere minus its two endpoint discs is unchanged.
Gluing this annulus times the moment interval proves that
the fixed three-dimensional cobordism is exactly
\[
                 W_F=(S^2\times I)\setminus
                             (\mathring D^3\sqcup\mathring D^3).
\]
The Cox reindexing preserves the two charts; it does not
permute or identify their boundary spheres.

The oriented normal three-plane bundle is trivial over this
three-dimensional cobordism, with its boundary product
structures retained.  On the two Cox fixed charts its
positive-weight complex normal bundle is respectively
\(L^{\oplus2}\) and \((L^{-1})^{\oplus2}\), so the
associated three-plane bundle has \(w_2=0\).  The other
two boundary spheres have the original \(SU(2)\) normal
model.  Restriction of \(H^2(W_F;\mathbb Z/2)\) to
the four boundary spheres is injective, so \(w_2\) vanishes
globally.  An oriented three-plane bundle on a three-dimensional
complex with this class zero is trivial: the obstruction on
two-cells is \(\pi_1(SO(3))=\mathbb Z/2\), and that
on three-cells is \(\pi_2(SO(3))=0\).  Boundary
trivializations differ by maps \(S^2\to SO(3)\), all
null-homotopic.  Collar homotopies therefore match the chosen
product trivializations at every boundary.  The actual circle
line restricted to the constant normal section has boundary
values \(0,0,+h_A,-h_B\).

Cap the two original-family sphere boundaries by \(D^3\),
as in their actual conifold construction, and join the two
inner-face spheres by the charged cap \(C\) above.  This
closes \(W_F\) to \(S^2\times S^1\).  Crossing with
the parameter circle closes the fixed four-dimensional base
to \(\bar X\).  The described local trivialization makes
it the product with that parameter circle, retaining the
possible parameter-independent clutch along the first
\(S^1\).  In particular it is a mapping torus over the
first circle of \(S^2\), crossed with the parameter circle.
All gluings preserve the sphere orientation.  Its signature
is zero, and its tangent first Pontryagin number is zero.
For the exact orientation comparison, capping the first two
boundaries leaves \(S^2\times I\).  Its line has degrees
\(+1\) and \(-1\) in the two complex boundary markings.
Its product transport consequently has degree \(-1\)
between those markings.  The joint cap has exactly that
antipodal base change at its second end.  The composite has
degree \(+1\), giving the oriented sphere mapping torus
above.  On the meridian sphere the first end map is also
antipodal, so both meridian orientations give the same
physical Euler value one.

Apply the explicit collar construction of P3K.7 to the
complement of this fixed locus in the master-space orbit
family.  At the two original-family boundaries it gives
their actual conifold mapping tori.  At the inner face it
gives \(Q_{\rm corr}\times S^1\).  It has one further
boundary, the normal sphere bundle \(B=S(\nu_{\bar X})\).
The meridian-one line is glued by (P3K.17.2) and (P3K.17.3)
through every collar, so it extends across this oriented
bordism.  This accounts for the boundary which would be
missed by capping the two inner fixed spheres separately.

The actual parameter-dependent Cox action is delta, which
acts trivially on the circle orbit and its normal rank-three
bundle.  Thus on the closed fixed base both its normal data
and the line data are pulled back from the three-dimensional
\(S^2\)-mapping-torus factor.  With the notation of P3K.7,
\(e_{\mathrm{vert}}=2x-\pi^*b\), it follows that
\(b^2=0\).  The sphere-bundle calculation, valid as an
oriented characteristic calculation, gives
\[
       \frac{4\langle x^3,[B]\rangle
                  -\langle p_1(TB)x,[B]\rangle}{24}
       =\frac{3\langle b^2,[\bar X]\rangle
                    -\langle p_1(T\bar X),[\bar X]\rangle}{24}
       =0.                                             \tag{P3K.17.5}
\]
No value of \(p_1(\nu_{\bar X})\) is assumed here.

For a closed oriented six-manifold with a complex line,
\(4x^3-p_1x\) is an oriented characteristic number and
vanishes on oriented boundaries carrying that line.  Hence
the two original-family return values are equal after
(P3K.17.4) and (P3K.17.5).  The actual endpoint manifolds
already have their prescribed spin structures, so their
integer \(A\) is the same characteristic expression.
The signs of the two additional boundary components do not
affect their zero contributions; the two moment faces have
opposite boundary orientations.  Together with P3F.32 this identifies the negative-diagonal
return index with zero for all the stated orders and twists.
The exact Hopf transfer P3F.22 gives its oriented zero class,
and P3F.24 gives the product mapping torus.
\end{proof}

\subsection{Audit record}

The initial draft and its explicit audit boundary are preserved in
\path{code/part3_fourfold/history/part3_flop_cobordism_audit.md}.
Independent blind-control audit checked the residual semifree weights,
both associated-line signs, the common Cox delta gauge, and the extra
normal-sphere term.  Root audit additionally required and then accepted
the transported truncations and the free period-orbit quotient before
the two-overlap Cech gauge cancellation.  The orientation transport
between the two complex fixed-sphere markings was explicitly retained.
The exact matrix and gauge certificate is \path{code/part3_fourfold/cox_flop.json}.
This proof treats the two stated OS49 diagonal fans; it does not assign
the same quotient model to every periodic unimodular fan.

\begin{statement}{P3K.18}{Proved}
In the actual whole-fibre \(I_1\) orbit projection of P3K.16,
every sphere member has primitive source-fibre multiplicities
at every logarithmic cap.  In particular no cap whose smooth
source quotient has a genuinely multiple central fibre occurs
in such a sphere member.
\end{statement}

\begin{proof}
The signed meridian matrix in P3K.16 has determinant
\(\pm1\); its explicit genus-one Heegaard map identifies
the target \(B\) with \(S^3\).  Let \(K\subset B\)
be one oriented cap-core knot.  Suppose that the source
fibre multiplicities on that core have a common divisor
\(\nu>1\).  The linking-number class of this single knot is
\[
            \eta\in H^1(S^3\setminus K;\mathbb Z/\nu),
\]
with value one on its positive meridian.  It is represented,
for example, by intersection with an oriented Seifert surface
for \(K\), reduced modulo \(\nu\).  Pull it back to
\(Y\setminus p^{-1}K\).

The relevant local normal exponent is that of the actual
three-dimensional target, not the old branched base coordinate.
In a cap chart of P3K.16 it is
\[
        \xi=e^{2\pi i rt/m'}u,
        \qquad u=z^g,
\]
the transverse disc coordinate of its target solid torus.
At a smooth point of a source component of multiplicity
\(\nu a_j\), choose a complex normal coordinate \(w\).
Then \(u=w^{\nu a_j}v\), where \(v\) is a nonvanishing
holomorphic unit.  Consequently a small positive component
meridian has linking number \(\nu a_j\) under \(p\).
The pulled-back class kills every such meridian modulo
\(\nu\).

The remaining singular strata of the central source fibre,
crossed with the suspension direction, have real codimension
at least four in the smooth five-manifold \(Y\).  A loop
or a homotopy disc can be perturbed off these strata; after
this perturbation, van Kampen says that filling the smooth
divisorial pieces only kills their normal meridians.  Thus
the pulled-back class extends to a homomorphism
\[
                         \pi_1(Y)\longrightarrow\mathbb Z/\nu.
\]
It is nonzero.  Indeed a meridian loop of \(K\) can first
be perturbed off all the other critical and cap-core circles.
Lift it through the connected regular elliptic fibres by the
horizontal-path construction of P3K.16, and close the lifted
path inside its starting fibre.  The resulting loop lies in
\(Y\setminus p^{-1}K\) and evaluates to one.  This
contradicts \(\pi_1(Y)=0\).  Therefore the common divisor
is one at every cap.
\end{proof}

\begin{statement}{P3K.19}{Proved}
For the actual affine root-line caps of P3K.15 in the sphere
orbit scope of P3K.16, the root fibre is reduced.  A singular
root fibre has no nontrivial finite zero-gamma kernel.  At a
smooth elliptic root fibre that kernel is an actual normal
root change, and the invariant source coordinates and line
frames cancel it exactly.  Thus the cap has a primitive
effective pair without an unrecorded source or line quotient.

At an original multiplicative fibre \(I_N\), a sphere cap
retains the following necessary restrictions on the original
source data \(M=\mathcal O(aG-O+D)\), \(Q=bG\):
\[
              \gcd(N,b[G])=1,
              \qquad a[G]=0\quad\text{in }\mathbb Z/N.       \tag{P3K.19.1}
\]
For the explicit generator markings in P3F.2 these are
\(\gcd(N,b)=1\) and \(N\mid a\).  These restrictions
hold at logarithmically replaced multiplicative fibres as
well as at retained ones.
\end{statement}

\begin{proof}
\par\medskip\noindent\textbf{The actual component maps and their multiplicities.}\enspace
Use \(A,B,C,H\) from P3K.15, with
\(g=\gcd(m,n)\), \(C=A^{m/g}B^{-n/g}\).
The components of the finite source quotient are precisely
the \(C\)-orbits of root components.  The residual
\(H\)-action is transitive on those quotient components
by P3K.15.  Since \(C,H\) generate the same group as
\(A,B\), the latter group is transitive on all root
components.  Both maps preserve the multiplicities of
the root fibre \(\operatorname{div}(z)\).  The root
source retains the pulled-back section of the original
elliptic surface, so these multiplicities have greatest
common divisor one: intersecting the fibre with that
section gives one.  Transitivity makes every multiplicity
equal, and hence all are one.

Let \(H_j\subset C\) be the stabilizer of a general
point on root component \(C_j\).  It is the subgroup
acting pointwise on that component.  Since \(C\) is
central and \(A,B\) are transitive, all these subgroups
are the same subgroup \(H_0\), of order \(h\).
At a general point its action on the normal parameter has
order exactly \(h\), because the root-base character is
faithful.  The finite quotient is therefore ramified to
order \(h\) along each root component.  The quotient
base coordinate is \(u=z^g\).  If \(\bar C\) is an
image component, comparing divisors in
\(\pi^*u=z^g\) gives its multiplicity exactly
\[
                              g/h.                         \tag{P3K.19.2}
\]
This computation includes component permutations and the
degree of the component map; it uses the generic normal
ramification index, not a declaration that the quotient is
unramified.  P3K.18 forces \(g/h=1\).  Hence \(C\)
fixes every root component pointwise.

At a singular point of a reduced root fibre, a nontrivial
finite-order surface automorphism cannot fix that fibre
germ pointwise.  Average a local coordinate map over the
finite group to linearize the action, with derivative the
identity at the base point of the averaging map.  The fixed
germ of a nonidentity linear map in two complex dimensions
is either a point or a smooth line.  Neither can contain
the reduced singular fibre germ.  The derivative must
therefore be the identity; the linearization then makes
the automorphism the identity in a neighbourhood, and
analytic continuation makes it the identity on the root
source.  This contradicts its faithful root-base character.
Thus \(g=1\) at a singular reduced root fibre.

At a smooth elliptic root fibre, the group fixes that curve
pointwise.  Near each of its points average a tangent
coordinate to obtain an invariant coordinate \(e\), and
keep the given normal base coordinate \(z\).  These are
coordinates by the inverse function theorem, and the
generator acts exactly by
\[
                         (e,z)\longmapsto(e,\zeta_g z).
\]
The actual quotient coordinates are \((e,u=z^g)\).
They identify the root source with the pullback of its
smooth quotient along \(u=z^g\).  The invariant local
line frames of P3K.15 identify the root line with the
pullback quotient line as well; no fibre character remains
because every point stabilizer acts trivially on the line.
These local maps agree on overlaps by their quotient
construction and give the required global map after
shrinking the source disc.  On the quotient,
\[
  \bar A=\bar H^{n/g},\qquad
  \bar B=\bar H^{m/g},\qquad
  q=u^{m/g}.
\]
Thus the root order and gamma numerator reduce to the
primitive pair \((m/g,n/g)\) by an exact source-and-line
quotient.  The marked meridians and any remaining integral
delta collar factor are still those transported in P3K.15.

\par\medskip\noindent\textbf{The multiplicative root and its integral component action.}\enspace
On the minimal \(m\)-root of an \(I_N\) neighbourhood,
the Tate period is \(z^{Nm}\), and the components of
the \(I_{Nm}\) fibre are labelled by
\(\operatorname{ord}_z Z\pmod{Nm}\).  The canonical
root deck sends \(z\) to \(\zeta_m z\) and preserves
these labels.  A translation by a Tate point of valuation
\(\alpha\) adds \(\alpha\) to them.  If the original
point \(Q\) has component \(c\in\mathbb Z/N\), its
root lift has component \(mc\).  The actual affine
root map \(A\) therefore acts by addition of some
integer \(\alpha\), and \(B\) by addition of \(mc\).
The norm identity \(A^m=B^n\) gives
\[
       m\alpha\equiv nmc\pmod{Nm},\qquad
                   \alpha\equiv nc\pmod N.          \tag{P3K.19.3}
\]
This is the actual source-component action of the root
deck and its translation.  An additional inversion would
change the marked linear monodromy from trace two to
trace minus two and is not the outer monodromy of this
\(I_N\) source cap.

Every component orbit is a coset of the subgroup generated
by \(\alpha,mc\) in \(\mathbb Z/(Nm)\).  Its number
is \(\gcd(Nm,\alpha,mc)\).  By (P3K.19.3) every
common divisor of \(N,c\) divides this number.  The
proved transitivity thus forces \(\gcd(N,c)=1\).

The root line is the pullback of the quotient line of
P3K.15 and consequently has degree zero on every root
component.  Its generic Picard point is the pulled-back
point \(aG\), whose component is \(ma[G]\pmod{Nm}\).
Its extension differs from the section line by an integral
vertical divisor, by the coherent-direct-image argument
of P3K.15.  The component-degree equations are therefore
\[
                  e_{ma[G]}-e_0+B_{Nm}d=0.
\]
Multiply by the cyclic component labels and sum.  The
sum of labels times \(B_{Nm}d\) is divisible by \(Nm\),
by the exact cyclic difference calculation of P3K.14.
It follows that \(ma[G]=0\pmod{Nm}\), or equivalently
\(a[G]=0\pmod N\).  This proves (P3K.19.1).

The argument applies to the minimal root source used in
the actual Tate cap.  A superfluous equivariant blowup
cannot evade the transitivity statement: blowing up a
node introduces a component of different multiplicity,
and a rational vertical minus-one component in a reduced
transitive fibre would make every component a rational
minus-one curve.  The fibre intersection equations would
then give each vertex weighted valency one, so the
connected fibre would consist of two rational curves
with one transverse intersection.  Its arithmetic genus
would be zero, contradicting the elliptic fibre genus
one.  Thus the transitive root is relatively minimal.
\end{proof}

These conclusions concern the actual root-line caps and the
global projection explicitly constructed in P3K.15–16.  They
do not assert a replacement of an arbitrary periodic cusp fan
by that fixed-sphere orbit model.  At such a fan the requisite
orbit comparison remains a separate part of the exhaustive
scope question.

\subsection{Whole-fibre zero sites and source reduction}

\begin{statement}{P3K.20}{Proved}
For a semistable multiplicative source fibre \(I_N\), a
whole-fibre zero of order \(e\) has period determinant
\(Ne\) and central Mumford fibre Euler number \(2Ne\),
for every periodic unimodular triangulation.  At an
\(I_s^*\) zero, \(s>0\), the smooth quadratic
semistable Mumford root, followed by its actual monomial
involution quotient and toric quotient resolutions, has
Euler contribution at least \(4se\).  Potentially good
zero sites are excluded from sphere members by P3K.12.

In this actual smooth-root toric construction, a sphere
member therefore has exactly one whole-fibre zero: a
simple \(I_1\) zero.  For a rational rank-one source,
writing its generic line point and translation as
\(aG\) and \(bG\), respectively, the height denominator
\(d=1/h(G)\) satisfies \(ab=d\) after choosing the
contracting generator and the sign of \(G\).
\end{statement}

\begin{proof}
\par\medskip\noindent\textbf{The actual line-frame determinant.}\enspace On the infinite
Tate chain let the source period and suspension period have
valuation vectors
\[
                    h=(N,\alpha),\qquad g=(-c,\beta).
\]
Write the fibre-coordinate graph of the actual line extension
as heights \(h_j\), so \(h_{j+N}=h_j+\alpha\).
An arbitrary vertical line divisor changes these heights
integrally; it is not omitted.  The line lift has the same
zero order \(e\) on every component precisely when its
actual fibre-coordinate transition has defect
\[
                      h_j+\beta-h_{j-c}=e.             \tag{P3K.20.1}
\]
Put \(r=\gcd(N,c)\), and iterate this identity \(N/r\)
times.  The period relation on the left gives
\[
      -\frac c r\alpha=\frac N r(\beta-e),
      \qquad
      \det(h,g)=N\beta+c\alpha=Ne.                    \tag{P3K.20.2}
\]
This is an identity of the actual component line frames.
It proves cancellation of the off-diagonal valuation terms
without changing or normalizing the vertical line divisor.

A fundamental polygon of this period lattice has area
\(Ne\).  Every unimodular triangle has area \(1/2\),
so it contains exactly \(2Ne\) triangles, counted with
their periodic boundary identifications.  In the central
Mumford fibre each triangle gives one zero-dimensional
torus orbit.  All other orbit strata are positive-dimensional
complex tori and have Euler number zero.  Additivity of
Euler number over the finite orbit stratification gives
\(2Ne\).  No vertical edge or semifree property of the
chosen delta circle is assumed in this calculation.

\par\medskip\noindent\textbf{The starred multiplicative quotient.}\enspace The quadratic
semistable root of an \(I_s^*\) fibre is \(I_{2s}\),
and the pulled-back whole-fibre zero has order \(2e\).
The determinant just proved is \(4se\).  The smooth
root Mumford fibre therefore has \(8se\) zero-dimensional
torus strata and Euler number \(8se\).

The quadratic deck is monomial on the central strata.
Indeed on the source nodal identity torus it is
\(Z\mapsto uZ^{-1}\), with a nonzero central unit
\(u\).  It commutes with the scalar circle on the line,
so in a line frame it sends \(V\mapsto f(Z)V\).
The source, line and their linearization are algebraic,
and this central torus map is regular and invertible.
Thus \(f\) is a unit of \(\mathbb C[Z,Z^{-1}]\),
so \(f(Z)=vZ^{\ell}\).  The induced maps on preserved
boundary torus strata are the restrictions of this same
integral monomial map.  Nonpreserved strata are permuted.
An invariant refinement can be used for descent; every
subdivision increases the number of maximal cones and
does not lower the preceding root Euler count.

There is also an explicit invariant unimodular height-one
choice, without a restriction to the rank-one generator
marking.  Write the actual monomial deck on valuations as
\(R(x,y)=(-x,y-jx)\), with \(j=-\ell\in\mathbb Z\).
It sends the source period \(h=(2s,\alpha)\) to
\(-h\); comparison of the second coordinate gives
\[
                         j=\alpha/s.
\]
Thus the integrality of this shear follows from the
actual deck morphism.  The whole-zero identity is now
\(2\beta+cj=4e\).  For odd \(j\), the integral
shear \(y'=y-(j-1)x/2\) gives
\(R(x,y')=(-x,y'-x)\).  The positive unit-square
triangulation is invariant, since its edge vectors
\(u,v,u+v\) go to \(-(u+v),v,-u\).
For even \(j\), use \(y'=y-jx/2\).  The reflection
is then \((-x,y')\) and the periods are
\((2s,0),(-c,2e)\).  Give the square with lower-left
corner \((v,k)\) its diagonal according to the parity
\[
                       v+c\left\lfloor k/(2e)\right\rfloor.
\]
Both periods preserve this parity.  Reflection sends the
square to \((-v-1,k)\), changes the parity, and exchanges
the two diagonals.  This is the required invariant
unimodular periodic triangulation.

In the marked rank-one case the generic point \(aG\)
has root Tate lift \(Z_G^a\), hence \(\alpha=as\),
\(c=bs\), and \(j=a\).  Changing that lift by a
Tate-period power \(Q_T^k\) is the actual line-frame
map \(V\mapsto Z^kV\).  Vertical component divisors
change the graph heights periodically and do not change
the quasiperiod.  These explicit maps recover the
rank-one spelling while the preceding deck argument
proves the starred-zero bound before that specialization.

The canonical-form normalization also identifies the
quotient-resolution comparison.  In Weierstrass coordinates
\(q=z^2,x=z^2X,y=z^3Y\), one has exactly
\[
                  (dx/y)\wedge dq=2(dX/Y)\wedge dz.
\]
After wedging the logarithmic line form, this is the
nowhere-zero Kulikov form
\(\Omega=d\log Z\wedge d\log V\wedge dz\).
The source inversion and \(z\mapsto-z\) give two
minus signs, while the additional \(Z^{\ell}\) in the line
map drops out of the wedge.  Thus the deck preserves
\(\Omega\).  Its nontrivial stabilizers on a smooth
root have eigenvalues \((-1,-1,+1)\), and their quotient
is a transverse \(A_1\) surface singularity along a fixed
curve.  In a resolution chart \(v=ut,w=ut^2\) for
\(v^2=uw\), the invariant two-form is
\(\tfrac12du\wedge dt\); it is nonvanishing.  This
is the actual crepant quotient resolution.

A divisor in this resolved model cannot be contracted to
a smooth model carrying the same nonvanishing form: the
Jacobian of a birational morphism vanishes along every
divisor mapped into codimension at least two, contradicting
the equality of those forms on their common open set.
Changes of the height-one Mumford triangulation are already
covered by the period-area count for each such fan.  In
particular a pure base-root contraction, whose pullback
canonical form has a ramification zero, is not being
silently used as a comparison of these normalized forms.

On an invariant torus stratum a fixed point equation is
\(x^{M-I}=a^{-1}\).  Its solution is empty or a coset
of the kernel of an integral torus homomorphism.  That
kernel is a finite union of complex tori.  Its Euler
number is nonnegative: positive-dimensional components
give zero and points give one.  Euler additivity therefore
gives \(\chi(X'^{\iota})\ge0\) on the root central fibre.
The finite-action orbit count, applied stratumwise, is
\[
 \chi(X'/\langle\iota\rangle)
       =\frac{\chi(X')+\chi(X'^{\iota})}{2}
       \ge4se.                                        \tag{P3K.20.3}
\]

For completeness, the quotient-resolution correction has
a definite sign in the stated toric construction.  Along
each fixed torus stratum the normal cyclic quotient has
constant diagonal characters.  Resolve it by a toric
subdivision.  A compact exceptional fibre has finitely
many torus orbits; its positive-dimensional orbits have
Euler number zero and it has at least one zero-dimensional
orbit.  Its Euler number is consequently at least one.
Euler additivity over the base fixed torus strata shows
that replacing a point-normal fibre by this exceptional
fibre adds a nonnegative quantity.  This proves the same
lower bound after the actual toric quotient resolutions.
Flips between unimodular height-one triangulations keep
the triangle count unchanged.

\par\medskip\noindent\textbf{The actual invariant-twist map at a semistable zero.}\enspace
The logarithmic modification in the supplied construction
\texttt{con-log} uses an extending root automorphism and
an invariant lattice translation, with punctured comparison
\(\sigma=(\log z)/(2\pi i)\,v\).  A whole-zero line
lift itself is not invertible on the zero fibre.  If the
same invariant-translation formula is interpreted on its
semistable Mumford torus, its quotient can nevertheless
be described exactly.  Full-rank unipotent monodromy has
invariant lattice equal to the primitive vanishing-torus
lattice.  Thus its twist has no gamma component.
On \(q=z^m\) write its deck as
\[
 (Z,V,z)\longmapsto
       (\zeta_m^r Z,\zeta_m^t V,\zeta_m z).
\]
The invariant coordinates are
\[
               Z'=Z/z^r,\qquad V'=V/z^t,\qquad q=z^m.
\]
They preserve every original period \(c(q)q^\Gamma\).
On cocharacters the quotient lattice and its exact
isomorphism are
\[
 \mathbb Z^3+\mathbb Z(r,t,1)/m
      \longrightarrow\mathbb Z^3,
 \qquad (x,y,s)\longmapsto(x-rs,y-ts,ms).
\]
The two horizontal basis vectors and \((r,t,1)/m\)
map to the three standard basis vectors.  Conjugating
the root period shear with translation \(m\Gamma\)
by this map gives the original shear \(\Gamma\).
Hence a smooth Kulikov target is again a Mumford fan
for the original period lattice and has the same Euler
count.  This gives the quotient directly, including
cases where resolving its root first would add curves.
No noninvariant affine operation at a zero germ has
been assumed beyond these actual published maps.

\par\medskip\noindent\textbf{The global zero count and height equation.}\enspace Away from
zero fibres, the line-complement construction has no fixed
point for the scalar circle, including its finite affine
quotients.  Its compact fibre Euler number is zero.
This also follows directly from the corresponding real
mapping-torus description.  On the smooth part of the
base the fibre is a complex torus and also has Euler
number zero.  Stratifying the base by its finitely many
special values therefore expresses the total Euler number
as the sum of the zero-site contributions above.

A smooth or potentially good zero site has the elliptic
fixed component excluded by P3K.12.  Every semistable
multiplicative zero contributes the positive even integer
\(2Ne\).  Every starred multiplicative zero in the
stated quotient-resolution construction contributes at
least four.  Since a six-sphere has Euler number two,
there is exactly one contribution, with \(N=e=1\).

The literal whole-fibre Hom relation and the telescoping
height calculation of P3K.14 give
\(\deg E=ab\,h(G)\).  Hence \(ab=d\).
For the pole version of the construction replace the
actual suspension generator \(B\) by \(B^{-1}\).
Their cyclic groups are identical, and the corresponding
period lattices are identical with the gamma generator
reversed.  Thus their toric quotient is identified by the
identity on its common charts; the inverse lift has zeros
in place of poles.  Finally replacing the name of the
generator \(G\) by \(-G\) makes the two coefficients
positive without changing the represented line point or
translation.  This gives the stated positive factorization.
\end{proof}

The starred-fibre bound retains the normalized canonical form
and its actual monomial quotient/resolution maps.  The
invariant-twist paragraph supplies the direct quotient map
for the stated logarithmic formula; it does not assume an
additional noninvariant zero-germ operation.  The later
Seifert restrictions still require the orbit map of K16;
the present Euler count itself applies to every periodic
unimodular Mumford triangulation and does not impose a
vertical-edge or semifree delta condition.

\subsection{The starred cap and the finite generic source reduction}

\begin{statement}{P3K.21}{Proved}
In the established sphere orbit and root-line cap scope, a
nonzero-linearization \(I_s^*\) source fibre can occur only
with \(s=1\), a primitive quadratic cap, and
\[
                         a\equiv0\pmod2,
                         \qquad b\equiv2\pmod4.       \tag{P3K.21.1}
\]
The two \(I_3^*\) rank-one source configurations are excluded.
Combined with P3K.19–20 and the actual source table, the
remaining generic source coefficients are the 26 pairs in
the 16 profiles displayed below.  This is a proved finite
restriction on the generic source data; it does not declare
all remaining affine cap tuples to be sphere tuples without
their integral gluing calculation.
\end{statement}

\begin{proof}
\par\medskip\noindent\textbf{The marked quadratic root.}\enspace The quadratic root of
\(I_s^*\) is the Tate \(I_{2s}\) fibre.  Its deck
acts on component labels by \(j\mapsto-j\), since
it is elliptic inversion and fixes the zero component.
The two invariant components are \(j=0,s\).  On each,
inversion has two fixed smooth points.  Near such a point
the two local coordinates have signs \((-1,-1)\), so
the quotient has equation \(v^2=uq\).  Resolving it
adds a multiplicity-one outer component.  The quotient
of the cycle is a path of \(s+1\) multiplicity-two
components, with two such outer leaves at each end.
This is the actual affine \(D_{s+4}\) graph.

The end \(j=0\) contains the zero section.  Remove
its corresponding outer leaf to obtain the finite
\(D_{s+4}\) intersection matrix.  Its inverse has
diagonal entry one at the other leaf of that end, and
\((s+4)/4\) at the two far-end leaves.  These entries
follow directly from the graph matrix; they are also the
local-height entries used in P3F.2.  The primitive
rank-one generator meets a spinor leaf, so its quadratic
root component is exactly \(s\), not zero.
The further degree-\(m/2\) root multiplies component
labels by \(m/2\).  Thus on an even \(m\)-root
the source fibre is \(I_{sm}\), the generator has
component \(sm/2\), and the canonical deck still
acts by reflection \(j\mapsto-j\).

An odd root leaves a starred fibre with unequal
multiplicities, excluded by the reduced-root conclusion
of P3K.19.  For an even root write the source component
action of the affine cap map as
\(A(j)=-j+\alpha\).  The suspension map acts as
\(B(j)=j+bsm/2\).  Its root fibre is singular, so
P3K.19 makes the pair \((m,n)\) primitive.  In
particular \(n\) is odd.  Since \(A^m\) is the
identity on components, the actual norm relation
\(A^m=B^n\) forces \(b\) even.  Then \(B\)
has trivial component action.  A reflection has orbits
of size at most two; transitivity forces \(sm\le2\).
Hence \(s=1,m=2\), and \(\alpha\) is odd.
The multidegree-zero root line has generic component
\(asm/2\); its component-Laplacian equation forces
\(a\) even, as in P3K.19.

There is one further unit constraint, which the component
labels alone would miss.  Write the quadratic Tate period
as \(Q_T=z^2u(z^2)\) and the generator as
\(Z_G=zv(z)\).  Its actual descent under inversion says
\(Z_G(-z)=Q_T/Z_G(z)\).  Thus
\(v(z)v(-z)=-u(z^2)\), and the identity-component
limit of \(2G\) is exactly \(-1\).  For the
affine reflection write
\(A(Z,z)=(R(z)/Z,-z)\), where
\(R=z^\alpha w(z)\).  Its square multiplies
\(Z\) by \(R(-z)/R(z)\), whose central value is
\((-1)^\alpha=-1\).  The central value of the
right-hand side \(B^n\) is
\((-1)^{bn/2}\).  Since \(n\) is odd, equality
forces \(b\equiv2\pmod4\).  This proves (P3K.21.1)
with both the component and unit parts of the norm retained.

\par\medskip\noindent\textbf{The finite generic factor list.}\enspace A simple whole
zero requires an \(I_1\) source value and \(ab=d\)
by P3K.20.  The three profiles with no \(I_1\) are
therefore removed.  P3K.19 imposes \(N\mid a\)
and \(\gcd(N,b)=1\) at every multiplicative fibre.
The preceding argument removes both \(I_3^*\) profiles
and imposes (P3K.21.1) at \(I_1^*\).  Enumerating the
positive divisors of each of the eight height denominators
now gives exactly the following list.


\begin{center}\small\begin{tabular}{c l l}\hline
OS row & Source fibre profile & Remaining \((a,b)\) \\
\hline
43 & \(III^*+3I_1\) & \((1,2),(2,1)\) \\
43 & \(III^*+II+I_1\) & \((1,2),(2,1)\) \\
45 & \(I_8+4I_1\) & \((8,1)\) \\
45 & \(I_8+II+2I_1\) & \((8,1)\) \\
47 & \(I_7+I_2+3I_1\) & \((14,1)\) \\
47 & \(I_7+I_2+II+I_1\) & \((14,1)\) \\
47 & \(I_7+III+2I_1\) & \((7,2),(14,1)\) \\
49 & \(IV^*+I_2+2I_1\) & \((2,3),(6,1)\) \\
49 & \(IV^*+III+I_1\) & \((1,6),(2,3),(3,2),(6,1)\) \\
50 & \(I_1^*+I_3+2I_1\) & \((6,2)\) \\
50 & \(I_1^*+IV+I_1\) & \((2,6),(6,2)\) \\
55 & \(I_5+I_4+3I_1\) & \((20,1)\) \\
55 & \(I_5+I_4+II+I_1\) & \((20,1)\) \\
56 & \(I_5+I_3+I_2+2I_1\) & \((30,1)\) \\
56 & \(I_5+I_3+III+I_1\) & \((15,2),(30,1)\) \\
56 & \(I_5+IV+I_2+I_1\) & \((10,3),(30,1)\) \\
\hline\end{tabular}\end{center}

The script \path{code/part3_fourfold/source_reduction.py} starts
from all 21 source profiles of P3F.1, retains their marked
height denominators, and executes these exact tests.  It
also exports the affine \(D\) matrices and their inverse
entries, so the starred marking is independently checkable.

\par\medskip\noindent\textbf{The indispensable vertical correction in row 50.}\enspace
The list allows arbitrary vertical line divisors, as the
literal construction does.  For its two row-50 factors,
the necessary correction is explicit.  Label the affine
\(D_5\) fibre by central vertices \(C_0,C_1\),
identity outer leaf \(L_0\), its order-two mate \(L_1\),
and spinor leaves \(R_0,R_1\).  The two left leaves
meet \(C_0\), the two right leaves meet \(C_1\),
and \(C_0,C_1\) meet.  Translation by \(2G\) or
\(6G\) swaps \(L_0,L_1\) and \(R_0,R_1\),
while fixing the two central components.

For \(a=2\) or \(6\), the section \(aG\) meets
\(L_1\), so the bare section-line degree vector is
\(e_{L_1}-e_{L_0}\).  If \(B_D\) is the actual
affine intersection matrix, then
\[
 e_{L_1}-e_{L_0}-B_D e_{L_0}
                  =e_{L_0}+e_{L_1}-e_{C_0}.          \tag{P3K.21.2}
\]
The right-hand side is invariant under the stated
translation.  Thus the actual line
\(\mathcal O(aG-O-L_0)\), with \(L_0\) here a
vertical fibre component, has the required invariant
component degrees.  At the other reducible fibre the
displayed factor pairs likewise have invariant component
degrees: at \(I_3\) the line point is in the identity
component; at \(IV\), either the line point is in the
identity component or the translation preserves all
components.  The Hom ratio is generically trivial and
has degree zero on every fibre component.  Its vertical
divisor is therefore a sum of whole fibres, and P3K.14
gives its base degree one.  Choosing its section places
the simple zero at the selected \(I_1\).

This proves that these corrected source line data satisfy
the literal whole-fibre Hom condition.  It does not
replace the remaining integral cap/gluing verification
for their total threefolds.
\end{proof}

The 26 pairs are a finite necessary source reduction in
the declared orbit and cap scope.  The full cap data still
include their exact cyclic norm, logarithmic site, integral
delta collar, Seifert determinant, and integral topology.
No unverified pair in this table is assigned a sphere or
a return index solely by this arithmetic list.

\subsection{Primitive-translation sources and the exact orbit filling}

\begin{statement}{P3K.22}{Proved}
For the following ten source profiles, let \(G\) be the marked
primitive Mordell--Weil generator, put \(d=1/h(G)\), and choose
\(M=\mathcal O(dG-O)\), translation \(Q=G\), a simple whole-fibre
zero at an \(I_1\), and one order-one logarithmic gluing with gamma
coefficient \(+1\) at a smooth source fibre.  Use the standard
vertical-edge unimodular cusp fan.  These are compact complex
threefolds diffeomorphic to \(S^6\).  Their positive scalar-parameter
loops have the indicated Jin integer and oriented \(\Theta_7\) class.


\begin{center}\small\begin{tabular}{c p{0.48\textwidth} c c c}\hline
OS row & Source profile & d & Jin integer A & Class mod28 \\
\hline
45 & \(I_8+4I_1\), \(I_8+II+2I_1\) & 8 & \(-3\) & 25 \\
47 & \(I_7+I_2+3I_1\), \(I_7+I_2+II+I_1\), \(I_7+III+2I_1\) & 14 & \(-6\) & 22 \\
55 & \(I_5+I_4+3I_1\), \(I_5+I_4+II+I_1\) & 20 & \(-9\) & 19 \\
56 & \(I_5+I_3+I_2+2I_1\), \(I_5+I_3+III+I_1\), \(I_5+IV+I_2+I_1\) & 30 & \(-14\) & 14 \\
\hline\end{tabular}\end{center}

The compact page is the indicated rational elliptic surface itself.
Its line is \(M\), with square \(-d\), and
\((dG)\cdot O=(d-2)/2\).  Every displayed mapping torus is therefore
not diffeomorphic to \(S^1\times S^6\), by the same oriented
mapping-torus detection used in P3K.10 and P3F.25.

There is also a restriction on additional smooth-site caps for these
primitive sources.  After the actual finite-kernel cancellation of
P3K.19, let their primitive cap data be
\((m_i,n_i,\lambda_{ai},\lambda_{bi})\), with
\(\gcd(m_i,n_i)=1\).  The exact orbit Mayer--Vietoris computation
below gives, whenever its first-homology gluing matrix is unimodular,
\[
        H_2(Y;\mathbb Z)=\bigoplus_i(\mathbb Z/m_i)^2.       \tag{P3K.22.1}
\]
Consequently a sphere member with only such smooth-site modifications
has \(m_i=1\) at every site.  For order-one sites the total gamma
coefficient \(\sum n_i=\pm1\) gives sphere topology, including
arbitrary integral source and delta collar coefficients.  The
individual index table above specifies the one-site positive-gamma
family; it does not identify uncomputed singular-source cap choices
with this family.
\end{statement}

\begin{proof}
\par\medskip\noindent\textbf{The source automorphisms and the line lift.}\enspace
P3F.1 constructs the source surfaces, and P3F.2 fixes their component
markings and proves \(G\cdot O=0\).  Translation and inversion
extend over their resolved singular fibres.  One direct proof keeps
the actual rational maps.  On the generic Weierstrass fibre use its
chord-and-tangent addition and \((x,y)\mapsto(x,-y)\).  Both preserve
the invariant differential up to the constant sign of inversion.
On a rational minimal elliptic surface the form
\((dx/y)\wedge dt\) has canonical divisor \(-F_\infty\), after a
base frame is chosen; the local minimal resolutions preserve this
relative differential.  Resolve the graph of either rational map,
with projections \(p,q:Z\to S\).  They are maps over the same base.
Pulling back that differential gives
\[
 -p^*F_\infty+E_p=-q^*F_\infty+E_q,
 \qquad p^*F_\infty=q^*F_\infty.
\]
Here the Jacobian discrepancy divisors have strictly positive
coefficient on every exceptional curve of their respective
compositions of point blowups.  Thus \(E_p=E_q\), so both
projections contract exactly the same curves and connected fibres.
The graph maps descend in both directions, giving mutually inverse
holomorphic automorphisms of \(S\).  This proves the required
extension, including II, III and IV fibres.  On an \(I_N\) chain
the actual Tate map is multiplication by the section unit and a
shift by its component label.  On III and IV, translation acts on
the dense smooth parts by addition in the component groups of
orders two and three; it therefore permutes the resolved components
by those same labels.  Inversion conjugates translation by \(G\)
to translation by \(-G\) on the whole surface.

The point \(dG\) is in the identity component at every reducible
fibre in the table.  Thus \(M\) has multidegree zero there.
The height identity gives
\((dG)\cdot O=(d-2)/2\), and its actual divisor class satisfies
\[
 (t_G)_*[M]-[M]=[F],\qquad t_G^*M\cong M(-F).           \tag{P3K.22.2}
\]
These equalities are also checked directly in the integer matrices
below.  The numerical equality is an equality of line bundles
because the rational surface has \(\operatorname{Pic}^0(S)=0\).
The resulting Hom line is \(\pi^*\mathcal O_{\mathbb P^1}(1)\).
Choose its section with its unique zero at the selected nodal fibre.
This gives the actual holomorphic lift with a simple whole-fibre
zero.

For an explicit meromorphic frame of that lift, put \(P_j=jG\).
On the generic cubic the quotient of the chord through
\(P_d,-G\) by the vertical line through \(P_{d-1}\) has divisor
\(P_d+(-G)-P_{d-1}-O\).  It is exactly the fibre divisor of
\(M\otimes(t_G^*M)^{-1}\).  Its remaining vertical orders form
whole fibres, since the component degrees vanish.  Multiplying by
the rational base function that moves this degree-one divisor to the
chosen \(I_1\) supplies the preceding section.  Thus its frame and
all local multipliers are obtained from the stated rational addition
map, rather than from an unspecified linearization.

For sufficiently small \(|\lambda|\), multiply this lift by
\(\lambda\).  In a fixed Hermitian metric it is uniformly strictly
contracting, away from its prescribed zero where the Mumford charts
are used.  Quotienting the nonzero line by its iterates is therefore
free and proper.  At the zero the actual valuation periods are
\((1,0),(0,1)\); the standard vertical-edge fan gives the smooth
index-one filling with its one fixed \(\mathbb P^1\), of normal
weights \((-1,+1)\).  The chart and line-gluing formulas are the
same ones used in P3K.10, with the present holomorphic unit logarithms
\(\Gamma,C\).  In particular, on a universal scalar-parameter
chart its cusp period and two origins are
\[
 \begin{aligned}
 \gamma&=(\Gamma,C+s+v),\\
 \sigma&=(\Gamma v,Cv+(v^2-v)/2),\\
 o&=(-\Gamma s,-Cs-s^2/2-sv),
 \end{aligned}                                                     \tag{P3K.22.3}
\]
and the closing ordinary line gauge is \(-s\delta\).  The exact
identities
\[
 \sigma(v+1)-\sigma(v)=\gamma-s\delta,
 \quad (\sigma-o)(v+\theta)-(\sigma-o)(v)
      =\theta\gamma+\tfrac12\theta(\theta-1)\delta
\]
hold for these actual units.  The degree-two cusp root used in
P3K.10 retains the half-integral ordinary-origin term and its integral
deck.  The smooth order-one cap is glued by
\((\log q)/(2\pi i)\,\gamma\); its scalar closing gauge is the
same actual base function used above.  Thus the gluing is holomorphic
in the scalar parameter and closes on its positive loop.  All filled
pieces are proper over their base discs; the contraction shell gives
properness on their common complement.  The compact base gives
compactness.  This establishes the analytic construction and the
positive-gamma parameter direction needed for the index calculation.

\par\medskip\noindent\textbf{Every integer matrix of the source mapping torus.}\enspace
Use the divisor basis
\[
 O,F,\{C_{i,j}:1\le j<N_i\},G,
\]
where \((N_i,k_i)\) is respectively
\((8,3)\), \((7,2),(2,1)\), \((5,2),(4,1)\), or
\((5,1),(3,1),(2,1)\).  A root \(A_1\) may denote III and
a root \(A_2\) may denote IV; their displayed component intersection
and permutation matrices are unchanged.  The Gram matrix has
\(O^2=G^2=-1\), \(O\cdot F=G\cdot F=1\),
\(G\cdot C_{i,j}=\delta_{j,k_i}\), root blocks
\(-A_{N_i-1}\), and all other entries zero.  Its determinant is
minus one in every case.  Since \(H_2(S;\mathbb Z)\) is free
unimodular of rank ten, these classes are an integral basis.

Let \(R_j\) be the vector whose \(i\)-th root block is column
\(jk_i\bmod N_i\) of \(A_{N_i-1}^{-1}\), or zero for component
zero.  Put
\[
 \varphi=G-O-F+R_1,\quad
 P_j=O+(P_j\cdot O+1)F+j\varphi-R_j,
 \quad
 P_j\cdot O=\frac{j^2/d-2+\sum_i\operatorname{contr}_i(jG)}2.
\]
The actual pushforward matrix \(T=(t_G)_*\) is now specified
column by column: \(O\mapsto G\), \(F\mapsto F\),
\(G\mapsto P_2\), and
\(C_{i,j}\mapsto C_{i,j+k_i}\), with
\(C_{i,0}=F-\sum_{j>0}C_{i,j}\).  All columns are integral.
The executed program verifies \(T^t\mathcal G T=\mathcal G\) and
\[
 \operatorname{Smith}(T-I)=(1,1,1,1,1,1,1,1,1,0),
 \quad\ker(T-I)=\mathbb ZF,
 \quad\operatorname{coker}(T-I)=\mathbb Z[O].             \tag{P3K.22.4}
\]
It exports every ten-by-ten matrix in
\path{code/part3_fourfold/primitive_source_translation.json}.

There is also a direct primitive-presentation check.  In the
coinvariants the components of each fibre become one class \(c_i\),
with \(N_i c_i=F\), and \(G=O\).  The remaining section relation
has coefficient \(P_2\cdot O-1\) on \(F\) and coefficient
\(\sum_j(2R_1-R_2)_{i,j}\) on \(c_i\).  Its Schur complement is
\(1/d\), while \(\prod_iN_i=d\).  The resulting determinant is
therefore \(\pm1\), not merely nonzero rationally.  For the four
rows the determinants are respectively \(-1,1,1,-1\).

The actual contracting lift covers \(t_{-G}\).  Its source matrix
has the same image and kernel, since
\(T^{-1}-I=-T^{-1}(T-I)\).  The mapping torus
\(W=T_{-G}(S)\) has the cellular cone complex with groups
\(\mathbb Z,\mathbb Z,\mathbb Z^{10},\mathbb Z^{10},
\mathbb Z,\mathbb Z\) in degrees zero through five; the sole
nonzero differential is \(T^{-1}-I:C_3\to C_2\).
Thus each \(H_i(W)\), \(0\le i\le5\), is \(\mathbb Z\),
and \(\pi_1(W)=\mathbb Z\gamma\).

The actual delta quotient before logarithmic gluing is this marked
mapping torus.  On its ordinary part the Hermitian logarithmic norm
is the suspension coordinate, as in P3K.13.  At the chosen \(I_1\),
the section \(G\) has a nonzero holomorphic Tate-unit representative
\(g(q)\).  Choose its logarithm on the disc.  The isotopy
\(Z\mapsto g(q)^u Z\) has node-chart formula
\((x,y)\mapsto(g(q)^u x,g(q)^{-u}y)\), so it extends through the
node for every real \(u\).  It identifies the source mapping torus
with the same \(I_1\)-surface-times-circle product of P3F.9, on
the actual boundary collar.  The bounded line units alter only the
logarithmic-norm coordinate and give the same product map.  This
supplies the comparison across the cusp, including its source
translation, rather than replacing that piece by a different fill.

\par\medskip\noindent\textbf{The primitive rim bases and the decisive involution.}\enspace
Remove \(r\) disjoint regular-fibre tubes
\(N_i=T^2\times D^2\times S^1_\gamma\) from \(W\), and call
the complement \(V\).  Write their boundary bases
\((a_i,b_i,\alpha_i,\gamma)\), with \(\alpha_i\) the positive
base meridian.  The Thom maps
\(H_j(W)\to\bigoplus_iH_{j-2}(T^3)\) are, in descending degrees,
\[
 \begin{array}{c|c}
 j&\text{map}\\\hline
 5&1\longmapsto(1,\ldots,1)\text{ in the }ab\gamma\text{ columns}\\
 4&1\longmapsto(ab)_1+\cdots+(ab)_r\\
 3&0\\
 2&1\longmapsto(1,\ldots,1).
 \end{array}                                                     \tag{P3K.22.5}
\]
For the zero map, the dual inclusion \(H_2(N_i)\to H_2(W)\)
is zero: the invariant class \(F\), which pairs one with \([O]\),
is pulled back from the source base and restricts to zero on the
whole tube core.  The other maps follow from the oriented fibre and
section intersections.  They are primitive diagonal maps.

The Thom sequence now gives free groups
\[
 H_1(V)=\mathbb Z^r,\quad H_2(V)=\mathbb Z^{3r},
 \quad H_3(V)=\mathbb Z^{3r},\quad H_4(V)=\mathbb Z^{r-1}.
\]
In degree two an integral basis is exactly
\(A_i=a_i\alpha_i,B_i=b_i\alpha_i,C_i=\alpha_i\gamma\).
In degree one use \(\gamma,\alpha_1,\ldots,\alpha_{r-1}\), with
\(\alpha_r=-\sum_{i<r}\alpha_i\).  In degree three use
\(z=ab\gamma\), the \(r-1\) independent classes \(ab\alpha_i\),
and all \(a_i\alpha_i\gamma,b_i\alpha_i\gamma\).  In degree four
use \(r-1\) boundary fundamental classes with their one sum relation.

The remaining boundary coefficients are fixed by an actual smooth
involution.  Inversion on the source conjugates the suspension map
to its inverse, so
\[
                         I(p,u)=(-p,-u)                         \tag{P3K.22.6}
\]
is a well-defined diffeomorphism of its mapping torus and preserves
every removed tube.  In a tube choose a real lift of the local
section and write its mapping-torus trivialization as
\(p'=p+uQ\), where \(Q\) is its signed suspension translation.
Then (P3K.22.6) is literally \((p',\alpha,u)\mapsto(-p',\alpha,-u)\).
Thus \((a,b,\alpha,\gamma)\mapsto(-a,-b,\alpha,-\gamma)\).
Every degree-two rim basis vector is odd, whereas
\(ab,a\gamma,b\gamma\) are even.  Naturality forces each image
of an even class to be two-torsion in \(H_2(V)\).  This group is
free by the primitive Thom calculation, so every such image is zero.
In particular the previously undetermined boundary linking integers
are exactly zero integrally.

The same argument gives the full degree-three boundary map.
The three rim classes \(ab\alpha,a\alpha\gamma,b\alpha\gamma\)
are even.  Each \(ab\gamma\) is odd and projects to the same
generator of \(H_3(W)\).  Subtracting two of them gives both an
odd class and a class in the even rim subgroup; freeness again
forces equality.  This proves the basis description and every
boundary coefficient used next.

\par\medskip\noindent\textbf{The actual cap maps and all mayer--vietoris matrices.}\enspace
On the orbit boundary a primitive smooth cap has meridian
\[
 \mu_i=m_i\alpha_i+n_i\gamma+\lambda_{ai}a_i+\lambda_{bi}b_i.
\]
The coefficient vector is the full cyclic norm from P3K.15; no
source-period term is discarded.  The smooth root source, followed
by its extending translation and base rotation, is a mapping torus
of a disc-times-torus diffeomorphism isotopic to the identity by
those explicit rotations and translations.  It is therefore
\(T^2\times D^2\times S^1\), with the displayed marked meridian.
Choose integers \(r_i,s_i\) with \(r_im_i+s_in_i=1\).
The four columns
\[
 a_i,\ b_i,\ \mu_i,\ \kappa_i=-s_i\alpha_i+r_i\gamma
\]
form an integral determinant-one basis.  Let \(Q_i\) be the
matrix of its inverse with the meridian row deleted.  Then
\(Q_i:\mathbb Z^4\to\mathbb Z^3\) is the actual cap inclusion
on first homology, and its higher maps are exactly
\(\bigwedge^jQ_i\).

For the exterior, in boundary wedge order
\((ab,a\alpha,a\gamma,b\alpha,b\gamma,\alpha\gamma)\), the
degree-two block is
\[
 i_2=\begin{pmatrix}
 0&1&0&0&0&0\\0&0&0&1&0&0\\0&0&0&0&0&1
 \end{pmatrix}.
\]
The degree-one map sends \(a_i,b_i\) to zero, \(\gamma\) to
the common gamma generator, and \(\alpha_i\) to the meridians with
their stated one sum relation.  The degree-three map in wedge order
\((ab\alpha,ab\gamma,a\alpha\gamma,b\alpha\gamma)\) sends
these columns to the corresponding rim class, the common \(z\),
and the two displayed independent rim classes.  The degree-four
map is \((I_{r-1}\mid-\mathbf1)\).  These descriptions specify
every entry of the integral matrices \(i_j\).

The Mayer--Vietoris map in degree \(j\) is the block matrix
\(D_j=\begin{pmatrix}i_j\\-\bigoplus_i\bigwedge^jQ_i\end{pmatrix}\); in degree zero it
is the column-incidence matrix \((\mathbf1,-I_r)^t\).
For one positive order-one cap with zero source collar coefficients,
the three nontrivial square matrices are explicitly
\[
 D_1=\begin{pmatrix}
 0&0&0&1\\-1&0&0&0\\0&-1&0&0\\0&0&1&-1
 \end{pmatrix},\qquad
 D_3=\begin{pmatrix}
 0&1&0&0\\0&0&1&0\\0&0&0&1\\1&-1&0&0
 \end{pmatrix},
\]
\[
 D_2=\begin{pmatrix}
 0&1&0&0&0&0\\0&0&0&1&0&0\\0&0&0&0&0&1\\
 -1&0&0&0&0&0\\0&1&-1&0&0&0\\0&0&0&1&-1&0
 \end{pmatrix}.
\]
All three are unimodular.  The degree-four domain is \(\mathbb Z\)
with zero target, giving the oriented top class of \(Y\).  Thus
\(H_*(Y;\mathbb Z)=H_*(S^5;\mathbb Z)\).

For general smooth-site data, eliminating the cap degree-two groups
leaves the exact block presentation
\[
 \begin{pmatrix}
 m_i&0&0&\lambda_{ai}\\
 0&m_i&0&\lambda_{bi}\\
 0&0&m_i&-n_i
 \end{pmatrix}                                                   \tag{P3K.22.7}
\]
on \((A_i,B_i,C_i)\).  Its cokernel is
\((\mathbb Z/m_i)^2\), since \(n_i\) is a unit modulo \(m_i\).
The first-homology determinant is the signed Seifert determinant
\(D=\sum_i n_i\prod_{j\ne i}m_j\), and the degree-three
determinant has the same absolute value.  These identities, every
matrix \(D_0,\ldots,D_4\), and 196 primitive two-cap cases with
nonzero source collars are executed and exported by
\path{code/part3_fourfold/primitive_orbit_filling.py} and its JSON.
This proves (P3K.22.1), including its torsion.

For order-one sites the fundamental-group calculation is equally
explicit.  The complement of the source fibres has fundamental
group generated by their base meridians: all elliptic cycles die
as in the rational source, and the section supplies the only sum
relation.  Thus
\(\pi_1(V)=\langle\alpha_i,\gamma\mid
\alpha_1\cdots\alpha_r=1,[\alpha_i,\gamma]=1\rangle\).
The cap imposes \(\alpha_i\gamma^{n_i}=1\); its source loops
already die.  Hence \(\pi_1(Y)=\mathbb Z/|\sum_i n_i|\), which
is zero for the stated sphere data.  The source meridian assertion
can also be checked directly: the missing meridian of a single
regular fibre bounds the punctured zero section, and van Kampen
recovers \(\pi_1(S)=0\).

\par\medskip\noindent\textbf{The actual hopf filling and its six-dimensional homology.}\enspace
In the one-site construction \(Y\) is simply connected with the
integral homology of \(S^5\), hence is the standard oriented
five-sphere by the same smooth classification used for P3F.9.
The fixed curve is the actual embedded sphere furnished by the
index-one cusp chart.  Its normal Hopf meridian has Euler value one.
For \(P=Y\setminus\operatorname{int}\nu F\), Alexander duality
gives \(H^*(P;\mathbb Z)=H^*(S^2;\mathbb Z)\); the circle
bundle on \(P\) therefore has \(H_*(E)=H_*(S^3)\) by the
integral Gysin sequence.  The meridian value one makes
\(\pi_2(P)\to\pi_1(S^1)\) surjective, so \(\pi_1(E)=0\).

The other piece of \(X\) is \(F\times D^4\), with boundary
\(F\times S^3\).  The positive-weight SU(2) normal trivialization
of the actual \(\mathcal O(-1)\oplus\mathcal O(-1)\) chart gives
this product and its Hopf map.  Its Mayer--Vietoris matrices are
\((1,-1)^t\) in degree zero, \((-1)\) in degree two,
\((1)\) in degree three, and the zero map from \(\mathbb Z\)
in degree five.  All omitted groups are zero.  Therefore
\(H_*(X;\mathbb Z)=H_*(S^6;\mathbb Z)\) and van Kampen gives
\(\pi_1(X)=0\).  The already used smooth six-sphere classification
and \(\Theta_6=0\) identify this actual smooth manifold with
\(S^6\).  The same calculation excludes the nonzero torsion in
(P3K.22.1) from a sphere member; the quotient of a semifree sphere
with this fixed sphere must have zero such torsion.

\par\medskip\noindent\textbf{The actual page and the signed return.}\enspace
The order-one gamma gluing makes the orbit base the ordinary Hopf
\(S^3\), by the signed meridian matrix of P3K.16.  Its regular
binding page has degree one over the source base.  The norm shell
of \(M^\times\) maps injectively into the orbit: a positive iterate
strictly decreases its norm, and a negative iterate increases it.
At the whole-zero cusp the line graph is horizontal and lies in
the chosen standard fan; its successive suspension translates have
strictly increasing graph defect one, so no new identification of
this shell occurs there.  It is disjoint from the delta fixed
sphere.  On the smooth logarithmic collar the displayed log-gamma
map is the exact inverse angular coordinate, so the shell closes
as the page \(S\setminus\operatorname{int}\nu F\), with its
original source line \(M\).  This is the same actual norm-shell
and binding-collar comparison proved for OS45 in P3K.10; all its
maps here use the present \(\Gamma,C\), and no finite cap or
component-supported zero is introduced.

Thus the capped page is \(S\), of \(\chi(\mathcal O)=1\) and
signature \(-8\), while its line square is
\(M^2=-2-2(dG\cdot O)=-d\).  The real compression, meridian-line
comparison and signed characteristic calculation of P3K.6--7 apply
to this actual page and give
\[
 A=\frac{12M^2-p_1(S)}{24}
   =1-\frac d2=-\frac{d-2}{2}.                        \tag{P3K.22.8}
\]
The identities (P3K.22.3) identify one positive scalar-parameter
loop with one positive gamma turn, retaining the integral line
gauge and the degree-two cusp-origin root.  Hence (P3K.22.8)
computes the stated actual loop, with no unproved division by a
parameter cover.  The orientation dictionary of P3K.7 and P3F.25
reduces this integer to the classes in the table.  This proves all
the asserted source, topology, page and return claims.
\end{proof}

The statement supplies ten explicit sphere families and an exact
obstruction to additional smooth finite caps in this primitive-source
scope.  Singular-source cap replacements and arbitrary cusp-fan return
changes retain their separate actual comparisons; they have not been
identified with these examples by the source-lattice calculation.

\begin{statement}{P3K.23}{Proved}
Consider one of the actual sphere families constructed above with a
semifree fixed curve \(F\simeq\mathbb P^1\), whose complex normal bundle
is \(\mathcal O(-1)\oplus\mathcal O(-1)\). Perform the toric flop of
this fixed edge in the periodic cusp model, retaining the actual
parameter clutch and all exterior gluings. Write the original positive
return's Jin index as
\[
             A=\frac{4k-\ell}{24},
\]
The positive return of the flopped family has oriented class
\[
                         -5A\pmod {28}                 \tag{K23.1}
\]
in the run's normalization of \(\Theta _7\). The flopped curve has two
isolated circle fixed points. Formula (K23.1) specifies its homotopy
sphere class; it does not assign a new Jin integer to that circle
action.

In particular, the four groups of primitive-source examples in
P3K.22 give the following flopped classes:


\begin{center}\small\begin{tabular}{c l l}\hline
\(d=1/h(G)\) & Original index \(A\) & Flopped class in \(\mathbb Z/28\) \\
\hline
\(8\) & \(-3\) & \(15\) \\
\(14\) & \(-6\) & \(2\) \\
\(20\) & \(-9\) & \(17\) \\
\(30\) & \(-14\) & \(14\) \\
\hline\end{tabular}\end{center}

The earlier degree-four OS47 example of index \(-4\) gives class
\(20\). Every original zero-index example remains class zero after
this fixed-edge flop. No formula for a flop at another toric edge is
asserted here.
\end{statement}

\begin{proof}
\par\medskip\noindent\textbf{The actual two local products and the parameter clutch.}\enspace
Write the conifold link as
\[
 K=(S^3_a\times S^3_c)/S^1,\qquad
 (a,c)\sim(e^{i\phi}a,e^{-i\phi}c).
\]
With \(\mathbb H\) oriented by \((1,i,j,k)\), put
\[
 u=a_0+j a_1,\quad v=c_0+j\overline{c_1},\quad
 q=uv,\quad h=u i u^{-1},\quad k=v^{-1} i v.
                                                        \tag{K23.2}
\]
The first resolution has the actual product coordinates \((h,q)\)
on its real normal sphere bundle; the second has \((k,q)\).
To invert the first map, choose a unit representative \(u\) of \(h\)
and put \(v=u^{-1}q\). Changing the representative changes only the
displayed auxiliary-circle orbit. For the second choose \(v\) and
put \(u=qv^{-1}\). These inverse formulas are smooth in the standard
quotient charts. The common boundary map is consequently
\[
             \tau(h,q)=(q^{-1}h q,q).                  \tag{K23.3}
\]
The fixed sphere is unknotted in the smooth six-sphere:
Haefliger's Theorem 1.2 and Corollary 6.6, with \(n=2,q=4\),
give the trivial embedding class \cite{Haefliger1966}.
Its exterior is therefore the actual product
\(E=D^3_y\times S^3_q\), with its given boundary marking extended
by the normal framing. The map (K23.3) extends over this exterior:
\[
 T(y,q)=(q^{-1}yq,q),\qquad
 T^{-1}(y,q)=(q yq^{-1},q).                           \tag{K23.4}
\]
These are smooth also at \(y=0\). Thus replacing the local resolution
conjugates the exterior return by \(T\).

This last sentence includes the actual parameter family. For a
local Cox angular action of weights
\((-r-s,0,r,s)\), write
\[
 a_\theta=e^{-i(r+s)\theta/2},\qquad
 b_\theta=e^{i(r-s)\theta/2}.
\]
The first product action and second product action are respectively
\[
 (h,q)\longmapsto
 (a_\theta h a_\theta^{-1},a_\theta q b_\theta),\qquad
 (k,q)\longmapsto
 (b_\theta^{-1}k b_\theta,a_\theta q b_\theta).
\]
Substitution in (K23.3) proves exact equivariance for every \(\theta\).
At \(2\pi\), both scalar endpoints have the same central sign;
the displayed product actions close even when \(r+s\) is odd.
Therefore the comparison of the two local family products is the
same map \(\tau\), independent of the parameter angle. The period
reindexing acts on the periodic collection of common conifold cells.
That collection has trivial cell stabilizer; choose one quotient
representative and extend every chart change equivariantly to its
translates. Extracting the integer winding from each bounded
nonzero toric unit leaves a logarithm on this representative. Contract
its positive modulus and its remaining logarithm in both resolution
charts at once, transporting the truncation collars by the resulting
common conifold isotopy. The integer winding is exactly the Cox
action just displayed. This retains the actual exterior
identifications, including their bounded units. No null-homotopy of
an odd Cox loop, or additional disc-supported correction, is inserted.

\par\medskip\noindent\textbf{Exact relative generators and conjugation.}\enspace
Choose a smooth degree-one map
\(\zeta:(D^3,S^2)\to(S^3,1)\), constant on a boundary collar,
which is equivariant for quaternionic conjugation and commutes
with \(y\). Explicitly one can use
\[
 \zeta(y)=\cos\phi(|y|)-\sin\phi(|y|)\frac{y}{|y|},
\]
where \(\phi\) decreases from \(\pi\) to zero, has the smooth odd
radial expansion at zero, and is flat near the boundary. Define
\[
 L(y,q)=(y,\zeta(y)q),\qquad R(y,q)=(y,q\zeta(y)).
                                                        \tag{K23.5}
\]
They commute and extend over \(D^3\times D^4\) by the same quaternion
multiplications. Their induced six-sphere classes are therefore zero.
They represent the left and right generators of
\(\pi_3SO(4)=\mathbb Z^2\). Extension by the identity to the double
injects the relative pseudo-isotopy group into the doubled group;
Krylov's Proposition 1 and Lemma 3 give, in this dimension, the
relative group \(\Theta _7\oplus\mathbb Z^2\)
\cite{Krylov2004}. Pseudo-isotopy is sufficient here, since extending
the resulting collar gives the same oriented sphere gluing.

On the double write
\[
 {\cal A}(s,t)=(tst^{-1},t),\quad
 {\cal B}(s,t)=(s,sts^{-1}),\quad
 {\cal C}=[{\cal A},{\cal B}].
\]
Krylov's Theorems 2--3 establish centrality and order \(28\) of
\({\cal C}\), and the relation
\({\cal A}Y=Y{\cal A}{\cal B}\), where \(Y(s,t)=(s,st)\)
\cite{Krylov2001}. In the marked relative products
\[
       T={\cal A}^{-1},\qquad L=Y,\qquad
       R=Y{\cal B}^{-1}.
\]
The first identity below also follows by literal substitution in
(K23.4)--(K23.5); the second then follows from the displayed central
relation:
\[
 TLT^{-1}=R,\qquad
 TRT^{-1}=L^{-1}R^2{\cal C}.
\]
Consequently
\[
 T(L^mR^n{\cal C}^s)T^{-1}
       =L^{-n}R^{m+2n}{\cal C}^{s+n}.                \tag{K23.6}
\]
All maps on the left are relative to the boundary. The injection
just cited retains this equality in the relative group. Transport
of a supporting oriented disc shows that \(T\) fixes the central
disc-supported subgroup.

\par\medskip\noindent\textbf{The actual integer coefficient from the hopf quotient.}\enspace
The orientations in (K23.2) are essential. At \(a_1=x+iy=0\),
the derivatives of \(h\) are \((-2k,-2j)\), so \(h\) reverses the
complex curve orientation. At \(u=1\),
\[
 q=\Re c_0+i\Im c_0+j\Re c_1+k\Im c_1
\]
preserves the complex normal orientation. Thus the complex
orientation of the exterior is \(D^3_y\times S^3_q\).
For the second resolution the derivatives of \(k\) are \((2k,-2j)\);
its base preserves and its normal coordinate reverses the respective
complex orientations. Its exterior has the same product orientation.

The physical delta action is \(q\mapsto q e^{-i\theta}\), with
Hopf quotient \(z=q i q^{-1}\). At \(q=1\), its circle vector is
\(-i\), and the quotient derivatives in the \(j,k\) directions are
\(-2k,2j\). Hence
\[
 [S^3_q]=-[S^1_\delta][S^2_z],\qquad
 [E]=[S^1_\delta][D^3_y][S^2_z].                    \tag{K23.7}
\]
The closed conifold quotient therefore has orientation
\(S^3_y\times S^2_z\), equivalently \(S^2_z\times S^3_y\).

Let \(t\) be the relative variation of the actual exterior return in
this product. The integral circle Gysin argument can be made
explicit. Put \(P=Y\setminus\operatorname{int}\nu F\), the orbit
complement before its conifold cap. On its physical mapping torus,
whose fourth cohomology vanishes, let \(\eta\)
be a circle connection with \(d\eta=x\), and take
\(d\beta=x^2\). The integral degree-three generator of the circle
total space is represented by
\(-\eta x+\pi^*\beta\); fibre integration and the meridian-one
condition fix its sign. Its boundary restriction in the complex
curve orientation is \(s+k_{\rm Jin}h_F\theta\):
Stokes gives \(k_{\rm phys}=-\int_{F\times S^1}\beta\), and the
physical mapping torus is the inverse of Jin's convention.
Since \(h=-h_F\), its restriction in (K23.2) is
\[
                         t=-k_{\rm Jin}.             \tag{K23.8}
\]
Equivalently, the boundary of a relative \(D^3\times[0,1]\) is
\(S^2\times[0,1]-D^3\times\{1\}+D^3\times\{0\}\), giving precisely
this variation coefficient. The integral Gysin sequence supplies
the integral class; choosing differential representatives does not
replace the integral assertion.

The generator \(L\) has variation \(+1\) and quotient
\((y,z)\mapsto(y,\operatorname{Ad}_{\zeta(y)}z)\). In Jin's
identification \((v,0)\sim(g(y)v,1)\), a connection is
\(A_\theta=-\theta\,dg\,g^{-1}\). With \(\omega=dg\,g^{-1}\),
\[
 F=-d\theta\,\omega+(\theta^2-\theta)\omega^2,\qquad
 \int_{S^3}\operatorname{tr}(\omega^3)
                      =-24\pi^2\deg g.
\]
The orientation is \(S^3\times S^1\), so direct integration gives
\(c_2=\deg g\) and \(p_1(\operatorname{ad}g)=-4\deg g\).
The vertical Euler class of this sphere bundle is \(2x\), and
the base \(S^3\times S^1\) has \(H^2=0\). Therefore
\[
                       (k_L,\ell_L)=(-1,-4).
\]
Composing the actual equivariant return with \(L^{-t}\) kills its
variation and gives
\[
                      k_0=0,\qquad \ell_0=\ell-4k.    \tag{K23.9}
\]
Additivity follows from the mapping-torus pair-of-pants bordism
carrying the actual Euler line.

Put \({\cal B}=LR^{-1}\). The relative class after (K23.9) is
\({\cal B}^b{\cal C}^s\). Its double is the actual principal circle
bundle over the closed conifold mapping torus. Its tangent bundle
is the pullback of the quotient tangent bundle plus a trivial
line. Since \(H^2=\mathbb Zx\), \(H^4=\mathbb Z[y\theta]\), and
\(x^2=0\), the integral Gysin sequence makes the pullback on \(H^4\)
injective.

Represent \({\cal C}^s\) on a disc disjoint from
\(S^3_y\times\{q=1\}\). The representative
\({\cal B}^b{\cal C}^s\) fixes that sphere pointwise and has
normal clutch \(\operatorname{Ad}_{\zeta(y)^b}\) on
\(S^3_y\times S^1\). Its Pontryagin number is \(-4b\), by the
preceding connection calculation. The pullback of \([y\theta]\)
evaluates \(1\) on this cycle. This retains the actual marking:
an isotopy of representatives is the identity on an initial fibre,
and the possible \(x\theta\) ambiguity in a quotient degree-three
lift pulls back to zero. Thus
\[
 \ell-4k=-4b,\qquad b=6A,\qquad n=-6A.              \tag{K23.10}
\]
Here \(n\) is the \(R\)-exponent after rewriting the original
relative class in \(L,R\). This calculation uses the actual Hopf
quotient before the flop; it makes no corresponding quotient claim
after the fixed curve has become two isolated points.

\par\medskip\noindent\textbf{The sign of the central sphere by an oriented plumbing.}\enspace
Let \(\Sigma_f\) denote \(D^4_s\times S^3_t\) glued to
\(S^3_s\times D^4_t\), right coordinate equal to \(f\) of the left
coordinate, with the orientations from
\(\partial(D^4_s\times D^4_t)\). Take three products
\(D_i=D^4_s\times D^4_t\), with orientations \(+,-,+\).
Glue their faces by
\[
 (s_1,t_1)\longmapsto(s_2,t_2)
       =(s_1,s_1t_1s_1^{-1}),\qquad
 (s_2,t_2)\longmapsto(s_3,t_3)
       =(t_2^{-1}s_2t_2,t_2).
\]
The first is on the \(S^3_s\times D^4_t\) faces, the second on the
\(D^4_s\times S^3_t\) faces. Their coordinate orientations agree,
while the adjacent eight-ball signs differ, as required for oriented
gluing. After rounding, the two remaining faces have boundary word
\({\cal A}^{-1}{\cal B}\). The map \({\cal A}\) extends the left
face and \({\cal B}^{-1}\) the right face, so the oriented boundary is
\[
 \Sigma_{{\cal B}^{-1}{\cal A}^{-1}{\cal B}{\cal A}}
                       =\Sigma_{{\cal C}^{-1}}.
\]

Orient the first zero-section sphere by \(ds\) in \(D_1\) and
\(-ds\) in \(D_2\), and the second by \(dt\) in \(D_2\) and
\(-dt\) in \(D_3\). At their unique intersection in \(D_2\) their
concatenated orientation is \(-ds\,dt\), the ambient orientation.
Each normal clutch fixes a real line. Thus
\[
 Q=\begin{pmatrix}0&1\\1&0\end{pmatrix},\qquad \sigma(W)=0.
\]
The normal clutches are \(\operatorname{Ad}_s\) and
\(\operatorname{Ad}_{t^{-1}}\). Their normal orientations on their
two hemispheres are respectively \(dt\) and \(-ds\).
Quaternionic conjugation reverses the four-plane orientation and
commutes with every adjoint action, so this second normal-frame
orientation introduces no extra Pontryagin sign. The two
Pontryagin evaluations are \((p,-p)\), with \(|p|=4\).
The eight-manifold has only two four-handles besides its zero-handle;
in particular it is spin. Its unimodular intersection form gives
\[
 p_1^2=(p,-p)Q^{-1}(p,-p)^{\mathsf T}=-32.
\]
The spin Eells--Kuiper formula consequently gives
\[
 \mu(\Sigma_{{\cal C}^{-1}})=-\frac1{28},\qquad
 \mu(\Sigma_{\cal C})=+\frac1{28}.                   \tag{K23.11}
\]

There are two explicit reversals in comparing this plumbing with
P3F.22. Its positive Hopf sphere \(u\) has orientation \(-q\),
by (K23.7), and its ordered factors are \((u,y)\), whereas the
plumbing factors are \((y,q)\). Thus its pieces identify as
\[
 D^4_u\times S^3_y=-(S^3_y\times D^4_q),\qquad
 S^3_u\times D^4_y=-(D^4_y\times S^3_q).
\]
The pieces exchange and the ambient orientation reverses. The
literal gluing map is inverted, giving
\(\lambda(F)=-\Sigma_{F^{-1}}\). For a disc-supported class,
putting its supporting disc in the common boundary and cancelling
the opposite product collars gives
\(\Sigma_{F^{-1}}=-\Sigma_F\). Therefore
\(\lambda({\cal C})=\Sigma_{\cal C}\), with the positive sign
in (K23.11) in the run's convention.

Finally \(L\) and \(R\) extend the seven-disc, so the central exponent
of the original relative class is its actual class \(A\bmod28\).
Equation (K23.6) adds \(n\) to that exponent, and (K23.10) gives
\(n=-6A\). This proves (K23.1) and the displayed examples.

All local orientation matrices and the Chern--Weil integral are
checked exactly in \path{code/part3_fourfold/fixed_edge_signs.py}. The
central word, hyperbolic pairing and rational Eells--Kuiper
arithmetic are checked in \path{code/part3_fourfold/commutator_orientation.py}.
These programs certify those algebraic steps; the actual family,
circle, collar and oriented gluing maps are the constructions in
this proof.
\end{proof}
